\documentclass[12pt,oneside,a4paper,article]{abntex2}
\usepackage[utf8]{inputenc}
\usepackage[T1]{fontenc}
\usepackage[brazil]{babel}
\usepackage{graphicx}
\usepackage{tabularx}
\usepackage{amsmath}
\usepackage{authblk}
\usepackage{url}
\usepackage[left=3cm,right=2cm,top=3cm,bottom=2cm]{geometry}
\usepackage{setspace}
\hypersetup{colorlinks=true,linkcolor=blue,urlcolor=blue,citecolor=blue}
\title{\textbf{SkillSwap: plataforma acessível de troca de conhecimento e aprendizagem entre pares}}
\providecommand{\theforeigntitle}{}
\providecommand{\htmladdnormallink}[2]{#2}
\author[1]{Andressa Freire Galvão\thanks{andressa.galvao@ucsal.edu.br}}
\author[1]{Deivide Sobral da Silva Bezerra\thanks{deivide.bezerra@ucsal.edu.br}}
\author[1]{Diego Gabriel Pinheiro Damasceno\thanks{diego.damasceno@ucsal.edu.br}}
\author[1]{Flávia Rodrigues Santos Silva\thanks{flaviarodrigues.silva@ucsal.edu.br}}
\author[1]{Jorlan Heider Da Fé Santos\thanks{jorlan.santos@ucsal.edu.br}}
\affil{Análise e Desenvolvimento de Sistemas \par Escola de Tecnologias \par Universidade Católica do Salvador (UCSAL) \par Av. Prof. Pinto de Aguiar, 2589 Pituaçu, CEP: 41740-090 \par Salvador/BA, Brasil}
\date{Agosto 2026}
\renewcommand*{\Authsep}{, }
\renewcommand*{\Authand}{, }
\renewcommand*{\Authands}{, }
\renewcommand*{\Affilfont}{\normalsize\normalfont}
\renewcommand*{\Authfont}{\bfseries}
\setlength{\affilsep}{1em}
\begin{document}
\begin{center}
\includegraphics[width=0.3\textwidth]{ucsal_logo.png}
\end{center}
{\let\newpage\relax\maketitle}
\begin{center}\small\textit{Declaração de autoria: utilizamos a IA somente para correção ortográfica. A revisão de clareza, organização editorial, autoria, as decisões, o código e a validação pertencem à equipe.}\end{center}
\begin{resumoumacoluna}
Este trabalho apresenta a SkillSwap, uma plataforma web de troca de habilidades que aproxima pessoas que ensinam e aprendem por meio de uma jornada acessível, explicável e auditável. O sistema integra descoberta, proposta bilateral, conversa, agenda, conclusão, avaliação, troca de conhecimento sem cobrança e conexão institucional com aprendizagem e empregabilidade. A implementação é exclusivamente de software e utiliza React/Vite, Spring Boot, PostgreSQL com Flyway, Firebase Authentication, notificações FCM opcionais, WebSockets protegidos e observabilidade. A pesquisa diferencia evidência bibliográfica, implementação e hipótese de impacto, propondo métricas de tempo, erros, conclusão, acessibilidade percebida, confiança e intenção de continuidade.
\vspace{\onelineskip}
\noindent\textbf{Palavras-chaves}: troca de habilidades, acessibilidade, aprendizagem entre pares, recomendação explicável, empregabilidade, segurança.
\end{resumoumacoluna}
\textual
\noindent\textbf{Base bibliográfica.} A análise dialoga com estudos do IEEE sobre plataformas de troca de conhecimento, aprendizagem adaptativa, acessibilidade e recomendação explicável [1--9].
\section{Leitura do projeto e escopo final}
Durante a consolidação do MVP, percebemos que o material recebido misturava produto, TCC, deploy e uma proposta de escopo físico que deixou de ser desejada.

No protótipo, esse entendimento aparece assim: O produto final é exclusivamente software: frontend React, backend Spring Boot, PostgreSQL/Flyway, Firebase, integrações opcionais e operação no Railway.

Na implementação, este requisito é tratado como parte do sistema e não como uma tela isolada. A interface orienta a próxima ação; a API confere identidade, papel, estado e limites; a persistência guarda somente o necessário; e as integrações externas ficam condicionadas às variáveis do ambiente. O escopo continua exclusivamente de software, sem depender de escopo físico.

O fluxo começa quando a pessoa encontra uma necessidade ou uma habilidade que deseja compartilhar. Em seguida, ela confirma a intenção, executa a ação principal e recebe uma resposta que pode ser consultada depois. Para nós, leitura do projeto e escopo final só está concluído quando o caminho é compreensível, existe recuperação para erro e o servidor registra uma evidência coerente.

A forma mais honesta de verificar este requisito é registrar: Matriz de escopo, arquivos do repositório e contratos de rotas. O critério de aceite deve ser escrito antes do piloto, com entrada, resultado esperado, responsável, data e ferramenta de verificação. Assim, a documentação separa conformidade técnica, observação de uso e hipótese de impacto.

A validação combina teste automatizado, inspeção manual e uma tarefa representativa. Vamos registrar tempo, erros, conclusão, compreensão e barreiras observadas por teclado, leitor de tela, zoom, contraste, texto, áudio ou conexão instável. O objetivo não é criar uma aparência de certeza, e sim manter um ciclo curto entre observar, corrigir e conferir novamente.

A equipe não trata este item como resolvido antes de enfrentar o risco: A equipe pode continuar recebendo pedidos antigos de escopo físico; este documento é a fonte final e o escopo físico fica fora do TCC. O próximo passo é transformar a pendência em issue, teste, decisão de reunião ou evidência de piloto. Este capítulo não promete superioridade nem receita: ele deixa um contrato verificável para que a equipe consiga descobrir se a hipótese se sustenta.

\section{Nota de autoria e uso de IA}
A discussão de nota de autoria e uso de ia começou por uma situação concreta: a documentação precisa declarar a participação humana e evitar atribuir à IA dados, resultados ou referências não conferidos.

Nossa escolha para este ponto foi direta: A IA foi utilizada somente para correção ortográfica. A revisão de clareza, organização editorial, autoria, as decisões, o código e a validação pertencem à equipe.

Na implementação, este requisito é tratado como parte do sistema e não como uma tela isolada. A interface orienta a próxima ação; a API confere identidade, papel, estado e limites; a persistência guarda somente o necessário; e as integrações externas ficam condicionadas às variáveis do ambiente. O escopo continua exclusivamente de software, sem depender de escopo físico.

O fluxo começa quando a pessoa encontra uma necessidade ou uma habilidade que deseja compartilhar. Em seguida, ela confirma a intenção, executa a ação principal e recebe uma resposta que pode ser consultada depois. Para nós, nota de autoria e uso de ia só está concluído quando o caminho é compreensível, existe recuperação para erro e o servidor registra uma evidência coerente.

O aceite não fica apenas no texto. A evidência prevista é: Declaração no início do documento e revisão final pelos cinco autores. O critério de aceite deve ser escrito antes do piloto, com entrada, resultado esperado, responsável, data e ferramenta de verificação. Assim, a documentação separa conformidade técnica, observação de uso e hipótese de impacto.

A validação combina teste automatizado, inspeção manual e uma tarefa representativa. Vamos registrar tempo, erros, conclusão, compreensão e barreiras observadas por teclado, leitor de tela, zoom, contraste, texto, áudio ou conexão instável. O objetivo não é criar uma aparência de certeza, e sim manter um ciclo curto entre observar, corrigir e conferir novamente.

Também registramos uma limitação, porque ela faz parte do resultado: A banca pode exigir detalhamento adicional; manter versões e evidências de revisão. O próximo passo é transformar a pendência em issue, teste, decisão de reunião ou evidência de piloto. Este capítulo não promete superioridade nem receita: ele deixa um contrato verificável para que a equipe consiga descobrir se a hipótese se sustenta.

\section{Conferência do release contra o código}
A discussão de conferência do release contra o código começou por uma situação concreta: o material de produto é amplo e poderia fazer a documentação parecer maior do que a implementação realmente entregue.

Nossa escolha para este ponto foi direta: A equipe conferiu o release pelas rotas Java, migrações Flyway, páginas React e contratos de autenticação. O documento classifica cada integração como implementada, opcional por variável de ambiente ou pendente de validação externa.

Na implementação, este requisito é tratado como parte do sistema e não como uma tela isolada. A interface orienta a próxima ação; a API confere identidade, papel, estado e limites; a persistência guarda somente o necessário; e as integrações externas ficam condicionadas às variáveis do ambiente. O escopo continua exclusivamente de software, sem depender de escopo físico.

O fluxo começa quando a pessoa encontra uma necessidade ou uma habilidade que deseja compartilhar. Em seguida, ela confirma a intenção, executa a ação principal e recebe uma resposta que pode ser consultada depois. Para nós, conferência do release contra o código só está concluído quando o caminho é compreensível, existe recuperação para erro e o servidor registra uma evidência coerente.

O aceite não fica apenas no texto. A evidência prevista é: Inventário em docs/CODE\_ALIGNMENT.md, revisão de controllers, frontend/src, V1-V4 e configuração sem segredos. O critério de aceite deve ser escrito antes do piloto, com entrada, resultado esperado, responsável, data e ferramenta de verificação. Assim, a documentação separa conformidade técnica, observação de uso e hipótese de impacto.

A validação combina teste automatizado, inspeção manual e uma tarefa representativa. Vamos registrar tempo, erros, conclusão, compreensão e barreiras observadas por teclado, leitor de tela, zoom, contraste, texto, áudio ou conexão instável. O objetivo não é criar uma aparência de certeza, e sim manter um ciclo curto entre observar, corrigir e conferir novamente.

Também registramos uma limitação, porque ela faz parte do resultado: Compilação e execução integral ainda dependem de banco, Firebase e provedores configurados; isso é declarado como limite, não como resultado já obtido. O próximo passo é transformar a pendência em issue, teste, decisão de reunião ou evidência de piloto. Este capítulo não promete superioridade nem receita: ele deixa um contrato verificável para que a equipe consiga descobrir se a hipótese se sustenta.

\section{Problema de negócio}
Durante a consolidação do MVP, percebemos que pessoas que querem aprender e ensinar encontram conteúdo isolado, pouca reciprocidade e dificuldade para transformar intenção em atividade concluída.

No protótipo, esse entendimento aparece assim: A SkillSwap organiza perfil, descoberta, proposta, conversa, conclusão e avaliação em uma jornada única.

Na implementação, este requisito é tratado como parte do sistema e não como uma tela isolada. A interface orienta a próxima ação; a API confere identidade, papel, estado e limites; a persistência guarda somente o necessário; e as integrações externas ficam condicionadas às variáveis do ambiente. O escopo continua exclusivamente de software, sem depender de escopo físico.

O fluxo começa quando a pessoa encontra uma necessidade ou uma habilidade que deseja compartilhar. Em seguida, ela confirma a intenção, executa a ação principal e recebe uma resposta que pode ser consultada depois. Para nós, problema de negócio só está concluído quando o caminho é compreensível, existe recuperação para erro e o servidor registra uma evidência coerente.

A forma mais honesta de verificar este requisito é registrar: Fluxo funcional do protótipo e indicadores de conversão entre etapas. O critério de aceite deve ser escrito antes do piloto, com entrada, resultado esperado, responsável, data e ferramenta de verificação. Assim, a documentação separa conformidade técnica, observação de uso e hipótese de impacto.

A validação combina teste automatizado, inspeção manual e uma tarefa representativa. Vamos registrar tempo, erros, conclusão, compreensão e barreiras observadas por teclado, leitor de tela, zoom, contraste, texto, áudio ou conexão instável. O objetivo não é criar uma aparência de certeza, e sim manter um ciclo curto entre observar, corrigir e conferir novamente.

A equipe não trata este item como resolvido antes de enfrentar o risco: Uma rede bilateral precisa de massa crítica e curadoria para evitar baixa oferta em habilidades raras. O próximo passo é transformar a pendência em issue, teste, decisão de reunião ou evidência de piloto. Este capítulo não promete superioridade nem receita: ele deixa um contrato verificável para que a equipe consiga descobrir se a hipótese se sustenta.

\section{Problema de inclusão}
Durante a consolidação do MVP, percebemos que interfaces sociais e educacionais frequentemente tratam acessibilidade como recurso posterior e deixam barreiras no fluxo principal.

No protótipo, esse entendimento aparece assim: Acessibilidade será requisito de aceite: teclado, contraste, foco, leitor de tela, voz, texto simples, legendas e alternativa ao VLibras.

Na implementação, este requisito é tratado como parte do sistema e não como uma tela isolada. A interface orienta a próxima ação; a API confere identidade, papel, estado e limites; a persistência guarda somente o necessário; e as integrações externas ficam condicionadas às variáveis do ambiente. O escopo continua exclusivamente de software, sem depender de escopo físico.

O fluxo começa quando a pessoa encontra uma necessidade ou uma habilidade que deseja compartilhar. Em seguida, ela confirma a intenção, executa a ação principal e recebe uma resposta que pode ser consultada depois. Para nós, problema de inclusão só está concluído quando o caminho é compreensível, existe recuperação para erro e o servidor registra uma evidência coerente.

A forma mais honesta de verificar este requisito é registrar: Matriz de acessibilidade manual, axe/WCAG e tarefas com tecnologias assistivas. O critério de aceite deve ser escrito antes do piloto, com entrada, resultado esperado, responsável, data e ferramenta de verificação. Assim, a documentação separa conformidade técnica, observação de uso e hipótese de impacto.

A validação combina teste automatizado, inspeção manual e uma tarefa representativa. Vamos registrar tempo, erros, conclusão, compreensão e barreiras observadas por teclado, leitor de tela, zoom, contraste, texto, áudio ou conexão instável. O objetivo não é criar uma aparência de certeza, e sim manter um ciclo curto entre observar, corrigir e conferir novamente.

A equipe não trata este item como resolvido antes de enfrentar o risco: Conformidade automática não garante experiência real; testes com usuários devem ocorrer antes do piloto. O próximo passo é transformar a pendência em issue, teste, decisão de reunião ou evidência de piloto. Este capítulo não promete superioridade nem receita: ele deixa um contrato verificável para que a equipe consiga descobrir se a hipótese se sustenta.

\section{Pergunta de pesquisa}
A discussão de pergunta de pesquisa começou por uma situação concreta: ainda não basta dizer que a plataforma é inclusiva ou inteligente; é preciso definir o que será medido.

Nossa escolha para este ponto foi direta: Pergunta: uma jornada multimodal e explicável reduz erros e abandono na realização de uma troca de habilidade, sem aumentar o tempo de tarefa?

Na implementação, este requisito é tratado como parte do sistema e não como uma tela isolada. A interface orienta a próxima ação; a API confere identidade, papel, estado e limites; a persistência guarda somente o necessário; e as integrações externas ficam condicionadas às variáveis do ambiente. O escopo continua exclusivamente de software, sem depender de escopo físico.

O fluxo começa quando a pessoa encontra uma necessidade ou uma habilidade que deseja compartilhar. Em seguida, ela confirma a intenção, executa a ação principal e recebe uma resposta que pode ser consultada depois. Para nós, pergunta de pesquisa só está concluído quando o caminho é compreensível, existe recuperação para erro e o servidor registra uma evidência coerente.

O aceite não fica apenas no texto. A evidência prevista é: Tempo, erros, conclusão, SUS, acessibilidade percebida e compreensão da explicação. O critério de aceite deve ser escrito antes do piloto, com entrada, resultado esperado, responsável, data e ferramenta de verificação. Assim, a documentação separa conformidade técnica, observação de uso e hipótese de impacto.

A validação combina teste automatizado, inspeção manual e uma tarefa representativa. Vamos registrar tempo, erros, conclusão, compreensão e barreiras observadas por teclado, leitor de tela, zoom, contraste, texto, áudio ou conexão instável. O objetivo não é criar uma aparência de certeza, e sim manter um ciclo curto entre observar, corrigir e conferir novamente.

Também registramos uma limitação, porque ela faz parte do resultado: Amostra pequena de TCC permite evidência preliminar, não generalização populacional. O próximo passo é transformar a pendência em issue, teste, decisão de reunião ou evidência de piloto. Este capítulo não promete superioridade nem receita: ele deixa um contrato verificável para que a equipe consiga descobrir se a hipótese se sustenta.

\section{Objetivo geral}
Na leitura que fizemos do projeto, o projeto precisava conectar impacto social, engenharia de software e possibilidade de startup.

Depois de comparar as alternativas, a equipe decidiu o seguinte: Projetar e entregar uma plataforma de troca de habilidades inclusiva, auditável e comercialmente testável, sem depender de escopo físico.

Na implementação, este requisito é tratado como parte do sistema e não como uma tela isolada. A interface orienta a próxima ação; a API confere identidade, papel, estado e limites; a persistência guarda somente o necessário; e as integrações externas ficam condicionadas às variáveis do ambiente. O escopo continua exclusivamente de software, sem depender de escopo físico.

O fluxo começa quando a pessoa encontra uma necessidade ou uma habilidade que deseja compartilhar. Em seguida, ela confirma a intenção, executa a ação principal e recebe uma resposta que pode ser consultada depois. Para nós, objetivo geral só está concluído quando o caminho é compreensível, existe recuperação para erro e o servidor registra uma evidência coerente.

Para conferir se a decisão funciona, vamos procurar esta evidência: Release reproduzível, documentação e demonstração do fluxo completo. O critério de aceite deve ser escrito antes do piloto, com entrada, resultado esperado, responsável, data e ferramenta de verificação. Assim, a documentação separa conformidade técnica, observação de uso e hipótese de impacto.

A validação combina teste automatizado, inspeção manual e uma tarefa representativa. Vamos registrar tempo, erros, conclusão, compreensão e barreiras observadas por teclado, leitor de tela, zoom, contraste, texto, áudio ou conexão instável. O objetivo não é criar uma aparência de certeza, e sim manter um ciclo curto entre observar, corrigir e conferir novamente.

O ponto que ainda merece cuidado é o seguinte: O objetivo não promete receita; a hipótese comercial será validada com entrevistas e piloto. O próximo passo é transformar a pendência em issue, teste, decisão de reunião ou evidência de piloto. Este capítulo não promete superioridade nem receita: ele deixa um contrato verificável para que a equipe consiga descobrir se a hipótese se sustenta.

\section{Objetivos específicos}
A discussão de objetivos específicos começou por uma situação concreta: recursos soltos dificultam priorização e avaliação.

Nossa escolha para este ponto foi direta: Mapear requisitos, corrigir segurança, consolidar dados, implementar acessibilidade, explicar recomendações e preparar deploy.

Na implementação, este requisito é tratado como parte do sistema e não como uma tela isolada. A interface orienta a próxima ação; a API confere identidade, papel, estado e limites; a persistência guarda somente o necessário; e as integrações externas ficam condicionadas às variáveis do ambiente. O escopo continua exclusivamente de software, sem depender de escopo físico.

O fluxo começa quando a pessoa encontra uma necessidade ou uma habilidade que deseja compartilhar. Em seguida, ela confirma a intenção, executa a ação principal e recebe uma resposta que pode ser consultada depois. Para nós, objetivos específicos só está concluído quando o caminho é compreensível, existe recuperação para erro e o servidor registra uma evidência coerente.

O aceite não fica apenas no texto. A evidência prevista é: Backlog priorizado, casos de uso, testes e checklist de release. O critério de aceite deve ser escrito antes do piloto, com entrada, resultado esperado, responsável, data e ferramenta de verificação. Assim, a documentação separa conformidade técnica, observação de uso e hipótese de impacto.

A validação combina teste automatizado, inspeção manual e uma tarefa representativa. Vamos registrar tempo, erros, conclusão, compreensão e barreiras observadas por teclado, leitor de tela, zoom, contraste, texto, áudio ou conexão instável. O objetivo não é criar uma aparência de certeza, e sim manter um ciclo curto entre observar, corrigir e conferir novamente.

Também registramos uma limitação, porque ela faz parte do resultado: Escopo grande pode estourar o cronograma; o MVP deve manter o fluxo perfeito. O próximo passo é transformar a pendência em issue, teste, decisão de reunião ou evidência de piloto. Este capítulo não promete superioridade nem receita: ele deixa um contrato verificável para que a equipe consiga descobrir se a hipótese se sustenta.

\section{Diferencial proposto}
A discussão de diferencial proposto começou por uma situação concreta: plataformas de cursos, redes sociais e job boards resolvem partes diferentes do problema.

Nossa escolha para este ponto foi direta: O recorte diferencial é a combinação de troca bilateral, acessibilidade por padrão, recomendação explicável, reputação e ponte entre aprendizagem e oportunidade.

Na implementação, este requisito é tratado como parte do sistema e não como uma tela isolada. A interface orienta a próxima ação; a API confere identidade, papel, estado e limites; a persistência guarda somente o necessário; e as integrações externas ficam condicionadas às variáveis do ambiente. O escopo continua exclusivamente de software, sem depender de escopo físico.

O fluxo começa quando a pessoa encontra uma necessidade ou uma habilidade que deseja compartilhar. Em seguida, ela confirma a intenção, executa a ação principal e recebe uma resposta que pode ser consultada depois. Para nós, diferencial proposto só está concluído quando o caminho é compreensível, existe recuperação para erro e o servidor registra uma evidência coerente.

O aceite não fica apenas no texto. A evidência prevista é: Tabela comparativa e experimento de usabilidade com baseline de fluxo. O critério de aceite deve ser escrito antes do piloto, com entrada, resultado esperado, responsável, data e ferramenta de verificação. Assim, a documentação separa conformidade técnica, observação de uso e hipótese de impacto.

A validação combina teste automatizado, inspeção manual e uma tarefa representativa. Vamos registrar tempo, erros, conclusão, compreensão e barreiras observadas por teclado, leitor de tela, zoom, contraste, texto, áudio ou conexão instável. O objetivo não é criar uma aparência de certeza, e sim manter um ciclo curto entre observar, corrigir e conferir novamente.

Também registramos uma limitação, porque ela faz parte do resultado: Diferencial é hipótese; superioridade só pode ser alegada depois de comparação controlada. O próximo passo é transformar a pendência em issue, teste, decisão de reunião ou evidência de piloto. Este capítulo não promete superioridade nem receita: ele deixa um contrato verificável para que a equipe consiga descobrir se a hipótese se sustenta.

\section{Literatura sobre matching}
Na leitura que fizemos do projeto, encontrar uma pessoa compatível é um problema de duas pontas e não apenas uma busca textual.

Depois de comparar as alternativas, a equipe decidiu o seguinte: Usar atributos de habilidade, disponibilidade, objetivo e preferências para gerar ranking compreensível.

Na implementação, este requisito é tratado como parte do sistema e não como uma tela isolada. A interface orienta a próxima ação; a API confere identidade, papel, estado e limites; a persistência guarda somente o necessário; e as integrações externas ficam condicionadas às variáveis do ambiente. O escopo continua exclusivamente de software, sem depender de escopo físico.

O fluxo começa quando a pessoa encontra uma necessidade ou uma habilidade que deseja compartilhar. Em seguida, ela confirma a intenção, executa a ação principal e recebe uma resposta que pode ser consultada depois. Para nós, literatura sobre matching só está concluído quando o caminho é compreensível, existe recuperação para erro e o servidor registra uma evidência coerente.

Para conferir se a decisão funciona, vamos procurar esta evidência: Chang, Li e Sun (2019) e logs de recomendações explicáveis. O critério de aceite deve ser escrito antes do piloto, com entrada, resultado esperado, responsável, data e ferramenta de verificação. Assim, a documentação separa conformidade técnica, observação de uso e hipótese de impacto.

A validação combina teste automatizado, inspeção manual e uma tarefa representativa. Vamos registrar tempo, erros, conclusão, compreensão e barreiras observadas por teclado, leitor de tela, zoom, contraste, texto, áudio ou conexão instável. O objetivo não é criar uma aparência de certeza, e sim manter um ciclo curto entre observar, corrigir e conferir novamente.

O ponto que ainda merece cuidado é o seguinte: Matching algorítmico pode reforçar popularidade; incluir exploração e controles de exposição. O próximo passo é transformar a pendência em issue, teste, decisão de reunião ou evidência de piloto. Este capítulo não promete superioridade nem receita: ele deixa um contrato verificável para que a equipe consiga descobrir se a hipótese se sustenta.

\section{Literatura sobre aprendizagem adaptativa}
Na leitura que fizemos do projeto, o excesso de cursos e caminhos possíveis dificulta a escolha do próximo passo.

Depois de comparar as alternativas, a equipe decidiu o seguinte: Começar com regras e sinais do domínio antes de modelos complexos; cada recomendação mostra os fatores usados.

Na implementação, este requisito é tratado como parte do sistema e não como uma tela isolada. A interface orienta a próxima ação; a API confere identidade, papel, estado e limites; a persistência guarda somente o necessário; e as integrações externas ficam condicionadas às variáveis do ambiente. O escopo continua exclusivamente de software, sem depender de escopo físico.

O fluxo começa quando a pessoa encontra uma necessidade ou uma habilidade que deseja compartilhar. Em seguida, ela confirma a intenção, executa a ação principal e recebe uma resposta que pode ser consultada depois. Para nós, literatura sobre aprendizagem adaptativa só está concluído quando o caminho é compreensível, existe recuperação para erro e o servidor registra uma evidência coerente.

Para conferir se a decisão funciona, vamos procurar esta evidência: Amin et al. (2023), requisitos de explicação e avaliação offline. O critério de aceite deve ser escrito antes do piloto, com entrada, resultado esperado, responsável, data e ferramenta de verificação. Assim, a documentação separa conformidade técnica, observação de uso e hipótese de impacto.

A validação combina teste automatizado, inspeção manual e uma tarefa representativa. Vamos registrar tempo, erros, conclusão, compreensão e barreiras observadas por teclado, leitor de tela, zoom, contraste, texto, áudio ou conexão instável. O objetivo não é criar uma aparência de certeza, e sim manter um ciclo curto entre observar, corrigir e conferir novamente.

O ponto que ainda merece cuidado é o seguinte: Não usar reforço em produção sem dados e protocolo de segurança; o TCC entrega uma versão explicável. O próximo passo é transformar a pendência em issue, teste, decisão de reunião ou evidência de piloto. Este capítulo não promete superioridade nem receita: ele deixa um contrato verificável para que a equipe consiga descobrir se a hipótese se sustenta.

\section{Literatura sobre recomendação interpretável}
Na leitura que fizemos do projeto, uma recomendação sem justificativa é difícil de aceitar, corrigir ou contestar.

Depois de comparar as alternativas, a equipe decidiu o seguinte: A API retorna fatores positivos, lacunas e alternativas, mantendo o usuário no controle.

Na implementação, este requisito é tratado como parte do sistema e não como uma tela isolada. A interface orienta a próxima ação; a API confere identidade, papel, estado e limites; a persistência guarda somente o necessário; e as integrações externas ficam condicionadas às variáveis do ambiente. O escopo continua exclusivamente de software, sem depender de escopo físico.

O fluxo começa quando a pessoa encontra uma necessidade ou uma habilidade que deseja compartilhar. Em seguida, ela confirma a intenção, executa a ação principal e recebe uma resposta que pode ser consultada depois. Para nós, literatura sobre recomendação interpretável só está concluído quando o caminho é compreensível, existe recuperação para erro e o servidor registra uma evidência coerente.

Para conferir se a decisão funciona, vamos procurar esta evidência: Yun et al. (2022), payload de explicação e teste de compreensão. O critério de aceite deve ser escrito antes do piloto, com entrada, resultado esperado, responsável, data e ferramenta de verificação. Assim, a documentação separa conformidade técnica, observação de uso e hipótese de impacto.

A validação combina teste automatizado, inspeção manual e uma tarefa representativa. Vamos registrar tempo, erros, conclusão, compreensão e barreiras observadas por teclado, leitor de tela, zoom, contraste, texto, áudio ou conexão instável. O objetivo não é criar uma aparência de certeza, e sim manter um ciclo curto entre observar, corrigir e conferir novamente.

O ponto que ainda merece cuidado é o seguinte: Explicações podem ser simplificações; registrar que são razões operacionais, não causalidade científica. O próximo passo é transformar a pendência em issue, teste, decisão de reunião ou evidência de piloto. Este capítulo não promete superioridade nem receita: ele deixa um contrato verificável para que a equipe consiga descobrir se a hipótese se sustenta.

\section{Literatura sobre acessibilidade}
Na leitura que fizemos do projeto, acessibilidade depende de interação e contexto, não apenas de uma auditoria estática.

Depois de comparar as alternativas, a equipe decidiu o seguinte: Projetar tarefas e testar com teclado, NVDA, TalkBack e VoiceOver; registrar limitações de ambiente.

Na implementação, este requisito é tratado como parte do sistema e não como uma tela isolada. A interface orienta a próxima ação; a API confere identidade, papel, estado e limites; a persistência guarda somente o necessário; e as integrações externas ficam condicionadas às variáveis do ambiente. O escopo continua exclusivamente de software, sem depender de escopo físico.

O fluxo começa quando a pessoa encontra uma necessidade ou uma habilidade que deseja compartilhar. Em seguida, ela confirma a intenção, executa a ação principal e recebe uma resposta que pode ser consultada depois. Para nós, literatura sobre acessibilidade só está concluído quando o caminho é compreensível, existe recuperação para erro e o servidor registra uma evidência coerente.

Para conferir se a decisão funciona, vamos procurar esta evidência: Lorgat et al. (2022), Trielsa e Angeline (2023), Kondoro et al. (2023). O critério de aceite deve ser escrito antes do piloto, com entrada, resultado esperado, responsável, data e ferramenta de verificação. Assim, a documentação separa conformidade técnica, observação de uso e hipótese de impacto.

A validação combina teste automatizado, inspeção manual e uma tarefa representativa. Vamos registrar tempo, erros, conclusão, compreensão e barreiras observadas por teclado, leitor de tela, zoom, contraste, texto, áudio ou conexão instável. O objetivo não é criar uma aparência de certeza, e sim manter um ciclo curto entre observar, corrigir e conferir novamente.

O ponto que ainda merece cuidado é o seguinte: O widget externo de Libras não substitui linguagem simples, legendas e validação linguística. O próximo passo é transformar a pendência em issue, teste, decisão de reunião ou evidência de piloto. Este capítulo não promete superioridade nem receita: ele deixa um contrato verificável para que a equipe consiga descobrir se a hipótese se sustenta.

\section{Literatura sobre aprendizagem e emprego}
Na leitura que fizemos do projeto, formação e oportunidade profissional ficam separadas em produtos diferentes.

Depois de comparar as alternativas, a equipe decidiu o seguinte: Registrar habilidades demonstradas, projetos, mentoria e vagas no mesmo perfil, com painel institucional.

Na implementação, este requisito é tratado como parte do sistema e não como uma tela isolada. A interface orienta a próxima ação; a API confere identidade, papel, estado e limites; a persistência guarda somente o necessário; e as integrações externas ficam condicionadas às variáveis do ambiente. O escopo continua exclusivamente de software, sem depender de escopo físico.

O fluxo começa quando a pessoa encontra uma necessidade ou uma habilidade que deseja compartilhar. Em seguida, ela confirma a intenção, executa a ação principal e recebe uma resposta que pode ser consultada depois. Para nós, literatura sobre aprendizagem e emprego só está concluído quando o caminho é compreensível, existe recuperação para erro e o servidor registra uma evidência coerente.

Para conferir se a decisão funciona, vamos procurar esta evidência: Srivastava et al. (2022), Jhalani et al. (2024) e piloto com instituição. O critério de aceite deve ser escrito antes do piloto, com entrada, resultado esperado, responsável, data e ferramenta de verificação. Assim, a documentação separa conformidade técnica, observação de uso e hipótese de impacto.

A validação combina teste automatizado, inspeção manual e uma tarefa representativa. Vamos registrar tempo, erros, conclusão, compreensão e barreiras observadas por teclado, leitor de tela, zoom, contraste, texto, áudio ou conexão instável. O objetivo não é criar uma aparência de certeza, e sim manter um ciclo curto entre observar, corrigir e conferir novamente.

O ponto que ainda merece cuidado é o seguinte: A plataforma não deve decidir contratação; ela organiza evidências e preferências. O próximo passo é transformar a pendência em issue, teste, decisão de reunião ou evidência de piloto. Este capítulo não promete superioridade nem receita: ele deixa um contrato verificável para que a equipe consiga descobrir se a hipótese se sustenta.

\section{Modelo B2B2C}
Na leitura que fizemos do projeto, usuários finais podem não pagar diretamente, enquanto instituições precisam de métricas de inclusão e capacitação.

Depois de comparar as alternativas, a equipe decidiu o seguinte: Oferecer implantação, painéis, trilhas institucionais e relatórios para universidades, empresas e ONGs, mantendo uso individual.

Na implementação, este requisito é tratado como parte do sistema e não como uma tela isolada. A interface orienta a próxima ação; a API confere identidade, papel, estado e limites; a persistência guarda somente o necessário; e as integrações externas ficam condicionadas às variáveis do ambiente. O escopo continua exclusivamente de software, sem depender de escopo físico.

O fluxo começa quando a pessoa encontra uma necessidade ou uma habilidade que deseja compartilhar. Em seguida, ela confirma a intenção, executa a ação principal e recebe uma resposta que pode ser consultada depois. Para nós, modelo b2b2c só está concluído quando o caminho é compreensível, existe recuperação para erro e o servidor registra uma evidência coerente.

Para conferir se a decisão funciona, vamos procurar esta evidência: Entrevistas, carta de proposta e métrica de ativação institucional. O critério de aceite deve ser escrito antes do piloto, com entrada, resultado esperado, responsável, data e ferramenta de verificação. Assim, a documentação separa conformidade técnica, observação de uso e hipótese de impacto.

A validação combina teste automatizado, inspeção manual e uma tarefa representativa. Vamos registrar tempo, erros, conclusão, compreensão e barreiras observadas por teclado, leitor de tela, zoom, contraste, texto, áudio ou conexão instável. O objetivo não é criar uma aparência de certeza, e sim manter um ciclo curto entre observar, corrigir e conferir novamente.

O ponto que ainda merece cuidado é o seguinte: Venda corporativa exige governança, SLA e proteção de dados além do MVP. O próximo passo é transformar a pendência em issue, teste, decisão de reunião ou evidência de piloto. Este capítulo não promete superioridade nem receita: ele deixa um contrato verificável para que a equipe consiga descobrir se a hipótese se sustenta.

\section{Personas inclusivas}
Durante a consolidação do MVP, percebemos que um usuário que usa leitor de tela tem necessidades diferentes de uma pessoa que prefere áudio ou texto simples.

No protótipo, esse entendimento aparece assim: Documentar quatro personas e suas jornadas, sem reduzir deficiência a um único perfil.

Na implementação, este requisito é tratado como parte do sistema e não como uma tela isolada. A interface orienta a próxima ação; a API confere identidade, papel, estado e limites; a persistência guarda somente o necessário; e as integrações externas ficam condicionadas às variáveis do ambiente. O escopo continua exclusivamente de software, sem depender de escopo físico.

O fluxo começa quando a pessoa encontra uma necessidade ou uma habilidade que deseja compartilhar. Em seguida, ela confirma a intenção, executa a ação principal e recebe uma resposta que pode ser consultada depois. Para nós, personas inclusivas só está concluído quando o caminho é compreensível, existe recuperação para erro e o servidor registra uma evidência coerente.

A forma mais honesta de verificar este requisito é registrar: Roteiros de tarefa, critérios de sucesso e feedback qualitativo. O critério de aceite deve ser escrito antes do piloto, com entrada, resultado esperado, responsável, data e ferramenta de verificação. Assim, a documentação separa conformidade técnica, observação de uso e hipótese de impacto.

A validação combina teste automatizado, inspeção manual e uma tarefa representativa. Vamos registrar tempo, erros, conclusão, compreensão e barreiras observadas por teclado, leitor de tela, zoom, contraste, texto, áudio ou conexão instável. O objetivo não é criar uma aparência de certeza, e sim manter um ciclo curto entre observar, corrigir e conferir novamente.

A equipe não trata este item como resolvido antes de enfrentar o risco: Personas são abstrações; não substituir participantes reais por suposições. O próximo passo é transformar a pendência em issue, teste, decisão de reunião ou evidência de piloto. Este capítulo não promete superioridade nem receita: ele deixa um contrato verificável para que a equipe consiga descobrir se a hipótese se sustenta.

\section{Jornada principal}
Durante a consolidação do MVP, percebemos que a ideia só se prova quando a pessoa chega de intenção à troca concluída.

No protótipo, esse entendimento aparece assim: Entrada, perfil, recomendação, proposta, chat, agenda, entrega, avaliação e histórico formam o happy path.

Na implementação, este requisito é tratado como parte do sistema e não como uma tela isolada. A interface orienta a próxima ação; a API confere identidade, papel, estado e limites; a persistência guarda somente o necessário; e as integrações externas ficam condicionadas às variáveis do ambiente. O escopo continua exclusivamente de software, sem depender de escopo físico.

O fluxo começa quando a pessoa encontra uma necessidade ou uma habilidade que deseja compartilhar. Em seguida, ela confirma a intenção, executa a ação principal e recebe uma resposta que pode ser consultada depois. Para nós, jornada principal só está concluído quando o caminho é compreensível, existe recuperação para erro e o servidor registra uma evidência coerente.

A forma mais honesta de verificar este requisito é registrar: Teste ponta a ponta com dados sintéticos e gravação sanitizada. O critério de aceite deve ser escrito antes do piloto, com entrada, resultado esperado, responsável, data e ferramenta de verificação. Assim, a documentação separa conformidade técnica, observação de uso e hipótese de impacto.

A validação combina teste automatizado, inspeção manual e uma tarefa representativa. Vamos registrar tempo, erros, conclusão, compreensão e barreiras observadas por teclado, leitor de tela, zoom, contraste, texto, áudio ou conexão instável. O objetivo não é criar uma aparência de certeza, e sim manter um ciclo curto entre observar, corrigir e conferir novamente.

A equipe não trata este item como resolvido antes de enfrentar o risco: Integrações externas podem falhar; sempre oferecer fallback e mensagem acionável. O próximo passo é transformar a pendência em issue, teste, decisão de reunião ou evidência de piloto. Este capítulo não promete superioridade nem receita: ele deixa um contrato verificável para que a equipe consiga descobrir se a hipótese se sustenta.

\section{Jornada administrativa}
A discussão de jornada administrativa começou por uma situação concreta: instituições precisam cuidar de segurança, moderação, indicadores e reconciliação.

Nossa escolha para este ponto foi direta: Console separado por papel, endpoints protegidos, auditoria de ações e dados agregados.

Na implementação, este requisito é tratado como parte do sistema e não como uma tela isolada. A interface orienta a próxima ação; a API confere identidade, papel, estado e limites; a persistência guarda somente o necessário; e as integrações externas ficam condicionadas às variáveis do ambiente. O escopo continua exclusivamente de software, sem depender de escopo físico.

O fluxo começa quando a pessoa encontra uma necessidade ou uma habilidade que deseja compartilhar. Em seguida, ela confirma a intenção, executa a ação principal e recebe uma resposta que pode ser consultada depois. Para nós, jornada administrativa só está concluído quando o caminho é compreensível, existe recuperação para erro e o servidor registra uma evidência coerente.

O aceite não fica apenas no texto. A evidência prevista é: Testes 403/200, trilha de auditoria e relatório operacional. O critério de aceite deve ser escrito antes do piloto, com entrada, resultado esperado, responsável, data e ferramenta de verificação. Assim, a documentação separa conformidade técnica, observação de uso e hipótese de impacto.

A validação combina teste automatizado, inspeção manual e uma tarefa representativa. Vamos registrar tempo, erros, conclusão, compreensão e barreiras observadas por teclado, leitor de tela, zoom, contraste, texto, áudio ou conexão instável. O objetivo não é criar uma aparência de certeza, e sim manter um ciclo curto entre observar, corrigir e conferir novamente.

Também registramos uma limitação, porque ela faz parte do resultado: Privilégios excessivos ou dados pessoais no painel aumentam risco LGPD. O próximo passo é transformar a pendência em issue, teste, decisão de reunião ou evidência de piloto. Este capítulo não promete superioridade nem receita: ele deixa um contrato verificável para que a equipe consiga descobrir se a hipótese se sustenta.

\section{Arquitetura modular}
A discussão de arquitetura modular começou por uma situação concreta: microsserviços prematuros aumentariam custo e dificultariam a demonstração.

Nossa escolha para este ponto foi direta: Monólito modular Spring Boot com módulos de identidade, comunidade, aprendizagem, notificações e administração.

Na implementação, este requisito é tratado como parte do sistema e não como uma tela isolada. A interface orienta a próxima ação; a API confere identidade, papel, estado e limites; a persistência guarda somente o necessário; e as integrações externas ficam condicionadas às variáveis do ambiente. O escopo continua exclusivamente de software, sem depender de escopo físico.

O fluxo começa quando a pessoa encontra uma necessidade ou uma habilidade que deseja compartilhar. Em seguida, ela confirma a intenção, executa a ação principal e recebe uma resposta que pode ser consultada depois. Para nós, arquitetura modular só está concluído quando o caminho é compreensível, existe recuperação para erro e o servidor registra uma evidência coerente.

O aceite não fica apenas no texto. A evidência prevista é: Pacotes Java, limites de serviço e documentação de dependências. O critério de aceite deve ser escrito antes do piloto, com entrada, resultado esperado, responsável, data e ferramenta de verificação. Assim, a documentação separa conformidade técnica, observação de uso e hipótese de impacto.

A validação combina teste automatizado, inspeção manual e uma tarefa representativa. Vamos registrar tempo, erros, conclusão, compreensão e barreiras observadas por teclado, leitor de tela, zoom, contraste, texto, áudio ou conexão instável. O objetivo não é criar uma aparência de certeza, e sim manter um ciclo curto entre observar, corrigir e conferir novamente.

Também registramos uma limitação, porque ela faz parte do resultado: Modularidade pode degradar se regras atravessarem camadas; aplicar revisão arquitetural. O próximo passo é transformar a pendência em issue, teste, decisão de reunião ou evidência de piloto. Este capítulo não promete superioridade nem receita: ele deixa um contrato verificável para que a equipe consiga descobrir se a hipótese se sustenta.

\begin{figure}[ht]
\centering
\fbox{\begin{minipage}{0.82\textwidth}
\centering
\textbf{Usuários e tecnologias assistivas}\\[4pt]
$\downarrow$\\[4pt]
React, TypeScript e Vite\\[4pt]
$\downarrow$\\[4pt]
Firebase Authentication e Spring Boot\\[4pt]
$\downarrow$\\[4pt]
PostgreSQL, Flyway, notificações e monitoramento
\end{minipage}}
\caption{Arquitetura exclusivamente de software da SkillSwap.}
\label{fig:arquitetura}
\end{figure}
\section{Frontend React}
Na leitura que fizemos do projeto, a interface precisa ser reativa, acessível e simples de evoluir.

Depois de comparar as alternativas, a equipe decidiu o seguinte: React/Vite/TypeScript com layout responsivo, componentes semânticos, React Query e adaptadores de API.

Na implementação, este requisito é tratado como parte do sistema e não como uma tela isolada. A interface orienta a próxima ação; a API confere identidade, papel, estado e limites; a persistência guarda somente o necessário; e as integrações externas ficam condicionadas às variáveis do ambiente. O escopo continua exclusivamente de software, sem depender de escopo físico.

O fluxo começa quando a pessoa encontra uma necessidade ou uma habilidade que deseja compartilhar. Em seguida, ela confirma a intenção, executa a ação principal e recebe uma resposta que pode ser consultada depois. Para nós, frontend react só está concluído quando o caminho é compreensível, existe recuperação para erro e o servidor registra uma evidência coerente.

Para conferir se a decisão funciona, vamos procurar esta evidência: Build, lint, testes Vitest e inspeção de navegação. O critério de aceite deve ser escrito antes do piloto, com entrada, resultado esperado, responsável, data e ferramenta de verificação. Assim, a documentação separa conformidade técnica, observação de uso e hipótese de impacto.

A validação combina teste automatizado, inspeção manual e uma tarefa representativa. Vamos registrar tempo, erros, conclusão, compreensão e barreiras observadas por teclado, leitor de tela, zoom, contraste, texto, áudio ou conexão instável. O objetivo não é criar uma aparência de certeza, e sim manter um ciclo curto entre observar, corrigir e conferir novamente.

O ponto que ainda merece cuidado é o seguinte: Mocks antigos podem mascarar contratos; endpoints reais devem ser smoke-tested. O próximo passo é transformar a pendência em issue, teste, decisão de reunião ou evidência de piloto. Este capítulo não promete superioridade nem receita: ele deixa um contrato verificável para que a equipe consiga descobrir se a hipótese se sustenta.

\section{Backend Spring Boot}
Durante a consolidação do MVP, percebemos que a API concentra autorização, transação, validação e integrações.

No protótipo, esse entendimento aparece assim: Spring Security Resource Server valida Firebase; serviços transacionais protegem regras e DTOs evitam exposição de entidades.

Na implementação, este requisito é tratado como parte do sistema e não como uma tela isolada. A interface orienta a próxima ação; a API confere identidade, papel, estado e limites; a persistência guarda somente o necessário; e as integrações externas ficam condicionadas às variáveis do ambiente. O escopo continua exclusivamente de software, sem depender de escopo físico.

O fluxo começa quando a pessoa encontra uma necessidade ou uma habilidade que deseja compartilhar. Em seguida, ela confirma a intenção, executa a ação principal e recebe uma resposta que pode ser consultada depois. Para nós, backend spring boot só está concluído quando o caminho é compreensível, existe recuperação para erro e o servidor registra uma evidência coerente.

A forma mais honesta de verificar este requisito é registrar: Testes de contexto, autorização e contratos HTTP. O critério de aceite deve ser escrito antes do piloto, com entrada, resultado esperado, responsável, data e ferramenta de verificação. Assim, a documentação separa conformidade técnica, observação de uso e hipótese de impacto.

A validação combina teste automatizado, inspeção manual e uma tarefa representativa. Vamos registrar tempo, erros, conclusão, compreensão e barreiras observadas por teclado, leitor de tela, zoom, contraste, texto, áudio ou conexão instável. O objetivo não é criar uma aparência de certeza, e sim manter um ciclo curto entre observar, corrigir e conferir novamente.

A equipe não trata este item como resolvido antes de enfrentar o risco: Falhas de configuração em produção podem bloquear login; manter health check e secrets validáveis. O próximo passo é transformar a pendência em issue, teste, decisão de reunião ou evidência de piloto. Este capítulo não promete superioridade nem receita: ele deixa um contrato verificável para que a equipe consiga descobrir se a hipótese se sustenta.

\section{PostgreSQL e Flyway}
Durante a consolidação do MVP, percebemos que dDL/DML manual causa divergência entre máquinas e deploys.

No protótipo, esse entendimento aparece assim: Flyway é a fonte oficial; Hibernate usa ddl-auto=validate e scripts legados não iniciam o banco.

Na implementação, este requisito é tratado como parte do sistema e não como uma tela isolada. A interface orienta a próxima ação; a API confere identidade, papel, estado e limites; a persistência guarda somente o necessário; e as integrações externas ficam condicionadas às variáveis do ambiente. O escopo continua exclusivamente de software, sem depender de escopo físico.

O fluxo começa quando a pessoa encontra uma necessidade ou uma habilidade que deseja compartilhar. Em seguida, ela confirma a intenção, executa a ação principal e recebe uma resposta que pode ser consultada depois. Para nós, postgresql e flyway só está concluído quando o caminho é compreensível, existe recuperação para erro e o servidor registra uma evidência coerente.

A forma mais honesta de verificar este requisito é registrar: V1-V4, checksum de migração e teste em banco vazio. O critério de aceite deve ser escrito antes do piloto, com entrada, resultado esperado, responsável, data e ferramenta de verificação. Assim, a documentação separa conformidade técnica, observação de uso e hipótese de impacto.

A validação combina teste automatizado, inspeção manual e uma tarefa representativa. Vamos registrar tempo, erros, conclusão, compreensão e barreiras observadas por teclado, leitor de tela, zoom, contraste, texto, áudio ou conexão instável. O objetivo não é criar uma aparência de certeza, e sim manter um ciclo curto entre observar, corrigir e conferir novamente.

A equipe não trata este item como resolvido antes de enfrentar o risco: Migrar banco existente exige backup, inspeção e migration incremental; nunca aplicar baseline cegamente. O próximo passo é transformar a pendência em issue, teste, decisão de reunião ou evidência de piloto. Este capítulo não promete superioridade nem receita: ele deixa um contrato verificável para que a equipe consiga descobrir se a hipótese se sustenta.

\section{Identidade Firebase}
Na leitura que fizemos do projeto, senha e provedor social não devem ser duplicados no backend.

Depois de comparar as alternativas, a equipe decidiu o seguinte: Firebase emite ID token; API valida issuer e sincroniza usuário local por e-mail verificado.

Na implementação, este requisito é tratado como parte do sistema e não como uma tela isolada. A interface orienta a próxima ação; a API confere identidade, papel, estado e limites; a persistência guarda somente o necessário; e as integrações externas ficam condicionadas às variáveis do ambiente. O escopo continua exclusivamente de software, sem depender de escopo físico.

O fluxo começa quando a pessoa encontra uma necessidade ou uma habilidade que deseja compartilhar. Em seguida, ela confirma a intenção, executa a ação principal e recebe uma resposta que pode ser consultada depois. Para nós, identidade firebase só está concluído quando o caminho é compreensível, existe recuperação para erro e o servidor registra uma evidência coerente.

Para conferir se a decisão funciona, vamos procurar esta evidência: Teste de token expirado, usuário ausente, banimento e papel ADMIN. O critério de aceite deve ser escrito antes do piloto, com entrada, resultado esperado, responsável, data e ferramenta de verificação. Assim, a documentação separa conformidade técnica, observação de uso e hipótese de impacto.

A validação combina teste automatizado, inspeção manual e uma tarefa representativa. Vamos registrar tempo, erros, conclusão, compreensão e barreiras observadas por teclado, leitor de tela, zoom, contraste, texto, áudio ou conexão instável. O objetivo não é criar uma aparência de certeza, e sim manter um ciclo curto entre observar, corrigir e conferir novamente.

O ponto que ainda merece cuidado é o seguinte: E-mail como vínculo requer política para alteração de e-mail e contas duplicadas. O próximo passo é transformar a pendência em issue, teste, decisão de reunião ou evidência de piloto. Este capítulo não promete superioridade nem receita: ele deixa um contrato verificável para que a equipe consiga descobrir se a hipótese se sustenta.

\section{Rotação de credenciais}
A discussão de rotação de credenciais começou por uma situação concreta: os arquivos enviados continham credenciais reais de banco e serviços.

Nossa escolha para este ponto foi direta: Entrega só contém placeholders e variáveis privadas; a equipe deve revogar e substituir credenciais antes do deploy.

Na implementação, este requisito é tratado como parte do sistema e não como uma tela isolada. A interface orienta a próxima ação; a API confere identidade, papel, estado e limites; a persistência guarda somente o necessário; e as integrações externas ficam condicionadas às variáveis do ambiente. O escopo continua exclusivamente de software, sem depender de escopo físico.

O fluxo começa quando a pessoa encontra uma necessidade ou uma habilidade que deseja compartilhar. Em seguida, ela confirma a intenção, executa a ação principal e recebe uma resposta que pode ser consultada depois. Para nós, rotação de credenciais só está concluído quando o caminho é compreensível, existe recuperação para erro e o servidor registra uma evidência coerente.

O aceite não fica apenas no texto. A evidência prevista é: Varredura de segredos, histórico do Git e checklist de Railway Secret Variables. O critério de aceite deve ser escrito antes do piloto, com entrada, resultado esperado, responsável, data e ferramenta de verificação. Assim, a documentação separa conformidade técnica, observação de uso e hipótese de impacto.

A validação combina teste automatizado, inspeção manual e uma tarefa representativa. Vamos registrar tempo, erros, conclusão, compreensão e barreiras observadas por teclado, leitor de tela, zoom, contraste, texto, áudio ou conexão instável. O objetivo não é criar uma aparência de certeza, e sim manter um ciclo curto entre observar, corrigir e conferir novamente.

Também registramos uma limitação, porque ela faz parte do resultado: Não é possível rotacionar conta externa sem acesso; o risco permanece aberto até a execução manual. O próximo passo é transformar a pendência em issue, teste, decisão de reunião ou evidência de piloto. Este capítulo não promete superioridade nem receita: ele deixa um contrato verificável para que a equipe consiga descobrir se a hipótese se sustenta.

\section{Sessão sem localStorage}
A discussão de sessão sem localstorage começou por uma situação concreta: token JWT no localStorage é exposto a scripts XSS e o UUID não deve ser fonte de identidade.

Nossa escolha para este ponto foi direta: Firebase gerencia sessão; Axios obtém ID token fresco do usuário atual e páginas buscam perfil por API.

Na implementação, este requisito é tratado como parte do sistema e não como uma tela isolada. A interface orienta a próxima ação; a API confere identidade, papel, estado e limites; a persistência guarda somente o necessário; e as integrações externas ficam condicionadas às variáveis do ambiente. O escopo continua exclusivamente de software, sem depender de escopo físico.

O fluxo começa quando a pessoa encontra uma necessidade ou uma habilidade que deseja compartilhar. Em seguida, ela confirma a intenção, executa a ação principal e recebe uma resposta que pode ser consultada depois. Para nós, sessão sem localstorage só está concluído quando o caminho é compreensível, existe recuperação para erro e o servidor registra uma evidência coerente.

O aceite não fica apenas no texto. A evidência prevista é: Busca por localStorage limitada a preferências visuais e inspeção de bundle. O critério de aceite deve ser escrito antes do piloto, com entrada, resultado esperado, responsável, data e ferramenta de verificação. Assim, a documentação separa conformidade técnica, observação de uso e hipótese de impacto.

A validação combina teste automatizado, inspeção manual e uma tarefa representativa. Vamos registrar tempo, erros, conclusão, compreensão e barreiras observadas por teclado, leitor de tela, zoom, contraste, texto, áudio ou conexão instável. O objetivo não é criar uma aparência de certeza, e sim manter um ciclo curto entre observar, corrigir e conferir novamente.

Também registramos uma limitação, porque ela faz parte do resultado: A persistência interna do Firebase ainda deve seguir política de CSP e origem confiável. O próximo passo é transformar a pendência em issue, teste, decisão de reunião ou evidência de piloto. Este capítulo não promete superioridade nem receita: ele deixa um contrato verificável para que a equipe consiga descobrir se a hipótese se sustenta.

\section{WebSocket sem token na URL}
Durante a consolidação do MVP, percebemos que query string é registrada em proxy, histórico e observabilidade.

No protótipo, esse entendimento aparece assim: Handshake aceita Authorization para clientes que suportam cabeçalho; navegador envia primeiro frame auth; handler valida antes de qualquer evento.

Na implementação, este requisito é tratado como parte do sistema e não como uma tela isolada. A interface orienta a próxima ação; a API confere identidade, papel, estado e limites; a persistência guarda somente o necessário; e as integrações externas ficam condicionadas às variáveis do ambiente. O escopo continua exclusivamente de software, sem depender de escopo físico.

O fluxo começa quando a pessoa encontra uma necessidade ou uma habilidade que deseja compartilhar. Em seguida, ela confirma a intenção, executa a ação principal e recebe uma resposta que pode ser consultada depois. Para nós, websocket sem token na url só está concluído quando o caminho é compreensível, existe recuperação para erro e o servidor registra uma evidência coerente.

A forma mais honesta de verificar este requisito é registrar: Teste sem auth, token inválido, token divergente e mensagem com senderId falso. O critério de aceite deve ser escrito antes do piloto, com entrada, resultado esperado, responsável, data e ferramenta de verificação. Assim, a documentação separa conformidade técnica, observação de uso e hipótese de impacto.

A validação combina teste automatizado, inspeção manual e uma tarefa representativa. Vamos registrar tempo, erros, conclusão, compreensão e barreiras observadas por teclado, leitor de tela, zoom, contraste, texto, áudio ou conexão instável. O objetivo não é criar uma aparência de certeza, e sim manter um ciclo curto entre observar, corrigir e conferir novamente.

A equipe não trata este item como resolvido antes de enfrentar o risco: A rota de handshake precisa permanecer pública só até o primeiro frame; timeout futuro é recomendado. O próximo passo é transformar a pendência em issue, teste, decisão de reunião ou evidência de piloto. Este capítulo não promete superioridade nem receita: ele deixa um contrato verificável para que a equipe consiga descobrir se a hipótese se sustenta.

\section{FCM autenticado}
Durante a consolidação do MVP, percebemos que o token de push precisa ser associado ao usuário autenticado, não recebido com uma identidade fornecida pelo cliente.

No protótipo, esse entendimento aparece assim: POST /api/devices/fcm-token recebe apenas o token e usa o usuário do JWT; DELETE revoga.

Na implementação, este requisito é tratado como parte do sistema e não como uma tela isolada. A interface orienta a próxima ação; a API confere identidade, papel, estado e limites; a persistência guarda somente o necessário; e as integrações externas ficam condicionadas às variáveis do ambiente. O escopo continua exclusivamente de software, sem depender de escopo físico.

O fluxo começa quando a pessoa encontra uma necessidade ou uma habilidade que deseja compartilhar. Em seguida, ela confirma a intenção, executa a ação principal e recebe uma resposta que pode ser consultada depois. Para nós, fcm autenticado só está concluído quando o caminho é compreensível, existe recuperação para erro e o servidor registra uma evidência coerente.

A forma mais honesta de verificar este requisito é registrar: Teste 204, 400, usuário ausente e token acima do limite. O critério de aceite deve ser escrito antes do piloto, com entrada, resultado esperado, responsável, data e ferramenta de verificação. Assim, a documentação separa conformidade técnica, observação de uso e hipótese de impacto.

A validação combina teste automatizado, inspeção manual e uma tarefa representativa. Vamos registrar tempo, erros, conclusão, compreensão e barreiras observadas por teclado, leitor de tela, zoom, contraste, texto, áudio ou conexão instável. O objetivo não é criar uma aparência de certeza, e sim manter um ciclo curto entre observar, corrigir e conferir novamente.

A equipe não trata este item como resolvido antes de enfrentar o risco: FCM multi-dispositivo pode exigir tabela própria na próxima versão. O próximo passo é transformar a pendência em issue, teste, decisão de reunião ou evidência de piloto. Este capítulo não promete superioridade nem receita: ele deixa um contrato verificável para que a equipe consiga descobrir se a hipótese se sustenta.

\section{Integrações comerciais desativadas}

A SkillSwap foi definida como uma plataforma gratuita para os usuários. Por esse motivo, Stripe, SwapCoins, carteira, gorjetas, recursos premium, limites pagos e cobrança por aulas não fazem parte do fluxo entregue. A tela financeira não aparece na navegação e as rotas comerciais são bloqueadas pela segurança da API.

As classes antigas desses módulos permanecem somente como registro da evolução do projeto. Elas não são chamadas pelo frontend e não podem ser utilizadas por uma conta comum. O acesso à plataforma não depende de compra, crédito ou assinatura.
\section{Notificações}
Durante a consolidação do MVP, percebemos que atividade social e propostas podem ser perdidas sem histórico e push.

No protótipo, esse entendimento aparece assim: Persistir notificação, expor unread/read e enviar FCM quando configuração estiver habilitada.

Na implementação, este requisito é tratado como parte do sistema e não como uma tela isolada. A interface orienta a próxima ação; a API confere identidade, papel, estado e limites; a persistência guarda somente o necessário; e as integrações externas ficam condicionadas às variáveis do ambiente. O escopo continua exclusivamente de software, sem depender de escopo físico.

O fluxo começa quando a pessoa encontra uma necessidade ou uma habilidade que deseja compartilhar. Em seguida, ela confirma a intenção, executa a ação principal e recebe uma resposta que pode ser consultada depois. Para nós, notificações só está concluído quando o caminho é compreensível, existe recuperação para erro e o servidor registra uma evidência coerente.

A forma mais honesta de verificar este requisito é registrar: Teste de criação, leitura, autorização por usuário e falha silenciosa de provedor. O critério de aceite deve ser escrito antes do piloto, com entrada, resultado esperado, responsável, data e ferramenta de verificação. Assim, a documentação separa conformidade técnica, observação de uso e hipótese de impacto.

A validação combina teste automatizado, inspeção manual e uma tarefa representativa. Vamos registrar tempo, erros, conclusão, compreensão e barreiras observadas por teclado, leitor de tela, zoom, contraste, texto, áudio ou conexão instável. O objetivo não é criar uma aparência de certeza, e sim manter um ciclo curto entre observar, corrigir e conferir novamente.

A equipe não trata este item como resolvido antes de enfrentar o risco: Nunca registrar tokens ou conteúdo sensível em logs. O próximo passo é transformar a pendência em issue, teste, decisão de reunião ou evidência de piloto. Este capítulo não promete superioridade nem receita: ele deixa um contrato verificável para que a equipe consiga descobrir se a hipótese se sustenta.

\section{Acessibilidade visual}
Na leitura que fizemos do projeto, contraste, zoom e foco ruim impedem a conclusão do fluxo.

Depois de comparar as alternativas, a equipe decidiu o seguinte: Usar semântica, estados de foco, contraste alto, tamanho de fonte e modo daltônico como preferências não sensíveis.

Na implementação, este requisito é tratado como parte do sistema e não como uma tela isolada. A interface orienta a próxima ação; a API confere identidade, papel, estado e limites; a persistência guarda somente o necessário; e as integrações externas ficam condicionadas às variáveis do ambiente. O escopo continua exclusivamente de software, sem depender de escopo físico.

O fluxo começa quando a pessoa encontra uma necessidade ou uma habilidade que deseja compartilhar. Em seguida, ela confirma a intenção, executa a ação principal e recebe uma resposta que pode ser consultada depois. Para nós, acessibilidade visual só está concluído quando o caminho é compreensível, existe recuperação para erro e o servidor registra uma evidência coerente.

Para conferir se a decisão funciona, vamos procurar esta evidência: Matriz WCAG, teclado e screenshot de estados. O critério de aceite deve ser escrito antes do piloto, com entrada, resultado esperado, responsável, data e ferramenta de verificação. Assim, a documentação separa conformidade técnica, observação de uso e hipótese de impacto.

A validação combina teste automatizado, inspeção manual e uma tarefa representativa. Vamos registrar tempo, erros, conclusão, compreensão e barreiras observadas por teclado, leitor de tela, zoom, contraste, texto, áudio ou conexão instável. O objetivo não é criar uma aparência de certeza, e sim manter um ciclo curto entre observar, corrigir e conferir novamente.

O ponto que ainda merece cuidado é o seguinte: Preferência local não substitui persistência no perfil quando o usuário troca de dispositivo. O próximo passo é transformar a pendência em issue, teste, decisão de reunião ou evidência de piloto. Este capítulo não promete superioridade nem receita: ele deixa um contrato verificável para que a equipe consiga descobrir se a hipótese se sustenta.

\section{Acessibilidade auditiva}
Na leitura que fizemos do projeto, vídeos e mensagens sem texto alternativo excluem pessoas surdas.

Depois de comparar as alternativas, a equipe decidiu o seguinte: Legendas, transcrição, textos de apoio, ícones com label e Libras como opção complementar.

Na implementação, este requisito é tratado como parte do sistema e não como uma tela isolada. A interface orienta a próxima ação; a API confere identidade, papel, estado e limites; a persistência guarda somente o necessário; e as integrações externas ficam condicionadas às variáveis do ambiente. O escopo continua exclusivamente de software, sem depender de escopo físico.

O fluxo começa quando a pessoa encontra uma necessidade ou uma habilidade que deseja compartilhar. Em seguida, ela confirma a intenção, executa a ação principal e recebe uma resposta que pode ser consultada depois. Para nós, acessibilidade auditiva só está concluído quando o caminho é compreensível, existe recuperação para erro e o servidor registra uma evidência coerente.

Para conferir se a decisão funciona, vamos procurar esta evidência: Tarefas sem áudio e avaliação de clareza das legendas. O critério de aceite deve ser escrito antes do piloto, com entrada, resultado esperado, responsável, data e ferramenta de verificação. Assim, a documentação separa conformidade técnica, observação de uso e hipótese de impacto.

A validação combina teste automatizado, inspeção manual e uma tarefa representativa. Vamos registrar tempo, erros, conclusão, compreensão e barreiras observadas por teclado, leitor de tela, zoom, contraste, texto, áudio ou conexão instável. O objetivo não é criar uma aparência de certeza, e sim manter um ciclo curto entre observar, corrigir e conferir novamente.

O ponto que ainda merece cuidado é o seguinte: Legendas automáticas precisam de revisão humana quando o conteúdo for institucional. O próximo passo é transformar a pendência em issue, teste, decisão de reunião ou evidência de piloto. Este capítulo não promete superioridade nem receita: ele deixa um contrato verificável para que a equipe consiga descobrir se a hipótese se sustenta.

\section{Acessibilidade cognitiva}
A discussão de acessibilidade cognitiva começou por uma situação concreta: interfaces densas e termos técnicos aumentam erro e abandono.

Nossa escolha para este ponto foi direta: Linguagem simples, hierarquia, confirmação antes de ações irreversíveis e mensagens de recuperação.

Na implementação, este requisito é tratado como parte do sistema e não como uma tela isolada. A interface orienta a próxima ação; a API confere identidade, papel, estado e limites; a persistência guarda somente o necessário; e as integrações externas ficam condicionadas às variáveis do ambiente. O escopo continua exclusivamente de software, sem depender de escopo físico.

O fluxo começa quando a pessoa encontra uma necessidade ou uma habilidade que deseja compartilhar. Em seguida, ela confirma a intenção, executa a ação principal e recebe uma resposta que pode ser consultada depois. Para nós, acessibilidade cognitiva só está concluído quando o caminho é compreensível, existe recuperação para erro e o servidor registra uma evidência coerente.

O aceite não fica apenas no texto. A evidência prevista é: Teste de compreensão e revisão com checklist de conteúdo claro. O critério de aceite deve ser escrito antes do piloto, com entrada, resultado esperado, responsável, data e ferramenta de verificação. Assim, a documentação separa conformidade técnica, observação de uso e hipótese de impacto.

A validação combina teste automatizado, inspeção manual e uma tarefa representativa. Vamos registrar tempo, erros, conclusão, compreensão e barreiras observadas por teclado, leitor de tela, zoom, contraste, texto, áudio ou conexão instável. O objetivo não é criar uma aparência de certeza, e sim manter um ciclo curto entre observar, corrigir e conferir novamente.

Também registramos uma limitação, porque ela faz parte do resultado: Simplificar não pode apagar informação legal ou financeira importante. O próximo passo é transformar a pendência em issue, teste, decisão de reunião ou evidência de piloto. Este capítulo não promete superioridade nem receita: ele deixa um contrato verificável para que a equipe consiga descobrir se a hipótese se sustenta.

\section{Privacidade LGPD}
Na leitura que fizemos do projeto, perfil, deficiência, localização e histórico podem ser dados pessoais sensíveis ou identificáveis.

Depois de comparar as alternativas, a equipe decidiu o seguinte: Minimização, consentimento, finalidade, exportação/exclusão, retenção e controle institucional.

Na implementação, este requisito é tratado como parte do sistema e não como uma tela isolada. A interface orienta a próxima ação; a API confere identidade, papel, estado e limites; a persistência guarda somente o necessário; e as integrações externas ficam condicionadas às variáveis do ambiente. O escopo continua exclusivamente de software, sem depender de escopo físico.

O fluxo começa quando a pessoa encontra uma necessidade ou uma habilidade que deseja compartilhar. Em seguida, ela confirma a intenção, executa a ação principal e recebe uma resposta que pode ser consultada depois. Para nós, privacidade lgpd só está concluído quando o caminho é compreensível, existe recuperação para erro e o servidor registra uma evidência coerente.

Para conferir se a decisão funciona, vamos procurar esta evidência: Registro de consentimento, política e testes de autorização. O critério de aceite deve ser escrito antes do piloto, com entrada, resultado esperado, responsável, data e ferramenta de verificação. Assim, a documentação separa conformidade técnica, observação de uso e hipótese de impacto.

A validação combina teste automatizado, inspeção manual e uma tarefa representativa. Vamos registrar tempo, erros, conclusão, compreensão e barreiras observadas por teclado, leitor de tela, zoom, contraste, texto, áudio ou conexão instável. O objetivo não é criar uma aparência de certeza, e sim manter um ciclo curto entre observar, corrigir e conferir novamente.

O ponto que ainda merece cuidado é o seguinte: Definições jurídicas devem ser revisadas por responsável da instituição. O próximo passo é transformar a pendência em issue, teste, decisão de reunião ou evidência de piloto. Este capítulo não promete superioridade nem receita: ele deixa um contrato verificável para que a equipe consiga descobrir se a hipótese se sustenta.

\section{Recomendação explicável}
Durante a consolidação do MVP, percebemos que ranking sem explicação pode reproduzir popularidade e desconfiança.

No protótipo, esse entendimento aparece assim: Pontuação inicial pondera habilidades, disponibilidade, objetivo e feedback; resposta inclui fatores e alternativas.

Na implementação, este requisito é tratado como parte do sistema e não como uma tela isolada. A interface orienta a próxima ação; a API confere identidade, papel, estado e limites; a persistência guarda somente o necessário; e as integrações externas ficam condicionadas às variáveis do ambiente. O escopo continua exclusivamente de software, sem depender de escopo físico.

O fluxo começa quando a pessoa encontra uma necessidade ou uma habilidade que deseja compartilhar. Em seguida, ela confirma a intenção, executa a ação principal e recebe uma resposta que pode ser consultada depois. Para nós, recomendação explicável só está concluído quando o caminho é compreensível, existe recuperação para erro e o servidor registra uma evidência coerente.

A forma mais honesta de verificar este requisito é registrar: Fixture com casos controlados e teste de razão exibida. O critério de aceite deve ser escrito antes do piloto, com entrada, resultado esperado, responsável, data e ferramenta de verificação. Assim, a documentação separa conformidade técnica, observação de uso e hipótese de impacto.

A validação combina teste automatizado, inspeção manual e uma tarefa representativa. Vamos registrar tempo, erros, conclusão, compreensão e barreiras observadas por teclado, leitor de tela, zoom, contraste, texto, áudio ou conexão instável. O objetivo não é criar uma aparência de certeza, e sim manter um ciclo curto entre observar, corrigir e conferir novamente.

A equipe não trata este item como resolvido antes de enfrentar o risco: Pontuação heurística não é prova de compatibilidade; permitir ajuste e denúncia. O próximo passo é transformar a pendência em issue, teste, decisão de reunião ou evidência de piloto. Este capítulo não promete superioridade nem receita: ele deixa um contrato verificável para que a equipe consiga descobrir se a hipótese se sustenta.

\section{Métricas de produto}
A discussão de métricas de produto começou por uma situação concreta: visualizações e curtidas não medem se a troca funcionou.

Nossa escolha para este ponto foi direta: Medir ativação, tempo até primeira proposta, aceite, conclusão, avaliação, retorno e acessibilidade percebida.

Na implementação, este requisito é tratado como parte do sistema e não como uma tela isolada. A interface orienta a próxima ação; a API confere identidade, papel, estado e limites; a persistência guarda somente o necessário; e as integrações externas ficam condicionadas às variáveis do ambiente. O escopo continua exclusivamente de software, sem depender de escopo físico.

O fluxo começa quando a pessoa encontra uma necessidade ou uma habilidade que deseja compartilhar. Em seguida, ela confirma a intenção, executa a ação principal e recebe uma resposta que pode ser consultada depois. Para nós, métricas de produto só está concluído quando o caminho é compreensível, existe recuperação para erro e o servidor registra uma evidência coerente.

O aceite não fica apenas no texto. A evidência prevista é: Eventos de produto sem conteúdo sensível e denominadores documentados. O critério de aceite deve ser escrito antes do piloto, com entrada, resultado esperado, responsável, data e ferramenta de verificação. Assim, a documentação separa conformidade técnica, observação de uso e hipótese de impacto.

A validação combina teste automatizado, inspeção manual e uma tarefa representativa. Vamos registrar tempo, erros, conclusão, compreensão e barreiras observadas por teclado, leitor de tela, zoom, contraste, texto, áudio ou conexão instável. O objetivo não é criar uma aparência de certeza, e sim manter um ciclo curto entre observar, corrigir e conferir novamente.

Também registramos uma limitação, porque ela faz parte do resultado: Não coletar mais dados do que o necessário para a pergunta de pesquisa. O próximo passo é transformar a pendência em issue, teste, decisão de reunião ou evidência de piloto. Este capítulo não promete superioridade nem receita: ele deixa um contrato verificável para que a equipe consiga descobrir se a hipótese se sustenta.

\section{Observabilidade}
Na leitura que fizemos do projeto, sem health, métricas e logs sanitizados a equipe não sabe se o deploy está saudável.

Depois de comparar as alternativas, a equipe decidiu o seguinte: Actuator health/info/prometheus, logs estruturados, métricas de API, banco, fila, webhook e WebSocket.

Na implementação, este requisito é tratado como parte do sistema e não como uma tela isolada. A interface orienta a próxima ação; a API confere identidade, papel, estado e limites; a persistência guarda somente o necessário; e as integrações externas ficam condicionadas às variáveis do ambiente. O escopo continua exclusivamente de software, sem depender de escopo físico.

O fluxo começa quando a pessoa encontra uma necessidade ou uma habilidade que deseja compartilhar. Em seguida, ela confirma a intenção, executa a ação principal e recebe uma resposta que pode ser consultada depois. Para nós, observabilidade só está concluído quando o caminho é compreensível, existe recuperação para erro e o servidor registra uma evidência coerente.

Para conferir se a decisão funciona, vamos procurar esta evidência: Dashboard local Prometheus/Grafana e alertas de disponibilidade. O critério de aceite deve ser escrito antes do piloto, com entrada, resultado esperado, responsável, data e ferramenta de verificação. Assim, a documentação separa conformidade técnica, observação de uso e hipótese de impacto.

A validação combina teste automatizado, inspeção manual e uma tarefa representativa. Vamos registrar tempo, erros, conclusão, compreensão e barreiras observadas por teclado, leitor de tela, zoom, contraste, texto, áudio ou conexão instável. O objetivo não é criar uma aparência de certeza, e sim manter um ciclo curto entre observar, corrigir e conferir novamente.

O ponto que ainda merece cuidado é o seguinte: Prometheus/Grafana podem ficar locais no TCC para controlar custo; produção precisa de alternativa. O próximo passo é transformar a pendência em issue, teste, decisão de reunião ou evidência de piloto. Este capítulo não promete superioridade nem receita: ele deixa um contrato verificável para que a equipe consiga descobrir se a hipótese se sustenta.

\section{Railway e custo}
Na leitura que fizemos do projeto, uma arquitetura completa pode exceder créditos gratuitos se todos os serviços forem ativados.

Depois de comparar as alternativas, a equipe decidiu o seguinte: MVP usa frontend, backend e PostgreSQL; Caffeine fica padrão; Redis, RabbitMQ e observabilidade pesada são opcionais.

Na implementação, este requisito é tratado como parte do sistema e não como uma tela isolada. A interface orienta a próxima ação; a API confere identidade, papel, estado e limites; a persistência guarda somente o necessário; e as integrações externas ficam condicionadas às variáveis do ambiente. O escopo continua exclusivamente de software, sem depender de escopo físico.

O fluxo começa quando a pessoa encontra uma necessidade ou uma habilidade que deseja compartilhar. Em seguida, ela confirma a intenção, executa a ação principal e recebe uma resposta que pode ser consultada depois. Para nós, railway e custo só está concluído quando o caminho é compreensível, existe recuperação para erro e o servidor registra uma evidência coerente.

Para conferir se a decisão funciona, vamos procurar esta evidência: Variáveis de ambiente, health check e estimativa mensal. O critério de aceite deve ser escrito antes do piloto, com entrada, resultado esperado, responsável, data e ferramenta de verificação. Assim, a documentação separa conformidade técnica, observação de uso e hipótese de impacto.

A validação combina teste automatizado, inspeção manual e uma tarefa representativa. Vamos registrar tempo, erros, conclusão, compreensão e barreiras observadas por teclado, leitor de tela, zoom, contraste, texto, áudio ou conexão instável. O objetivo não é criar uma aparência de certeza, e sim manter um ciclo curto entre observar, corrigir e conferir novamente.

O ponto que ainda merece cuidado é o seguinte: Subdomínio e plano podem mudar; revisar preço atual antes do contrato. O próximo passo é transformar a pendência em issue, teste, decisão de reunião ou evidência de piloto. Este capítulo não promete superioridade nem receita: ele deixa um contrato verificável para que a equipe consiga descobrir se a hipótese se sustenta.

\section{CI/CD}
A discussão de ci/cd começou por uma situação concreta: build local não é evidência de entrega reproduzível.

Nossa escolha para este ponto foi direta: Pipeline executa lint, testes frontend, build Vite, testes Maven, validação Flyway, scan de segredos e imagem Docker.

Na implementação, este requisito é tratado como parte do sistema e não como uma tela isolada. A interface orienta a próxima ação; a API confere identidade, papel, estado e limites; a persistência guarda somente o necessário; e as integrações externas ficam condicionadas às variáveis do ambiente. O escopo continua exclusivamente de software, sem depender de escopo físico.

O fluxo começa quando a pessoa encontra uma necessidade ou uma habilidade que deseja compartilhar. Em seguida, ela confirma a intenção, executa a ação principal e recebe uma resposta que pode ser consultada depois. Para nós, ci/cd só está concluído quando o caminho é compreensível, existe recuperação para erro e o servidor registra uma evidência coerente.

O aceite não fica apenas no texto. A evidência prevista é: Workflow GitHub Actions e artefatos de build. O critério de aceite deve ser escrito antes do piloto, com entrada, resultado esperado, responsável, data e ferramenta de verificação. Assim, a documentação separa conformidade técnica, observação de uso e hipótese de impacto.

A validação combina teste automatizado, inspeção manual e uma tarefa representativa. Vamos registrar tempo, erros, conclusão, compreensão e barreiras observadas por teclado, leitor de tela, zoom, contraste, texto, áudio ou conexão instável. O objetivo não é criar uma aparência de certeza, e sim manter um ciclo curto entre observar, corrigir e conferir novamente.

Também registramos uma limitação, porque ela faz parte do resultado: Runner pode não ter navegador ou serviço externo; separar testes unitários de integração. O próximo passo é transformar a pendência em issue, teste, decisão de reunião ou evidência de piloto. Este capítulo não promete superioridade nem receita: ele deixa um contrato verificável para que a equipe consiga descobrir se a hipótese se sustenta.

\section{Backup e rollback}
Na leitura que fizemos do projeto, banco e integrações externas exigem recuperação testada, não apenas esperança.

Depois de comparar as alternativas, a equipe decidiu o seguinte: Backup antes de migration, restore em ambiente isolado, versionamento de imagem e rollback documentado.

Na implementação, este requisito é tratado como parte do sistema e não como uma tela isolada. A interface orienta a próxima ação; a API confere identidade, papel, estado e limites; a persistência guarda somente o necessário; e as integrações externas ficam condicionadas às variáveis do ambiente. O escopo continua exclusivamente de software, sem depender de escopo físico.

O fluxo começa quando a pessoa encontra uma necessidade ou uma habilidade que deseja compartilhar. Em seguida, ela confirma a intenção, executa a ação principal e recebe uma resposta que pode ser consultada depois. Para nós, backup e rollback só está concluído quando o caminho é compreensível, existe recuperação para erro e o servidor registra uma evidência coerente.

Para conferir se a decisão funciona, vamos procurar esta evidência: Ensaio com checklist e tempos de RPO/RTO. O critério de aceite deve ser escrito antes do piloto, com entrada, resultado esperado, responsável, data e ferramenta de verificação. Assim, a documentação separa conformidade técnica, observação de uso e hipótese de impacto.

A validação combina teste automatizado, inspeção manual e uma tarefa representativa. Vamos registrar tempo, erros, conclusão, compreensão e barreiras observadas por teclado, leitor de tela, zoom, contraste, texto, áudio ou conexão instável. O objetivo não é criar uma aparência de certeza, e sim manter um ciclo curto entre observar, corrigir e conferir novamente.

O ponto que ainda merece cuidado é o seguinte: Rollback de schema pode ser destrutivo; preferir migrations compatíveis para frente. O próximo passo é transformar a pendência em issue, teste, decisão de reunião ou evidência de piloto. Este capítulo não promete superioridade nem receita: ele deixa um contrato verificável para que a equipe consiga descobrir se a hipótese se sustenta.

\section{Testes de autorização}
Durante a consolidação do MVP, percebemos que uma tela escondida não protege endpoint.

No protótipo, esse entendimento aparece assim: Cada recurso verifica papel, proprietário, vínculo e estado; IDs enviados pelo cliente nunca substituem o principal.

Na implementação, este requisito é tratado como parte do sistema e não como uma tela isolada. A interface orienta a próxima ação; a API confere identidade, papel, estado e limites; a persistência guarda somente o necessário; e as integrações externas ficam condicionadas às variáveis do ambiente. O escopo continua exclusivamente de software, sem depender de escopo físico.

O fluxo começa quando a pessoa encontra uma necessidade ou uma habilidade que deseja compartilhar. Em seguida, ela confirma a intenção, executa a ação principal e recebe uma resposta que pode ser consultada depois. Para nós, testes de autorização só está concluído quando o caminho é compreensível, existe recuperação para erro e o servidor registra uma evidência coerente.

A forma mais honesta de verificar este requisito é registrar: Matriz de 200, 401, 403 e 404 para rotas críticas. O critério de aceite deve ser escrito antes do piloto, com entrada, resultado esperado, responsável, data e ferramenta de verificação. Assim, a documentação separa conformidade técnica, observação de uso e hipótese de impacto.

A validação combina teste automatizado, inspeção manual e uma tarefa representativa. Vamos registrar tempo, erros, conclusão, compreensão e barreiras observadas por teclado, leitor de tela, zoom, contraste, texto, áudio ou conexão instável. O objetivo não é criar uma aparência de certeza, e sim manter um ciclo curto entre observar, corrigir e conferir novamente.

A equipe não trata este item como resolvido antes de enfrentar o risco: Cobertura deve incluir WebSocket, webhook, administração e dados de outros usuários. O próximo passo é transformar a pendência em issue, teste, decisão de reunião ou evidência de piloto. Este capítulo não promete superioridade nem receita: ele deixa um contrato verificável para que a equipe consiga descobrir se a hipótese se sustenta.

\section{Teste de carga básico}
A discussão de teste de carga básico começou por uma situação concreta: latência percebida e concorrência podem quebrar saldo, feed ou chat.

Nossa escolha para este ponto foi direta: Testar leitura de feed, criação de proposta, confirmação de aula e geração de certificado em baixa escala.

Na implementação, este requisito é tratado como parte do sistema e não como uma tela isolada. A interface orienta a próxima ação; a API confere identidade, papel, estado e limites; a persistência guarda somente o necessário; e as integrações externas ficam condicionadas às variáveis do ambiente. O escopo continua exclusivamente de software, sem depender de escopo físico.

O fluxo começa quando a pessoa encontra uma necessidade ou uma habilidade que deseja compartilhar. Em seguida, ela confirma a intenção, executa a ação principal e recebe uma resposta que pode ser consultada depois. Para nós, teste de carga básico só está concluído quando o caminho é compreensível, existe recuperação para erro e o servidor registra uma evidência coerente.

O aceite não fica apenas no texto. A evidência prevista é: p95, taxa de erro e inconsistência de saldo. O critério de aceite deve ser escrito antes do piloto, com entrada, resultado esperado, responsável, data e ferramenta de verificação. Assim, a documentação separa conformidade técnica, observação de uso e hipótese de impacto.

A validação combina teste automatizado, inspeção manual e uma tarefa representativa. Vamos registrar tempo, erros, conclusão, compreensão e barreiras observadas por teclado, leitor de tela, zoom, contraste, texto, áudio ou conexão instável. O objetivo não é criar uma aparência de certeza, e sim manter um ciclo curto entre observar, corrigir e conferir novamente.

Também registramos uma limitação, porque ela faz parte do resultado: Não executar carga em serviço pago sem autorização e janela controlada. O próximo passo é transformar a pendência em issue, teste, decisão de reunião ou evidência de piloto. Este capítulo não promete superioridade nem receita: ele deixa um contrato verificável para que a equipe consiga descobrir se a hipótese se sustenta.

\section{Plano de validação}
Na leitura que fizemos do projeto, o TCC precisa comprovar o que foi implementado e separar resultado de proposta.

Depois de comparar as alternativas, a equipe decidiu o seguinte: Relatar método, participantes, tarefas, métricas, limitações, consentimento e resultados reais somente depois de coletados.

Na implementação, este requisito é tratado como parte do sistema e não como uma tela isolada. A interface orienta a próxima ação; a API confere identidade, papel, estado e limites; a persistência guarda somente o necessário; e as integrações externas ficam condicionadas às variáveis do ambiente. O escopo continua exclusivamente de software, sem depender de escopo físico.

O fluxo começa quando a pessoa encontra uma necessidade ou uma habilidade que deseja compartilhar. Em seguida, ela confirma a intenção, executa a ação principal e recebe uma resposta que pode ser consultada depois. Para nós, plano de validação só está concluído quando o caminho é compreensível, existe recuperação para erro e o servidor registra uma evidência coerente.

Para conferir se a decisão funciona, vamos procurar esta evidência: Protocolos, planilhas e evidências versionadas. O critério de aceite deve ser escrito antes do piloto, com entrada, resultado esperado, responsável, data e ferramenta de verificação. Assim, a documentação separa conformidade técnica, observação de uso e hipótese de impacto.

A validação combina teste automatizado, inspeção manual e uma tarefa representativa. Vamos registrar tempo, erros, conclusão, compreensão e barreiras observadas por teclado, leitor de tela, zoom, contraste, texto, áudio ou conexão instável. O objetivo não é criar uma aparência de certeza, e sim manter um ciclo curto entre observar, corrigir e conferir novamente.

O ponto que ainda merece cuidado é o seguinte: Não inventar amostra, ganho, acurácia ou satisfação. O próximo passo é transformar a pendência em issue, teste, decisão de reunião ou evidência de piloto. Este capítulo não promete superioridade nem receita: ele deixa um contrato verificável para que a equipe consiga descobrir se a hipótese se sustenta.

\section{Roadmap de startup}
Durante a consolidação do MVP, percebemos que a visão de startup precisa ser executável por etapas.

No protótipo, esse entendimento aparece assim: Fase 1 valida troca; fase 2 piloto institucional; fase 3 integrações; fase 4 escala e modelos de recomendação.

Na implementação, este requisito é tratado como parte do sistema e não como uma tela isolada. A interface orienta a próxima ação; a API confere identidade, papel, estado e limites; a persistência guarda somente o necessário; e as integrações externas ficam condicionadas às variáveis do ambiente. O escopo continua exclusivamente de software, sem depender de escopo físico.

O fluxo começa quando a pessoa encontra uma necessidade ou uma habilidade que deseja compartilhar. Em seguida, ela confirma a intenção, executa a ação principal e recebe uma resposta que pode ser consultada depois. Para nós, roadmap de startup só está concluído quando o caminho é compreensível, existe recuperação para erro e o servidor registra uma evidência coerente.

A forma mais honesta de verificar este requisito é registrar: OKRs, responsáveis, custo e hipótese por fase. O critério de aceite deve ser escrito antes do piloto, com entrada, resultado esperado, responsável, data e ferramenta de verificação. Assim, a documentação separa conformidade técnica, observação de uso e hipótese de impacto.

A validação combina teste automatizado, inspeção manual e uma tarefa representativa. Vamos registrar tempo, erros, conclusão, compreensão e barreiras observadas por teclado, leitor de tela, zoom, contraste, texto, áudio ou conexão instável. O objetivo não é criar uma aparência de certeza, e sim manter um ciclo curto entre observar, corrigir e conferir novamente.

A equipe não trata este item como resolvido antes de enfrentar o risco: Diversificar demais antes de validar o fluxo principal dilui caixa e foco. O próximo passo é transformar a pendência em issue, teste, decisão de reunião ou evidência de piloto. Este capítulo não promete superioridade nem receita: ele deixa um contrato verificável para que a equipe consiga descobrir se a hipótese se sustenta.

\section{Critério de sucesso}
Durante a consolidação do MVP, percebemos que uma entrega pode ter muitas telas e ainda não resolver o problema.

No protótipo, esse entendimento aparece assim: O MVP será considerado pronto quando uma pessoa completar uma troca, outra receber crédito/registro, a equipe auditar o fluxo e usuários com diferentes modos concluírem tarefas principais.

Na implementação, este requisito é tratado como parte do sistema e não como uma tela isolada. A interface orienta a próxima ação; a API confere identidade, papel, estado e limites; a persistência guarda somente o necessário; e as integrações externas ficam condicionadas às variáveis do ambiente. O escopo continua exclusivamente de software, sem depender de escopo físico.

O fluxo começa quando a pessoa encontra uma necessidade ou uma habilidade que deseja compartilhar. Em seguida, ela confirma a intenção, executa a ação principal e recebe uma resposta que pode ser consultada depois. Para nós, critério de sucesso só está concluído quando o caminho é compreensível, existe recuperação para erro e o servidor registra uma evidência coerente.

A forma mais honesta de verificar este requisito é registrar: Demo gravada, testes e relatório de lacunas. O critério de aceite deve ser escrito antes do piloto, com entrada, resultado esperado, responsável, data e ferramenta de verificação. Assim, a documentação separa conformidade técnica, observação de uso e hipótese de impacto.

A validação combina teste automatizado, inspeção manual e uma tarefa representativa. Vamos registrar tempo, erros, conclusão, compreensão e barreiras observadas por teclado, leitor de tela, zoom, contraste, texto, áudio ou conexão instável. O objetivo não é criar uma aparência de certeza, e sim manter um ciclo curto entre observar, corrigir e conferir novamente.

A equipe não trata este item como resolvido antes de enfrentar o risco: Pronto para TCC não significa pronto para produção comercial sem piloto e suporte. O próximo passo é transformar a pendência em issue, teste, decisão de reunião ou evidência de piloto. Este capítulo não promete superioridade nem receita: ele deixa um contrato verificável para que a equipe consiga descobrir se a hipótese se sustenta.

\section{Funcionalidade 01: Cadastro por e-mail}
Na leitura que fizemos do projeto, esta página especifica o comportamento de cadastro por e-mail no produto.

Depois de comparar as alternativas, a equipe decidiu o seguinte: Implementação esperada: Firebase cria a identidade; backend provisiona o perfil; nenhuma senha é copiada para SkillSwap.

Na implementação, este requisito é tratado como parte do sistema e não como uma tela isolada. A interface orienta a próxima ação; a API confere identidade, papel, estado e limites; a persistência guarda somente o necessário; e as integrações externas ficam condicionadas às variáveis do ambiente. O escopo continua exclusivamente de software, sem depender de escopo físico.

O fluxo começa quando a pessoa encontra uma necessidade ou uma habilidade que deseja compartilhar. Em seguida, ela confirma a intenção, executa a ação principal e recebe uma resposta que pode ser consultada depois. Para nós, funcionalidade 01: cadastro por e-mail só está concluído quando o caminho é compreensível, existe recuperação para erro e o servidor registra uma evidência coerente.

Para conferir se a decisão funciona, vamos procurar esta evidência: Critério de aceite: a funcionalidade deve funcionar com usuário autenticado quando houver dados pessoais, exibir feedback compreensível e preservar autorização no servidor. O critério de aceite deve ser escrito antes do piloto, com entrada, resultado esperado, responsável, data e ferramenta de verificação. Assim, a documentação separa conformidade técnica, observação de uso e hipótese de impacto.

A validação combina teste automatizado, inspeção manual e uma tarefa representativa. Vamos registrar tempo, erros, conclusão, compreensão e barreiras observadas por teclado, leitor de tela, zoom, contraste, texto, áudio ou conexão instável. O objetivo não é criar uma aparência de certeza, e sim manter um ciclo curto entre observar, corrigir e conferir novamente.

O ponto que ainda merece cuidado é o seguinte: Evidência: rota, serviço, componente, teste ou checklist correspondente; resultados de usuários só entram depois da execução do protocolo. O próximo passo é transformar a pendência em issue, teste, decisão de reunião ou evidência de piloto. Este capítulo não promete superioridade nem receita: ele deixa um contrato verificável para que a equipe consiga descobrir se a hipótese se sustenta.

\section{Funcionalidade 02: Login Google}
Durante a consolidação do MVP, percebemos que esta página especifica o comportamento de login google no produto.

No protótipo, esse entendimento aparece assim: Implementação esperada: Provedor social reduz fricção e continua sob as regras de verificação do Firebase.

Na implementação, este requisito é tratado como parte do sistema e não como uma tela isolada. A interface orienta a próxima ação; a API confere identidade, papel, estado e limites; a persistência guarda somente o necessário; e as integrações externas ficam condicionadas às variáveis do ambiente. O escopo continua exclusivamente de software, sem depender de escopo físico.

O fluxo começa quando a pessoa encontra uma necessidade ou uma habilidade que deseja compartilhar. Em seguida, ela confirma a intenção, executa a ação principal e recebe uma resposta que pode ser consultada depois. Para nós, funcionalidade 02: login google só está concluído quando o caminho é compreensível, existe recuperação para erro e o servidor registra uma evidência coerente.

A forma mais honesta de verificar este requisito é registrar: Critério de aceite: a funcionalidade deve funcionar com usuário autenticado quando houver dados pessoais, exibir feedback compreensível e preservar autorização no servidor. O critério de aceite deve ser escrito antes do piloto, com entrada, resultado esperado, responsável, data e ferramenta de verificação. Assim, a documentação separa conformidade técnica, observação de uso e hipótese de impacto.

A validação combina teste automatizado, inspeção manual e uma tarefa representativa. Vamos registrar tempo, erros, conclusão, compreensão e barreiras observadas por teclado, leitor de tela, zoom, contraste, texto, áudio ou conexão instável. O objetivo não é criar uma aparência de certeza, e sim manter um ciclo curto entre observar, corrigir e conferir novamente.

A equipe não trata este item como resolvido antes de enfrentar o risco: Evidência: rota, serviço, componente, teste ou checklist correspondente; resultados de usuários só entram depois da execução do protocolo. O próximo passo é transformar a pendência em issue, teste, decisão de reunião ou evidência de piloto. Este capítulo não promete superioridade nem receita: ele deixa um contrato verificável para que a equipe consiga descobrir se a hipótese se sustenta.

\section{Funcionalidade 03: Perfil público}
Durante a consolidação do MVP, percebemos que esta página especifica o comportamento de perfil público no produto.

No protótipo, esse entendimento aparece assim: Implementação esperada: Mostra nome, username, bio, habilidades e evidências sem expor dados privados.

Na implementação, este requisito é tratado como parte do sistema e não como uma tela isolada. A interface orienta a próxima ação; a API confere identidade, papel, estado e limites; a persistência guarda somente o necessário; e as integrações externas ficam condicionadas às variáveis do ambiente. O escopo continua exclusivamente de software, sem depender de escopo físico.

O fluxo começa quando a pessoa encontra uma necessidade ou uma habilidade que deseja compartilhar. Em seguida, ela confirma a intenção, executa a ação principal e recebe uma resposta que pode ser consultada depois. Para nós, funcionalidade 03: perfil público só está concluído quando o caminho é compreensível, existe recuperação para erro e o servidor registra uma evidência coerente.

A forma mais honesta de verificar este requisito é registrar: Critério de aceite: a funcionalidade deve funcionar com usuário autenticado quando houver dados pessoais, exibir feedback compreensível e preservar autorização no servidor. O critério de aceite deve ser escrito antes do piloto, com entrada, resultado esperado, responsável, data e ferramenta de verificação. Assim, a documentação separa conformidade técnica, observação de uso e hipótese de impacto.

A validação combina teste automatizado, inspeção manual e uma tarefa representativa. Vamos registrar tempo, erros, conclusão, compreensão e barreiras observadas por teclado, leitor de tela, zoom, contraste, texto, áudio ou conexão instável. O objetivo não é criar uma aparência de certeza, e sim manter um ciclo curto entre observar, corrigir e conferir novamente.

A equipe não trata este item como resolvido antes de enfrentar o risco: Evidência: rota, serviço, componente, teste ou checklist correspondente; resultados de usuários só entram depois da execução do protocolo. O próximo passo é transformar a pendência em issue, teste, decisão de reunião ou evidência de piloto. Este capítulo não promete superioridade nem receita: ele deixa um contrato verificável para que a equipe consiga descobrir se a hipótese se sustenta.

\section{Funcionalidade 04: Preferências}
Na leitura que fizemos do projeto, esta página especifica o comportamento de preferências no produto.

Depois de comparar as alternativas, a equipe decidiu o seguinte: Implementação esperada: Tema, fonte e contraste ficam locais; dados de acessibilidade persistentes exigem consentimento.

Na implementação, este requisito é tratado como parte do sistema e não como uma tela isolada. A interface orienta a próxima ação; a API confere identidade, papel, estado e limites; a persistência guarda somente o necessário; e as integrações externas ficam condicionadas às variáveis do ambiente. O escopo continua exclusivamente de software, sem depender de escopo físico.

O fluxo começa quando a pessoa encontra uma necessidade ou uma habilidade que deseja compartilhar. Em seguida, ela confirma a intenção, executa a ação principal e recebe uma resposta que pode ser consultada depois. Para nós, funcionalidade 04: preferências só está concluído quando o caminho é compreensível, existe recuperação para erro e o servidor registra uma evidência coerente.

Para conferir se a decisão funciona, vamos procurar esta evidência: Critério de aceite: a funcionalidade deve funcionar com usuário autenticado quando houver dados pessoais, exibir feedback compreensível e preservar autorização no servidor. O critério de aceite deve ser escrito antes do piloto, com entrada, resultado esperado, responsável, data e ferramenta de verificação. Assim, a documentação separa conformidade técnica, observação de uso e hipótese de impacto.

A validação combina teste automatizado, inspeção manual e uma tarefa representativa. Vamos registrar tempo, erros, conclusão, compreensão e barreiras observadas por teclado, leitor de tela, zoom, contraste, texto, áudio ou conexão instável. O objetivo não é criar uma aparência de certeza, e sim manter um ciclo curto entre observar, corrigir e conferir novamente.

O ponto que ainda merece cuidado é o seguinte: Evidência: rota, serviço, componente, teste ou checklist correspondente; resultados de usuários só entram depois da execução do protocolo. O próximo passo é transformar a pendência em issue, teste, decisão de reunião ou evidência de piloto. Este capítulo não promete superioridade nem receita: ele deixa um contrato verificável para que a equipe consiga descobrir se a hipótese se sustenta.

\section{Funcionalidade 05: Busca}
Durante a consolidação do MVP, percebemos que esta página especifica o comportamento de busca no produto.

No protótipo, esse entendimento aparece assim: Implementação esperada: Consulta pessoas e habilidades com paginação, normalização e limites de entrada.

Na implementação, este requisito é tratado como parte do sistema e não como uma tela isolada. A interface orienta a próxima ação; a API confere identidade, papel, estado e limites; a persistência guarda somente o necessário; e as integrações externas ficam condicionadas às variáveis do ambiente. O escopo continua exclusivamente de software, sem depender de escopo físico.

O fluxo começa quando a pessoa encontra uma necessidade ou uma habilidade que deseja compartilhar. Em seguida, ela confirma a intenção, executa a ação principal e recebe uma resposta que pode ser consultada depois. Para nós, funcionalidade 05: busca só está concluído quando o caminho é compreensível, existe recuperação para erro e o servidor registra uma evidência coerente.

A forma mais honesta de verificar este requisito é registrar: Critério de aceite: a funcionalidade deve funcionar com usuário autenticado quando houver dados pessoais, exibir feedback compreensível e preservar autorização no servidor. O critério de aceite deve ser escrito antes do piloto, com entrada, resultado esperado, responsável, data e ferramenta de verificação. Assim, a documentação separa conformidade técnica, observação de uso e hipótese de impacto.

A validação combina teste automatizado, inspeção manual e uma tarefa representativa. Vamos registrar tempo, erros, conclusão, compreensão e barreiras observadas por teclado, leitor de tela, zoom, contraste, texto, áudio ou conexão instável. O objetivo não é criar uma aparência de certeza, e sim manter um ciclo curto entre observar, corrigir e conferir novamente.

A equipe não trata este item como resolvido antes de enfrentar o risco: Evidência: rota, serviço, componente, teste ou checklist correspondente; resultados de usuários só entram depois da execução do protocolo. O próximo passo é transformar a pendência em issue, teste, decisão de reunião ou evidência de piloto. Este capítulo não promete superioridade nem receita: ele deixa um contrato verificável para que a equipe consiga descobrir se a hipótese se sustenta.

\section{Funcionalidade 06: Feed}
A discussão de funcionalidade 06: feed começou por uma situação concreta: esta página especifica o comportamento de feed no produto.

Nossa escolha para este ponto foi direta: Implementação esperada: Publicações, comentários, curtidas e comunidades ampliam descoberta sem substituir o fluxo de troca.

Na implementação, este requisito é tratado como parte do sistema e não como uma tela isolada. A interface orienta a próxima ação; a API confere identidade, papel, estado e limites; a persistência guarda somente o necessário; e as integrações externas ficam condicionadas às variáveis do ambiente. O escopo continua exclusivamente de software, sem depender de escopo físico.

O fluxo começa quando a pessoa encontra uma necessidade ou uma habilidade que deseja compartilhar. Em seguida, ela confirma a intenção, executa a ação principal e recebe uma resposta que pode ser consultada depois. Para nós, funcionalidade 06: feed só está concluído quando o caminho é compreensível, existe recuperação para erro e o servidor registra uma evidência coerente.

O aceite não fica apenas no texto. A evidência prevista é: Critério de aceite: a funcionalidade deve funcionar com usuário autenticado quando houver dados pessoais, exibir feedback compreensível e preservar autorização no servidor. O critério de aceite deve ser escrito antes do piloto, com entrada, resultado esperado, responsável, data e ferramenta de verificação. Assim, a documentação separa conformidade técnica, observação de uso e hipótese de impacto.

A validação combina teste automatizado, inspeção manual e uma tarefa representativa. Vamos registrar tempo, erros, conclusão, compreensão e barreiras observadas por teclado, leitor de tela, zoom, contraste, texto, áudio ou conexão instável. O objetivo não é criar uma aparência de certeza, e sim manter um ciclo curto entre observar, corrigir e conferir novamente.

Também registramos uma limitação, porque ela faz parte do resultado: Evidência: rota, serviço, componente, teste ou checklist correspondente; resultados de usuários só entram depois da execução do protocolo. O próximo passo é transformar a pendência em issue, teste, decisão de reunião ou evidência de piloto. Este capítulo não promete superioridade nem receita: ele deixa um contrato verificável para que a equipe consiga descobrir se a hipótese se sustenta.

\section{Funcionalidade 07: Comunidades}
Na leitura que fizemos do projeto, esta página especifica o comportamento de comunidades no produto.

Depois de comparar as alternativas, a equipe decidiu o seguinte: Implementação esperada: Agrupam interesses e permitem contexto social para a recomendação.

Na implementação, este requisito é tratado como parte do sistema e não como uma tela isolada. A interface orienta a próxima ação; a API confere identidade, papel, estado e limites; a persistência guarda somente o necessário; e as integrações externas ficam condicionadas às variáveis do ambiente. O escopo continua exclusivamente de software, sem depender de escopo físico.

O fluxo começa quando a pessoa encontra uma necessidade ou uma habilidade que deseja compartilhar. Em seguida, ela confirma a intenção, executa a ação principal e recebe uma resposta que pode ser consultada depois. Para nós, funcionalidade 07: comunidades só está concluído quando o caminho é compreensível, existe recuperação para erro e o servidor registra uma evidência coerente.

Para conferir se a decisão funciona, vamos procurar esta evidência: Critério de aceite: a funcionalidade deve funcionar com usuário autenticado quando houver dados pessoais, exibir feedback compreensível e preservar autorização no servidor. O critério de aceite deve ser escrito antes do piloto, com entrada, resultado esperado, responsável, data e ferramenta de verificação. Assim, a documentação separa conformidade técnica, observação de uso e hipótese de impacto.

A validação combina teste automatizado, inspeção manual e uma tarefa representativa. Vamos registrar tempo, erros, conclusão, compreensão e barreiras observadas por teclado, leitor de tela, zoom, contraste, texto, áudio ou conexão instável. O objetivo não é criar uma aparência de certeza, e sim manter um ciclo curto entre observar, corrigir e conferir novamente.

O ponto que ainda merece cuidado é o seguinte: Evidência: rota, serviço, componente, teste ou checklist correspondente; resultados de usuários só entram depois da execução do protocolo. O próximo passo é transformar a pendência em issue, teste, decisão de reunião ou evidência de piloto. Este capítulo não promete superioridade nem receita: ele deixa um contrato verificável para que a equipe consiga descobrir se a hipótese se sustenta.

\section{Funcionalidade 08: Proposta}
A discussão de funcionalidade 08: proposta começou por uma situação concreta: esta página especifica o comportamento de proposta no produto.

Nossa escolha para este ponto foi direta: Implementação esperada: Registra remetente, destinatário, habilidade, mensagem, estado e timestamps.

Na implementação, este requisito é tratado como parte do sistema e não como uma tela isolada. A interface orienta a próxima ação; a API confere identidade, papel, estado e limites; a persistência guarda somente o necessário; e as integrações externas ficam condicionadas às variáveis do ambiente. O escopo continua exclusivamente de software, sem depender de escopo físico.

O fluxo começa quando a pessoa encontra uma necessidade ou uma habilidade que deseja compartilhar. Em seguida, ela confirma a intenção, executa a ação principal e recebe uma resposta que pode ser consultada depois. Para nós, funcionalidade 08: proposta só está concluído quando o caminho é compreensível, existe recuperação para erro e o servidor registra uma evidência coerente.

O aceite não fica apenas no texto. A evidência prevista é: Critério de aceite: a funcionalidade deve funcionar com usuário autenticado quando houver dados pessoais, exibir feedback compreensível e preservar autorização no servidor. O critério de aceite deve ser escrito antes do piloto, com entrada, resultado esperado, responsável, data e ferramenta de verificação. Assim, a documentação separa conformidade técnica, observação de uso e hipótese de impacto.

A validação combina teste automatizado, inspeção manual e uma tarefa representativa. Vamos registrar tempo, erros, conclusão, compreensão e barreiras observadas por teclado, leitor de tela, zoom, contraste, texto, áudio ou conexão instável. O objetivo não é criar uma aparência de certeza, e sim manter um ciclo curto entre observar, corrigir e conferir novamente.

Também registramos uma limitação, porque ela faz parte do resultado: Evidência: rota, serviço, componente, teste ou checklist correspondente; resultados de usuários só entram depois da execução do protocolo. O próximo passo é transformar a pendência em issue, teste, decisão de reunião ou evidência de piloto. Este capítulo não promete superioridade nem receita: ele deixa um contrato verificável para que a equipe consiga descobrir se a hipótese se sustenta.

\section{Funcionalidade 09: Negociação}
A discussão de funcionalidade 09: negociação começou por uma situação concreta: esta página especifica o comportamento de negociação no produto.

Nossa escolha para este ponto foi direta: Implementação esperada: Permite aceitar, recusar, bloquear e concluir, com regras de transição explícitas.

Na implementação, este requisito é tratado como parte do sistema e não como uma tela isolada. A interface orienta a próxima ação; a API confere identidade, papel, estado e limites; a persistência guarda somente o necessário; e as integrações externas ficam condicionadas às variáveis do ambiente. O escopo continua exclusivamente de software, sem depender de escopo físico.

O fluxo começa quando a pessoa encontra uma necessidade ou uma habilidade que deseja compartilhar. Em seguida, ela confirma a intenção, executa a ação principal e recebe uma resposta que pode ser consultada depois. Para nós, funcionalidade 09: negociação só está concluído quando o caminho é compreensível, existe recuperação para erro e o servidor registra uma evidência coerente.

O aceite não fica apenas no texto. A evidência prevista é: Critério de aceite: a funcionalidade deve funcionar com usuário autenticado quando houver dados pessoais, exibir feedback compreensível e preservar autorização no servidor. O critério de aceite deve ser escrito antes do piloto, com entrada, resultado esperado, responsável, data e ferramenta de verificação. Assim, a documentação separa conformidade técnica, observação de uso e hipótese de impacto.

A validação combina teste automatizado, inspeção manual e uma tarefa representativa. Vamos registrar tempo, erros, conclusão, compreensão e barreiras observadas por teclado, leitor de tela, zoom, contraste, texto, áudio ou conexão instável. O objetivo não é criar uma aparência de certeza, e sim manter um ciclo curto entre observar, corrigir e conferir novamente.

Também registramos uma limitação, porque ela faz parte do resultado: Evidência: rota, serviço, componente, teste ou checklist correspondente; resultados de usuários só entram depois da execução do protocolo. O próximo passo é transformar a pendência em issue, teste, decisão de reunião ou evidência de piloto. Este capítulo não promete superioridade nem receita: ele deixa um contrato verificável para que a equipe consiga descobrir se a hipótese se sustenta.

\section{Funcionalidade 10: Chat HTTP}
Na leitura que fizemos do projeto, esta página especifica o comportamento de chat http no produto.

Depois de comparar as alternativas, a equipe decidiu o seguinte: Implementação esperada: Persistência é feita pelo serviço; WebSocket atua como canal de atualização em tempo real.

Na implementação, este requisito é tratado como parte do sistema e não como uma tela isolada. A interface orienta a próxima ação; a API confere identidade, papel, estado e limites; a persistência guarda somente o necessário; e as integrações externas ficam condicionadas às variáveis do ambiente. O escopo continua exclusivamente de software, sem depender de escopo físico.

O fluxo começa quando a pessoa encontra uma necessidade ou uma habilidade que deseja compartilhar. Em seguida, ela confirma a intenção, executa a ação principal e recebe uma resposta que pode ser consultada depois. Para nós, funcionalidade 10: chat http só está concluído quando o caminho é compreensível, existe recuperação para erro e o servidor registra uma evidência coerente.

Para conferir se a decisão funciona, vamos procurar esta evidência: Critério de aceite: a funcionalidade deve funcionar com usuário autenticado quando houver dados pessoais, exibir feedback compreensível e preservar autorização no servidor. O critério de aceite deve ser escrito antes do piloto, com entrada, resultado esperado, responsável, data e ferramenta de verificação. Assim, a documentação separa conformidade técnica, observação de uso e hipótese de impacto.

A validação combina teste automatizado, inspeção manual e uma tarefa representativa. Vamos registrar tempo, erros, conclusão, compreensão e barreiras observadas por teclado, leitor de tela, zoom, contraste, texto, áudio ou conexão instável. O objetivo não é criar uma aparência de certeza, e sim manter um ciclo curto entre observar, corrigir e conferir novamente.

O ponto que ainda merece cuidado é o seguinte: Evidência: rota, serviço, componente, teste ou checklist correspondente; resultados de usuários só entram depois da execução do protocolo. O próximo passo é transformar a pendência em issue, teste, decisão de reunião ou evidência de piloto. Este capítulo não promete superioridade nem receita: ele deixa um contrato verificável para que a equipe consiga descobrir se a hipótese se sustenta.

\section{Funcionalidade 11: Chat WebSocket}
Durante a consolidação do MVP, percebemos que esta página especifica o comportamento de chat websocket no produto.

No protótipo, esse entendimento aparece assim: Implementação esperada: Autenticação ocorre por Authorization ou primeiro frame idToken; senderId é conferido no servidor.

Na implementação, este requisito é tratado como parte do sistema e não como uma tela isolada. A interface orienta a próxima ação; a API confere identidade, papel, estado e limites; a persistência guarda somente o necessário; e as integrações externas ficam condicionadas às variáveis do ambiente. O escopo continua exclusivamente de software, sem depender de escopo físico.

O fluxo começa quando a pessoa encontra uma necessidade ou uma habilidade que deseja compartilhar. Em seguida, ela confirma a intenção, executa a ação principal e recebe uma resposta que pode ser consultada depois. Para nós, funcionalidade 11: chat websocket só está concluído quando o caminho é compreensível, existe recuperação para erro e o servidor registra uma evidência coerente.

A forma mais honesta de verificar este requisito é registrar: Critério de aceite: a funcionalidade deve funcionar com usuário autenticado quando houver dados pessoais, exibir feedback compreensível e preservar autorização no servidor. O critério de aceite deve ser escrito antes do piloto, com entrada, resultado esperado, responsável, data e ferramenta de verificação. Assim, a documentação separa conformidade técnica, observação de uso e hipótese de impacto.

A validação combina teste automatizado, inspeção manual e uma tarefa representativa. Vamos registrar tempo, erros, conclusão, compreensão e barreiras observadas por teclado, leitor de tela, zoom, contraste, texto, áudio ou conexão instável. O objetivo não é criar uma aparência de certeza, e sim manter um ciclo curto entre observar, corrigir e conferir novamente.

A equipe não trata este item como resolvido antes de enfrentar o risco: Evidência: rota, serviço, componente, teste ou checklist correspondente; resultados de usuários só entram depois da execução do protocolo. O próximo passo é transformar a pendência em issue, teste, decisão de reunião ou evidência de piloto. Este capítulo não promete superioridade nem receita: ele deixa um contrato verificável para que a equipe consiga descobrir se a hipótese se sustenta.

\section{Funcionalidade 12: Aula}
A discussão de funcionalidade 12: aula começou por uma situação concreta: esta página especifica o comportamento de aula no produto.

Nossa escolha para este ponto foi direta: Implementação esperada: Agenda e lembrete organizam compromisso, mas não presumem integração paga de vídeo.

Na implementação, este requisito é tratado como parte do sistema e não como uma tela isolada. A interface orienta a próxima ação; a API confere identidade, papel, estado e limites; a persistência guarda somente o necessário; e as integrações externas ficam condicionadas às variáveis do ambiente. O escopo continua exclusivamente de software, sem depender de escopo físico.

O fluxo começa quando a pessoa encontra uma necessidade ou uma habilidade que deseja compartilhar. Em seguida, ela confirma a intenção, executa a ação principal e recebe uma resposta que pode ser consultada depois. Para nós, funcionalidade 12: aula só está concluído quando o caminho é compreensível, existe recuperação para erro e o servidor registra uma evidência coerente.

O aceite não fica apenas no texto. A evidência prevista é: Critério de aceite: a funcionalidade deve funcionar com usuário autenticado quando houver dados pessoais, exibir feedback compreensível e preservar autorização no servidor. O critério de aceite deve ser escrito antes do piloto, com entrada, resultado esperado, responsável, data e ferramenta de verificação. Assim, a documentação separa conformidade técnica, observação de uso e hipótese de impacto.

A validação combina teste automatizado, inspeção manual e uma tarefa representativa. Vamos registrar tempo, erros, conclusão, compreensão e barreiras observadas por teclado, leitor de tela, zoom, contraste, texto, áudio ou conexão instável. O objetivo não é criar uma aparência de certeza, e sim manter um ciclo curto entre observar, corrigir e conferir novamente.

Também registramos uma limitação, porque ela faz parte do resultado: Evidência: rota, serviço, componente, teste ou checklist correspondente; resultados de usuários só entram depois da execução do protocolo. O próximo passo é transformar a pendência em issue, teste, decisão de reunião ou evidência de piloto. Este capítulo não promete superioridade nem receita: ele deixa um contrato verificável para que a equipe consiga descobrir se a hipótese se sustenta.

\section{Funcionalidade 13: Vídeo}
A discussão de funcionalidade 13: vídeo começou por uma situação concreta: esta página especifica o comportamento de vídeo no produto.

Nossa escolha para este ponto foi direta: Implementação esperada: LiveKit é opcional e depende de segredos; a funcionalidade pode ficar desligada no MVP.

Na implementação, este requisito é tratado como parte do sistema e não como uma tela isolada. A interface orienta a próxima ação; a API confere identidade, papel, estado e limites; a persistência guarda somente o necessário; e as integrações externas ficam condicionadas às variáveis do ambiente. O escopo continua exclusivamente de software, sem depender de escopo físico.

O fluxo começa quando a pessoa encontra uma necessidade ou uma habilidade que deseja compartilhar. Em seguida, ela confirma a intenção, executa a ação principal e recebe uma resposta que pode ser consultada depois. Para nós, funcionalidade 13: vídeo só está concluído quando o caminho é compreensível, existe recuperação para erro e o servidor registra uma evidência coerente.

O aceite não fica apenas no texto. A evidência prevista é: Critério de aceite: a funcionalidade deve funcionar com usuário autenticado quando houver dados pessoais, exibir feedback compreensível e preservar autorização no servidor. O critério de aceite deve ser escrito antes do piloto, com entrada, resultado esperado, responsável, data e ferramenta de verificação. Assim, a documentação separa conformidade técnica, observação de uso e hipótese de impacto.

A validação combina teste automatizado, inspeção manual e uma tarefa representativa. Vamos registrar tempo, erros, conclusão, compreensão e barreiras observadas por teclado, leitor de tela, zoom, contraste, texto, áudio ou conexão instável. O objetivo não é criar uma aparência de certeza, e sim manter um ciclo curto entre observar, corrigir e conferir novamente.

Também registramos uma limitação, porque ela faz parte do resultado: Evidência: rota, serviço, componente, teste ou checklist correspondente; resultados de usuários só entram depois da execução do protocolo. O próximo passo é transformar a pendência em issue, teste, decisão de reunião ou evidência de piloto. Este capítulo não promete superioridade nem receita: ele deixa um contrato verificável para que a equipe consiga descobrir se a hipótese se sustenta.

\section{Funcionalidade 14: Cursos}
Na leitura que fizemos do projeto, esta página especifica o comportamento de cursos no produto.

Depois de comparar as alternativas, a equipe decidiu o seguinte: Implementação esperada: Trilhas representam conteúdo e evolução; materiais devem ter alternativa textual e legenda.

Na implementação, este requisito é tratado como parte do sistema e não como uma tela isolada. A interface orienta a próxima ação; a API confere identidade, papel, estado e limites; a persistência guarda somente o necessário; e as integrações externas ficam condicionadas às variáveis do ambiente. O escopo continua exclusivamente de software, sem depender de escopo físico.

O fluxo começa quando a pessoa encontra uma necessidade ou uma habilidade que deseja compartilhar. Em seguida, ela confirma a intenção, executa a ação principal e recebe uma resposta que pode ser consultada depois. Para nós, funcionalidade 14: cursos só está concluído quando o caminho é compreensível, existe recuperação para erro e o servidor registra uma evidência coerente.

Para conferir se a decisão funciona, vamos procurar esta evidência: Critério de aceite: a funcionalidade deve funcionar com usuário autenticado quando houver dados pessoais, exibir feedback compreensível e preservar autorização no servidor. O critério de aceite deve ser escrito antes do piloto, com entrada, resultado esperado, responsável, data e ferramenta de verificação. Assim, a documentação separa conformidade técnica, observação de uso e hipótese de impacto.

A validação combina teste automatizado, inspeção manual e uma tarefa representativa. Vamos registrar tempo, erros, conclusão, compreensão e barreiras observadas por teclado, leitor de tela, zoom, contraste, texto, áudio ou conexão instável. O objetivo não é criar uma aparência de certeza, e sim manter um ciclo curto entre observar, corrigir e conferir novamente.

O ponto que ainda merece cuidado é o seguinte: Evidência: rota, serviço, componente, teste ou checklist correspondente; resultados de usuários só entram depois da execução do protocolo. O próximo passo é transformar a pendência em issue, teste, decisão de reunião ou evidência de piloto. Este capítulo não promete superioridade nem receita: ele deixa um contrato verificável para que a equipe consiga descobrir se a hipótese se sustenta.

\section{Funcionalidade 15: Ranking}
A discussão de funcionalidade 15: ranking começou por uma situação concreta: esta página especifica o comportamento de ranking no produto.

Nossa escolha para este ponto foi direta: Implementação esperada: Reputação auxilia descoberta, mas não deve excluir perfis novos ou grupos minoritários.

Na implementação, este requisito é tratado como parte do sistema e não como uma tela isolada. A interface orienta a próxima ação; a API confere identidade, papel, estado e limites; a persistência guarda somente o necessário; e as integrações externas ficam condicionadas às variáveis do ambiente. O escopo continua exclusivamente de software, sem depender de escopo físico.

O fluxo começa quando a pessoa encontra uma necessidade ou uma habilidade que deseja compartilhar. Em seguida, ela confirma a intenção, executa a ação principal e recebe uma resposta que pode ser consultada depois. Para nós, funcionalidade 15: ranking só está concluído quando o caminho é compreensível, existe recuperação para erro e o servidor registra uma evidência coerente.

O aceite não fica apenas no texto. A evidência prevista é: Critério de aceite: a funcionalidade deve funcionar com usuário autenticado quando houver dados pessoais, exibir feedback compreensível e preservar autorização no servidor. O critério de aceite deve ser escrito antes do piloto, com entrada, resultado esperado, responsável, data e ferramenta de verificação. Assim, a documentação separa conformidade técnica, observação de uso e hipótese de impacto.

A validação combina teste automatizado, inspeção manual e uma tarefa representativa. Vamos registrar tempo, erros, conclusão, compreensão e barreiras observadas por teclado, leitor de tela, zoom, contraste, texto, áudio ou conexão instável. O objetivo não é criar uma aparência de certeza, e sim manter um ciclo curto entre observar, corrigir e conferir novamente.

Também registramos uma limitação, porque ela faz parte do resultado: Evidência: rota, serviço, componente, teste ou checklist correspondente; resultados de usuários só entram depois da execução do protocolo. O próximo passo é transformar a pendência em issue, teste, decisão de reunião ou evidência de piloto. Este capítulo não promete superioridade nem receita: ele deixa um contrato verificável para que a equipe consiga descobrir se a hipótese se sustenta.

\section{Funcionalidade 16: ELOS}
Na leitura que fizemos do projeto, esta página especifica o comportamento de elos no produto.

Depois de comparar as alternativas, a equipe decidiu o seguinte: Implementação esperada: Chatbot responde em linguagem simples e comunica limites; não decide contratação ou diagnóstico.

Na implementação, este requisito é tratado como parte do sistema e não como uma tela isolada. A interface orienta a próxima ação; a API confere identidade, papel, estado e limites; a persistência guarda somente o necessário; e as integrações externas ficam condicionadas às variáveis do ambiente. O escopo continua exclusivamente de software, sem depender de escopo físico.

O fluxo começa quando a pessoa encontra uma necessidade ou uma habilidade que deseja compartilhar. Em seguida, ela confirma a intenção, executa a ação principal e recebe uma resposta que pode ser consultada depois. Para nós, funcionalidade 16: elos só está concluído quando o caminho é compreensível, existe recuperação para erro e o servidor registra uma evidência coerente.

Para conferir se a decisão funciona, vamos procurar esta evidência: Critério de aceite: a funcionalidade deve funcionar com usuário autenticado quando houver dados pessoais, exibir feedback compreensível e preservar autorização no servidor. O critério de aceite deve ser escrito antes do piloto, com entrada, resultado esperado, responsável, data e ferramenta de verificação. Assim, a documentação separa conformidade técnica, observação de uso e hipótese de impacto.

A validação combina teste automatizado, inspeção manual e uma tarefa representativa. Vamos registrar tempo, erros, conclusão, compreensão e barreiras observadas por teclado, leitor de tela, zoom, contraste, texto, áudio ou conexão instável. O objetivo não é criar uma aparência de certeza, e sim manter um ciclo curto entre observar, corrigir e conferir novamente.

O ponto que ainda merece cuidado é o seguinte: Evidência: rota, serviço, componente, teste ou checklist correspondente; resultados de usuários só entram depois da execução do protocolo. O próximo passo é transformar a pendência em issue, teste, decisão de reunião ou evidência de piloto. Este capítulo não promete superioridade nem receita: ele deixa um contrato verificável para que a equipe consiga descobrir se a hipótese se sustenta.

\section{Módulos comerciais fora do escopo}

A SkillSwap não cobra pelos recursos entregues. Carteira, SwapCoins, pacotes, checkout, transferências, Stripe, gorjetas e recursos premium não fazem parte do fluxo desta versão. A tela financeira foi retirada da navegação e as rotas comerciais são bloqueadas na API.

Essa decisão deixa o acesso igual para todas as pessoas e mantém o foco na troca de conhecimento e na acessibilidade. O código legado desses módulos não é chamado pelo frontend e permanece apenas como registro da evolução do projeto.
\section{Funcionalidade 22: Administração}
Durante a consolidação do MVP, percebemos que esta página especifica o comportamento de administração no produto.

No protótipo, esse entendimento aparece assim: Implementação esperada: Painel agrupa métricas e ações com papel ADMIN e auditoria.

Na implementação, este requisito é tratado como parte do sistema e não como uma tela isolada. A interface orienta a próxima ação; a API confere identidade, papel, estado e limites; a persistência guarda somente o necessário; e as integrações externas ficam condicionadas às variáveis do ambiente. O escopo continua exclusivamente de software, sem depender de escopo físico.

O fluxo começa quando a pessoa encontra uma necessidade ou uma habilidade que deseja compartilhar. Em seguida, ela confirma a intenção, executa a ação principal e recebe uma resposta que pode ser consultada depois. Para nós, funcionalidade 22: administração só está concluído quando o caminho é compreensível, existe recuperação para erro e o servidor registra uma evidência coerente.

A forma mais honesta de verificar este requisito é registrar: Critério de aceite: a funcionalidade deve funcionar com usuário autenticado quando houver dados pessoais, exibir feedback compreensível e preservar autorização no servidor. O critério de aceite deve ser escrito antes do piloto, com entrada, resultado esperado, responsável, data e ferramenta de verificação. Assim, a documentação separa conformidade técnica, observação de uso e hipótese de impacto.

A validação combina teste automatizado, inspeção manual e uma tarefa representativa. Vamos registrar tempo, erros, conclusão, compreensão e barreiras observadas por teclado, leitor de tela, zoom, contraste, texto, áudio ou conexão instável. O objetivo não é criar uma aparência de certeza, e sim manter um ciclo curto entre observar, corrigir e conferir novamente.

A equipe não trata este item como resolvido antes de enfrentar o risco: Evidência: rota, serviço, componente, teste ou checklist correspondente; resultados de usuários só entram depois da execução do protocolo. O próximo passo é transformar a pendência em issue, teste, decisão de reunião ou evidência de piloto. Este capítulo não promete superioridade nem receita: ele deixa um contrato verificável para que a equipe consiga descobrir se a hipótese se sustenta.

\section{Funcionalidade 23: FCM}
A discussão de funcionalidade 23: fcm começou por uma situação concreta: esta página especifica o comportamento de fcm no produto.

Nossa escolha para este ponto foi direta: Implementação esperada: Token é registrado pelo usuário autenticado e pode ser revogado.

Na implementação, este requisito é tratado como parte do sistema e não como uma tela isolada. A interface orienta a próxima ação; a API confere identidade, papel, estado e limites; a persistência guarda somente o necessário; e as integrações externas ficam condicionadas às variáveis do ambiente. O escopo continua exclusivamente de software, sem depender de escopo físico.

O fluxo começa quando a pessoa encontra uma necessidade ou uma habilidade que deseja compartilhar. Em seguida, ela confirma a intenção, executa a ação principal e recebe uma resposta que pode ser consultada depois. Para nós, funcionalidade 23: fcm só está concluído quando o caminho é compreensível, existe recuperação para erro e o servidor registra uma evidência coerente.

O aceite não fica apenas no texto. A evidência prevista é: Critério de aceite: a funcionalidade deve funcionar com usuário autenticado quando houver dados pessoais, exibir feedback compreensível e preservar autorização no servidor. O critério de aceite deve ser escrito antes do piloto, com entrada, resultado esperado, responsável, data e ferramenta de verificação. Assim, a documentação separa conformidade técnica, observação de uso e hipótese de impacto.

A validação combina teste automatizado, inspeção manual e uma tarefa representativa. Vamos registrar tempo, erros, conclusão, compreensão e barreiras observadas por teclado, leitor de tela, zoom, contraste, texto, áudio ou conexão instável. O objetivo não é criar uma aparência de certeza, e sim manter um ciclo curto entre observar, corrigir e conferir novamente.

Também registramos uma limitação, porque ela faz parte do resultado: Evidência: rota, serviço, componente, teste ou checklist correspondente; resultados de usuários só entram depois da execução do protocolo. O próximo passo é transformar a pendência em issue, teste, decisão de reunião ou evidência de piloto. Este capítulo não promete superioridade nem receita: ele deixa um contrato verificável para que a equipe consiga descobrir se a hipótese se sustenta.

\section{Funcionalidade 24: Notificações}
Na leitura que fizemos do projeto, esta página especifica o comportamento de notificações no produto.

Depois de comparar as alternativas, a equipe decidiu o seguinte: Implementação esperada: Histórico unread/read complementa push e email.

Na implementação, este requisito é tratado como parte do sistema e não como uma tela isolada. A interface orienta a próxima ação; a API confere identidade, papel, estado e limites; a persistência guarda somente o necessário; e as integrações externas ficam condicionadas às variáveis do ambiente. O escopo continua exclusivamente de software, sem depender de escopo físico.

O fluxo começa quando a pessoa encontra uma necessidade ou uma habilidade que deseja compartilhar. Em seguida, ela confirma a intenção, executa a ação principal e recebe uma resposta que pode ser consultada depois. Para nós, funcionalidade 24: notificações só está concluído quando o caminho é compreensível, existe recuperação para erro e o servidor registra uma evidência coerente.

Para conferir se a decisão funciona, vamos procurar esta evidência: Critério de aceite: a funcionalidade deve funcionar com usuário autenticado quando houver dados pessoais, exibir feedback compreensível e preservar autorização no servidor. O critério de aceite deve ser escrito antes do piloto, com entrada, resultado esperado, responsável, data e ferramenta de verificação. Assim, a documentação separa conformidade técnica, observação de uso e hipótese de impacto.

A validação combina teste automatizado, inspeção manual e uma tarefa representativa. Vamos registrar tempo, erros, conclusão, compreensão e barreiras observadas por teclado, leitor de tela, zoom, contraste, texto, áudio ou conexão instável. O objetivo não é criar uma aparência de certeza, e sim manter um ciclo curto entre observar, corrigir e conferir novamente.

O ponto que ainda merece cuidado é o seguinte: Evidência: rota, serviço, componente, teste ou checklist correspondente; resultados de usuários só entram depois da execução do protocolo. O próximo passo é transformar a pendência em issue, teste, decisão de reunião ou evidência de piloto. Este capítulo não promete superioridade nem receita: ele deixa um contrato verificável para que a equipe consiga descobrir se a hipótese se sustenta.

\section{Funcionalidade 25: Actuator}
Na leitura que fizemos do projeto, esta página especifica o comportamento de actuator no produto.

Depois de comparar as alternativas, a equipe decidiu o seguinte: Implementação esperada: Health e Prometheus sustentam operação e diagnóstico.

Na implementação, este requisito é tratado como parte do sistema e não como uma tela isolada. A interface orienta a próxima ação; a API confere identidade, papel, estado e limites; a persistência guarda somente o necessário; e as integrações externas ficam condicionadas às variáveis do ambiente. O escopo continua exclusivamente de software, sem depender de escopo físico.

O fluxo começa quando a pessoa encontra uma necessidade ou uma habilidade que deseja compartilhar. Em seguida, ela confirma a intenção, executa a ação principal e recebe uma resposta que pode ser consultada depois. Para nós, funcionalidade 25: actuator só está concluído quando o caminho é compreensível, existe recuperação para erro e o servidor registra uma evidência coerente.

Para conferir se a decisão funciona, vamos procurar esta evidência: Critério de aceite: a funcionalidade deve funcionar com usuário autenticado quando houver dados pessoais, exibir feedback compreensível e preservar autorização no servidor. O critério de aceite deve ser escrito antes do piloto, com entrada, resultado esperado, responsável, data e ferramenta de verificação. Assim, a documentação separa conformidade técnica, observação de uso e hipótese de impacto.

A validação combina teste automatizado, inspeção manual e uma tarefa representativa. Vamos registrar tempo, erros, conclusão, compreensão e barreiras observadas por teclado, leitor de tela, zoom, contraste, texto, áudio ou conexão instável. O objetivo não é criar uma aparência de certeza, e sim manter um ciclo curto entre observar, corrigir e conferir novamente.

O ponto que ainda merece cuidado é o seguinte: Evidência: rota, serviço, componente, teste ou checklist correspondente; resultados de usuários só entram depois da execução do protocolo. O próximo passo é transformar a pendência em issue, teste, decisão de reunião ou evidência de piloto. Este capítulo não promete superioridade nem receita: ele deixa um contrato verificável para que a equipe consiga descobrir se a hipótese se sustenta.

\section{Funcionalidade 26: Caffeine}
Durante a consolidação do MVP, percebemos que esta página especifica o comportamento de caffeine no produto.

No protótipo, esse entendimento aparece assim: Implementação esperada: Cache local reduz custo de Redis no primeiro deploy.

Na implementação, este requisito é tratado como parte do sistema e não como uma tela isolada. A interface orienta a próxima ação; a API confere identidade, papel, estado e limites; a persistência guarda somente o necessário; e as integrações externas ficam condicionadas às variáveis do ambiente. O escopo continua exclusivamente de software, sem depender de escopo físico.

O fluxo começa quando a pessoa encontra uma necessidade ou uma habilidade que deseja compartilhar. Em seguida, ela confirma a intenção, executa a ação principal e recebe uma resposta que pode ser consultada depois. Para nós, funcionalidade 26: caffeine só está concluído quando o caminho é compreensível, existe recuperação para erro e o servidor registra uma evidência coerente.

A forma mais honesta de verificar este requisito é registrar: Critério de aceite: a funcionalidade deve funcionar com usuário autenticado quando houver dados pessoais, exibir feedback compreensível e preservar autorização no servidor. O critério de aceite deve ser escrito antes do piloto, com entrada, resultado esperado, responsável, data e ferramenta de verificação. Assim, a documentação separa conformidade técnica, observação de uso e hipótese de impacto.

A validação combina teste automatizado, inspeção manual e uma tarefa representativa. Vamos registrar tempo, erros, conclusão, compreensão e barreiras observadas por teclado, leitor de tela, zoom, contraste, texto, áudio ou conexão instável. O objetivo não é criar uma aparência de certeza, e sim manter um ciclo curto entre observar, corrigir e conferir novamente.

A equipe não trata este item como resolvido antes de enfrentar o risco: Evidência: rota, serviço, componente, teste ou checklist correspondente; resultados de usuários só entram depois da execução do protocolo. O próximo passo é transformar a pendência em issue, teste, decisão de reunião ou evidência de piloto. Este capítulo não promete superioridade nem receita: ele deixa um contrato verificável para que a equipe consiga descobrir se a hipótese se sustenta.

\section{Funcionalidade 27: RabbitMQ}
Durante a consolidação do MVP, percebemos que esta página especifica o comportamento de rabbitmq no produto.

No protótipo, esse entendimento aparece assim: Implementação esperada: O processamento comercial está fora do escopo e as rotas correspondentes permanecem bloqueadas.

Na implementação, este requisito é tratado como parte do sistema e não como uma tela isolada. A interface orienta a próxima ação; a API confere identidade, papel, estado e limites; a persistência guarda somente o necessário; e as integrações externas ficam condicionadas às variáveis do ambiente. O escopo continua exclusivamente de software, sem depender de escopo físico.

O fluxo começa quando a pessoa encontra uma necessidade ou uma habilidade que deseja compartilhar. Em seguida, ela confirma a intenção, executa a ação principal e recebe uma resposta que pode ser consultada depois. Para nós, funcionalidade 27: rabbitmq só está concluído quando o caminho é compreensível, existe recuperação para erro e o servidor registra uma evidência coerente.

A forma mais honesta de verificar este requisito é registrar: Critério de aceite: a funcionalidade deve funcionar com usuário autenticado quando houver dados pessoais, exibir feedback compreensível e preservar autorização no servidor. O critério de aceite deve ser escrito antes do piloto, com entrada, resultado esperado, responsável, data e ferramenta de verificação. Assim, a documentação separa conformidade técnica, observação de uso e hipótese de impacto.

A validação combina teste automatizado, inspeção manual e uma tarefa representativa. Vamos registrar tempo, erros, conclusão, compreensão e barreiras observadas por teclado, leitor de tela, zoom, contraste, texto, áudio ou conexão instável. O objetivo não é criar uma aparência de certeza, e sim manter um ciclo curto entre observar, corrigir e conferir novamente.

A equipe não trata este item como resolvido antes de enfrentar o risco: Evidência: rota, serviço, componente, teste ou checklist correspondente; resultados de usuários só entram depois da execução do protocolo. O próximo passo é transformar a pendência em issue, teste, decisão de reunião ou evidência de piloto. Este capítulo não promete superioridade nem receita: ele deixa um contrato verificável para que a equipe consiga descobrir se a hipótese se sustenta.

\section{Funcionalidade 28: Flyway}
Durante a consolidação do MVP, percebemos que esta página especifica o comportamento de flyway no produto.

No protótipo, esse entendimento aparece assim: Implementação esperada: Migrações versionam o contrato de banco e impedem drift.

Na implementação, este requisito é tratado como parte do sistema e não como uma tela isolada. A interface orienta a próxima ação; a API confere identidade, papel, estado e limites; a persistência guarda somente o necessário; e as integrações externas ficam condicionadas às variáveis do ambiente. O escopo continua exclusivamente de software, sem depender de escopo físico.

O fluxo começa quando a pessoa encontra uma necessidade ou uma habilidade que deseja compartilhar. Em seguida, ela confirma a intenção, executa a ação principal e recebe uma resposta que pode ser consultada depois. Para nós, funcionalidade 28: flyway só está concluído quando o caminho é compreensível, existe recuperação para erro e o servidor registra uma evidência coerente.

A forma mais honesta de verificar este requisito é registrar: Critério de aceite: a funcionalidade deve funcionar com usuário autenticado quando houver dados pessoais, exibir feedback compreensível e preservar autorização no servidor. O critério de aceite deve ser escrito antes do piloto, com entrada, resultado esperado, responsável, data e ferramenta de verificação. Assim, a documentação separa conformidade técnica, observação de uso e hipótese de impacto.

A validação combina teste automatizado, inspeção manual e uma tarefa representativa. Vamos registrar tempo, erros, conclusão, compreensão e barreiras observadas por teclado, leitor de tela, zoom, contraste, texto, áudio ou conexão instável. O objetivo não é criar uma aparência de certeza, e sim manter um ciclo curto entre observar, corrigir e conferir novamente.

A equipe não trata este item como resolvido antes de enfrentar o risco: Evidência: rota, serviço, componente, teste ou checklist correspondente; resultados de usuários só entram depois da execução do protocolo. O próximo passo é transformar a pendência em issue, teste, decisão de reunião ou evidência de piloto. Este capítulo não promete superioridade nem receita: ele deixa um contrato verificável para que a equipe consiga descobrir se a hipótese se sustenta.

\section{Funcionalidade 29: CORS}
Durante a consolidação do MVP, percebemos que esta página especifica o comportamento de cors no produto.

No protótipo, esse entendimento aparece assim: Implementação esperada: Origem permitida vem de variável; nunca usar wildcard com dados autenticados.

Na implementação, este requisito é tratado como parte do sistema e não como uma tela isolada. A interface orienta a próxima ação; a API confere identidade, papel, estado e limites; a persistência guarda somente o necessário; e as integrações externas ficam condicionadas às variáveis do ambiente. O escopo continua exclusivamente de software, sem depender de escopo físico.

O fluxo começa quando a pessoa encontra uma necessidade ou uma habilidade que deseja compartilhar. Em seguida, ela confirma a intenção, executa a ação principal e recebe uma resposta que pode ser consultada depois. Para nós, funcionalidade 29: cors só está concluído quando o caminho é compreensível, existe recuperação para erro e o servidor registra uma evidência coerente.

A forma mais honesta de verificar este requisito é registrar: Critério de aceite: a funcionalidade deve funcionar com usuário autenticado quando houver dados pessoais, exibir feedback compreensível e preservar autorização no servidor. O critério de aceite deve ser escrito antes do piloto, com entrada, resultado esperado, responsável, data e ferramenta de verificação. Assim, a documentação separa conformidade técnica, observação de uso e hipótese de impacto.

A validação combina teste automatizado, inspeção manual e uma tarefa representativa. Vamos registrar tempo, erros, conclusão, compreensão e barreiras observadas por teclado, leitor de tela, zoom, contraste, texto, áudio ou conexão instável. O objetivo não é criar uma aparência de certeza, e sim manter um ciclo curto entre observar, corrigir e conferir novamente.

A equipe não trata este item como resolvido antes de enfrentar o risco: Evidência: rota, serviço, componente, teste ou checklist correspondente; resultados de usuários só entram depois da execução do protocolo. O próximo passo é transformar a pendência em issue, teste, decisão de reunião ou evidência de piloto. Este capítulo não promete superioridade nem receita: ele deixa um contrato verificável para que a equipe consiga descobrir se a hipótese se sustenta.

\section{Funcionalidade 30: Logs}
A discussão de funcionalidade 30: logs começou por uma situação concreta: esta página especifica o comportamento de logs no produto.

Nossa escolha para este ponto foi direta: Implementação esperada: Logs não incluem token, senha, chave privada ou payload sensível.

Na implementação, este requisito é tratado como parte do sistema e não como uma tela isolada. A interface orienta a próxima ação; a API confere identidade, papel, estado e limites; a persistência guarda somente o necessário; e as integrações externas ficam condicionadas às variáveis do ambiente. O escopo continua exclusivamente de software, sem depender de escopo físico.

O fluxo começa quando a pessoa encontra uma necessidade ou uma habilidade que deseja compartilhar. Em seguida, ela confirma a intenção, executa a ação principal e recebe uma resposta que pode ser consultada depois. Para nós, funcionalidade 30: logs só está concluído quando o caminho é compreensível, existe recuperação para erro e o servidor registra uma evidência coerente.

O aceite não fica apenas no texto. A evidência prevista é: Critério de aceite: a funcionalidade deve funcionar com usuário autenticado quando houver dados pessoais, exibir feedback compreensível e preservar autorização no servidor. O critério de aceite deve ser escrito antes do piloto, com entrada, resultado esperado, responsável, data e ferramenta de verificação. Assim, a documentação separa conformidade técnica, observação de uso e hipótese de impacto.

A validação combina teste automatizado, inspeção manual e uma tarefa representativa. Vamos registrar tempo, erros, conclusão, compreensão e barreiras observadas por teclado, leitor de tela, zoom, contraste, texto, áudio ou conexão instável. O objetivo não é criar uma aparência de certeza, e sim manter um ciclo curto entre observar, corrigir e conferir novamente.

Também registramos uma limitação, porque ela faz parte do resultado: Evidência: rota, serviço, componente, teste ou checklist correspondente; resultados de usuários só entram depois da execução do protocolo. O próximo passo é transformar a pendência em issue, teste, decisão de reunião ou evidência de piloto. Este capítulo não promete superioridade nem receita: ele deixa um contrato verificável para que a equipe consiga descobrir se a hipótese se sustenta.

\section{Funcionalidade 31: Documentação}
Durante a consolidação do MVP, percebemos que esta página especifica o comportamento de documentação no produto.

No protótipo, esse entendimento aparece assim: Implementação esperada: README, checklist, arquitetura, fontes e contrato devem acompanhar o código.

Na implementação, este requisito é tratado como parte do sistema e não como uma tela isolada. A interface orienta a próxima ação; a API confere identidade, papel, estado e limites; a persistência guarda somente o necessário; e as integrações externas ficam condicionadas às variáveis do ambiente. O escopo continua exclusivamente de software, sem depender de escopo físico.

O fluxo começa quando a pessoa encontra uma necessidade ou uma habilidade que deseja compartilhar. Em seguida, ela confirma a intenção, executa a ação principal e recebe uma resposta que pode ser consultada depois. Para nós, funcionalidade 31: documentação só está concluído quando o caminho é compreensível, existe recuperação para erro e o servidor registra uma evidência coerente.

A forma mais honesta de verificar este requisito é registrar: Critério de aceite: a funcionalidade deve funcionar com usuário autenticado quando houver dados pessoais, exibir feedback compreensível e preservar autorização no servidor. O critério de aceite deve ser escrito antes do piloto, com entrada, resultado esperado, responsável, data e ferramenta de verificação. Assim, a documentação separa conformidade técnica, observação de uso e hipótese de impacto.

A validação combina teste automatizado, inspeção manual e uma tarefa representativa. Vamos registrar tempo, erros, conclusão, compreensão e barreiras observadas por teclado, leitor de tela, zoom, contraste, texto, áudio ou conexão instável. O objetivo não é criar uma aparência de certeza, e sim manter um ciclo curto entre observar, corrigir e conferir novamente.

A equipe não trata este item como resolvido antes de enfrentar o risco: Evidência: rota, serviço, componente, teste ou checklist correspondente; resultados de usuários só entram depois da execução do protocolo. O próximo passo é transformar a pendência em issue, teste, decisão de reunião ou evidência de piloto. Este capítulo não promete superioridade nem receita: ele deixa um contrato verificável para que a equipe consiga descobrir se a hipótese se sustenta.

\section{Teste 01: Token ausente}
A discussão de teste 01: token ausente começou por uma situação concreta: objetivo: validar token ausente no ambiente de demonstração.

Nossa escolha para este ponto foi direta: Resultado esperado: API retorna 401 e não cria sessão local.

Na implementação, este requisito é tratado como parte do sistema e não como uma tela isolada. A interface orienta a próxima ação; a API confere identidade, papel, estado e limites; a persistência guarda somente o necessário; e as integrações externas ficam condicionadas às variáveis do ambiente. O escopo continua exclusivamente de software, sem depender de escopo físico.

O fluxo começa quando a pessoa encontra uma necessidade ou uma habilidade que deseja compartilhar. Em seguida, ela confirma a intenção, executa a ação principal e recebe uma resposta que pode ser consultada depois. Para nós, teste 01: token ausente só está concluído quando o caminho é compreensível, existe recuperação para erro e o servidor registra uma evidência coerente.

O aceite não fica apenas no texto. A evidência prevista é: Procedimento: preparar dados sintéticos, executar a ação, verificar resposta HTTP ou evento, conferir banco/ledger quando aplicável e registrar versão do código. O critério de aceite deve ser escrito antes do piloto, com entrada, resultado esperado, responsável, data e ferramenta de verificação. Assim, a documentação separa conformidade técnica, observação de uso e hipótese de impacto.

A validação combina teste automatizado, inspeção manual e uma tarefa representativa. Vamos registrar tempo, erros, conclusão, compreensão e barreiras observadas por teclado, leitor de tela, zoom, contraste, texto, áudio ou conexão instável. O objetivo não é criar uma aparência de certeza, e sim manter um ciclo curto entre observar, corrigir e conferir novamente.

Também registramos uma limitação, porque ela faz parte do resultado: Classificação: aprovado, reprovado ou bloqueado. Não marcar como aprovado sem evidência reproduzível; bloquear quando depender de credencial ou serviço externo ainda não configurado. O próximo passo é transformar a pendência em issue, teste, decisão de reunião ou evidência de piloto. Este capítulo não promete superioridade nem receita: ele deixa um contrato verificável para que a equipe consiga descobrir se a hipótese se sustenta.

\section{Teste 02: Token expirado}
Na leitura que fizemos do projeto, objetivo: validar token expirado no ambiente de demonstração.

Depois de comparar as alternativas, a equipe decidiu o seguinte: Resultado esperado: API retorna 401 e o frontend renova pelo Firebase.

Na implementação, este requisito é tratado como parte do sistema e não como uma tela isolada. A interface orienta a próxima ação; a API confere identidade, papel, estado e limites; a persistência guarda somente o necessário; e as integrações externas ficam condicionadas às variáveis do ambiente. O escopo continua exclusivamente de software, sem depender de escopo físico.

O fluxo começa quando a pessoa encontra uma necessidade ou uma habilidade que deseja compartilhar. Em seguida, ela confirma a intenção, executa a ação principal e recebe uma resposta que pode ser consultada depois. Para nós, teste 02: token expirado só está concluído quando o caminho é compreensível, existe recuperação para erro e o servidor registra uma evidência coerente.

Para conferir se a decisão funciona, vamos procurar esta evidência: Procedimento: preparar dados sintéticos, executar a ação, verificar resposta HTTP ou evento, conferir banco/ledger quando aplicável e registrar versão do código. O critério de aceite deve ser escrito antes do piloto, com entrada, resultado esperado, responsável, data e ferramenta de verificação. Assim, a documentação separa conformidade técnica, observação de uso e hipótese de impacto.

A validação combina teste automatizado, inspeção manual e uma tarefa representativa. Vamos registrar tempo, erros, conclusão, compreensão e barreiras observadas por teclado, leitor de tela, zoom, contraste, texto, áudio ou conexão instável. O objetivo não é criar uma aparência de certeza, e sim manter um ciclo curto entre observar, corrigir e conferir novamente.

O ponto que ainda merece cuidado é o seguinte: Classificação: aprovado, reprovado ou bloqueado. Não marcar como aprovado sem evidência reproduzível; bloquear quando depender de credencial ou serviço externo ainda não configurado. O próximo passo é transformar a pendência em issue, teste, decisão de reunião ou evidência de piloto. Este capítulo não promete superioridade nem receita: ele deixa um contrato verificável para que a equipe consiga descobrir se a hipótese se sustenta.

\section{Teste 03: Usuário banido}
A discussão de teste 03: usuário banido começou por uma situação concreta: objetivo: validar usuário banido no ambiente de demonstração.

Nossa escolha para este ponto foi direta: Resultado esperado: Filtro bloqueia rotas autenticadas e registra decisão.

Na implementação, este requisito é tratado como parte do sistema e não como uma tela isolada. A interface orienta a próxima ação; a API confere identidade, papel, estado e limites; a persistência guarda somente o necessário; e as integrações externas ficam condicionadas às variáveis do ambiente. O escopo continua exclusivamente de software, sem depender de escopo físico.

O fluxo começa quando a pessoa encontra uma necessidade ou uma habilidade que deseja compartilhar. Em seguida, ela confirma a intenção, executa a ação principal e recebe uma resposta que pode ser consultada depois. Para nós, teste 03: usuário banido só está concluído quando o caminho é compreensível, existe recuperação para erro e o servidor registra uma evidência coerente.

O aceite não fica apenas no texto. A evidência prevista é: Procedimento: preparar dados sintéticos, executar a ação, verificar resposta HTTP ou evento, conferir banco/ledger quando aplicável e registrar versão do código. O critério de aceite deve ser escrito antes do piloto, com entrada, resultado esperado, responsável, data e ferramenta de verificação. Assim, a documentação separa conformidade técnica, observação de uso e hipótese de impacto.

A validação combina teste automatizado, inspeção manual e uma tarefa representativa. Vamos registrar tempo, erros, conclusão, compreensão e barreiras observadas por teclado, leitor de tela, zoom, contraste, texto, áudio ou conexão instável. O objetivo não é criar uma aparência de certeza, e sim manter um ciclo curto entre observar, corrigir e conferir novamente.

Também registramos uma limitação, porque ela faz parte do resultado: Classificação: aprovado, reprovado ou bloqueado. Não marcar como aprovado sem evidência reproduzível; bloquear quando depender de credencial ou serviço externo ainda não configurado. O próximo passo é transformar a pendência em issue, teste, decisão de reunião ou evidência de piloto. Este capítulo não promete superioridade nem receita: ele deixa um contrato verificável para que a equipe consiga descobrir se a hipótese se sustenta.

\section{Teste 04: Papel admin}
Durante a consolidação do MVP, percebemos que objetivo: validar papel admin no ambiente de demonstração.

No protótipo, esse entendimento aparece assim: Resultado esperado: Endpoint administrativo aceita ADMIN e rejeita USER com 403.

Na implementação, este requisito é tratado como parte do sistema e não como uma tela isolada. A interface orienta a próxima ação; a API confere identidade, papel, estado e limites; a persistência guarda somente o necessário; e as integrações externas ficam condicionadas às variáveis do ambiente. O escopo continua exclusivamente de software, sem depender de escopo físico.

O fluxo começa quando a pessoa encontra uma necessidade ou uma habilidade que deseja compartilhar. Em seguida, ela confirma a intenção, executa a ação principal e recebe uma resposta que pode ser consultada depois. Para nós, teste 04: papel admin só está concluído quando o caminho é compreensível, existe recuperação para erro e o servidor registra uma evidência coerente.

A forma mais honesta de verificar este requisito é registrar: Procedimento: preparar dados sintéticos, executar a ação, verificar resposta HTTP ou evento, conferir banco/ledger quando aplicável e registrar versão do código. O critério de aceite deve ser escrito antes do piloto, com entrada, resultado esperado, responsável, data e ferramenta de verificação. Assim, a documentação separa conformidade técnica, observação de uso e hipótese de impacto.

A validação combina teste automatizado, inspeção manual e uma tarefa representativa. Vamos registrar tempo, erros, conclusão, compreensão e barreiras observadas por teclado, leitor de tela, zoom, contraste, texto, áudio ou conexão instável. O objetivo não é criar uma aparência de certeza, e sim manter um ciclo curto entre observar, corrigir e conferir novamente.

A equipe não trata este item como resolvido antes de enfrentar o risco: Classificação: aprovado, reprovado ou bloqueado. Não marcar como aprovado sem evidência reproduzível; bloquear quando depender de credencial ou serviço externo ainda não configurado. O próximo passo é transformar a pendência em issue, teste, decisão de reunião ou evidência de piloto. Este capítulo não promete superioridade nem receita: ele deixa um contrato verificável para que a equipe consiga descobrir se a hipótese se sustenta.

\section{Teste 05: Perfil de terceiro}
A discussão de teste 05: perfil de terceiro começou por uma situação concreta: objetivo: validar perfil de terceiro no ambiente de demonstração.

Nossa escolha para este ponto foi direta: Resultado esperado: Dados privados não são retornados em consulta pública.

Na implementação, este requisito é tratado como parte do sistema e não como uma tela isolada. A interface orienta a próxima ação; a API confere identidade, papel, estado e limites; a persistência guarda somente o necessário; e as integrações externas ficam condicionadas às variáveis do ambiente. O escopo continua exclusivamente de software, sem depender de escopo físico.

O fluxo começa quando a pessoa encontra uma necessidade ou uma habilidade que deseja compartilhar. Em seguida, ela confirma a intenção, executa a ação principal e recebe uma resposta que pode ser consultada depois. Para nós, teste 05: perfil de terceiro só está concluído quando o caminho é compreensível, existe recuperação para erro e o servidor registra uma evidência coerente.

O aceite não fica apenas no texto. A evidência prevista é: Procedimento: preparar dados sintéticos, executar a ação, verificar resposta HTTP ou evento, conferir banco/ledger quando aplicável e registrar versão do código. O critério de aceite deve ser escrito antes do piloto, com entrada, resultado esperado, responsável, data e ferramenta de verificação. Assim, a documentação separa conformidade técnica, observação de uso e hipótese de impacto.

A validação combina teste automatizado, inspeção manual e uma tarefa representativa. Vamos registrar tempo, erros, conclusão, compreensão e barreiras observadas por teclado, leitor de tela, zoom, contraste, texto, áudio ou conexão instável. O objetivo não é criar uma aparência de certeza, e sim manter um ciclo curto entre observar, corrigir e conferir novamente.

Também registramos uma limitação, porque ela faz parte do resultado: Classificação: aprovado, reprovado ou bloqueado. Não marcar como aprovado sem evidência reproduzível; bloquear quando depender de credencial ou serviço externo ainda não configurado. O próximo passo é transformar a pendência em issue, teste, decisão de reunião ou evidência de piloto. Este capítulo não promete superioridade nem receita: ele deixa um contrato verificável para que a equipe consiga descobrir se a hipótese se sustenta.

\section{Teste 06: Proposta de terceiro}
Durante a consolidação do MVP, percebemos que objetivo: validar proposta de terceiro no ambiente de demonstração.

No protótipo, esse entendimento aparece assim: Resultado esperado: Ação em proposta alheia é rejeitada pelo proprietário/regra.

Na implementação, este requisito é tratado como parte do sistema e não como uma tela isolada. A interface orienta a próxima ação; a API confere identidade, papel, estado e limites; a persistência guarda somente o necessário; e as integrações externas ficam condicionadas às variáveis do ambiente. O escopo continua exclusivamente de software, sem depender de escopo físico.

O fluxo começa quando a pessoa encontra uma necessidade ou uma habilidade que deseja compartilhar. Em seguida, ela confirma a intenção, executa a ação principal e recebe uma resposta que pode ser consultada depois. Para nós, teste 06: proposta de terceiro só está concluído quando o caminho é compreensível, existe recuperação para erro e o servidor registra uma evidência coerente.

A forma mais honesta de verificar este requisito é registrar: Procedimento: preparar dados sintéticos, executar a ação, verificar resposta HTTP ou evento, conferir banco/ledger quando aplicável e registrar versão do código. O critério de aceite deve ser escrito antes do piloto, com entrada, resultado esperado, responsável, data e ferramenta de verificação. Assim, a documentação separa conformidade técnica, observação de uso e hipótese de impacto.

A validação combina teste automatizado, inspeção manual e uma tarefa representativa. Vamos registrar tempo, erros, conclusão, compreensão e barreiras observadas por teclado, leitor de tela, zoom, contraste, texto, áudio ou conexão instável. O objetivo não é criar uma aparência de certeza, e sim manter um ciclo curto entre observar, corrigir e conferir novamente.

A equipe não trata este item como resolvido antes de enfrentar o risco: Classificação: aprovado, reprovado ou bloqueado. Não marcar como aprovado sem evidência reproduzível; bloquear quando depender de credencial ou serviço externo ainda não configurado. O próximo passo é transformar a pendência em issue, teste, decisão de reunião ou evidência de piloto. Este capítulo não promete superioridade nem receita: ele deixa um contrato verificável para que a equipe consiga descobrir se a hipótese se sustenta.

\section{Teste 07: WebSocket sem primeiro frame}
A discussão de teste 07: websocket sem primeiro frame começou por uma situação concreta: objetivo: validar websocket sem primeiro frame no ambiente de demonstração.

Nossa escolha para este ponto foi direta: Resultado esperado: Nenhuma mensagem de negócio é transmitida antes de autenticar.

Na implementação, este requisito é tratado como parte do sistema e não como uma tela isolada. A interface orienta a próxima ação; a API confere identidade, papel, estado e limites; a persistência guarda somente o necessário; e as integrações externas ficam condicionadas às variáveis do ambiente. O escopo continua exclusivamente de software, sem depender de escopo físico.

O fluxo começa quando a pessoa encontra uma necessidade ou uma habilidade que deseja compartilhar. Em seguida, ela confirma a intenção, executa a ação principal e recebe uma resposta que pode ser consultada depois. Para nós, teste 07: websocket sem primeiro frame só está concluído quando o caminho é compreensível, existe recuperação para erro e o servidor registra uma evidência coerente.

O aceite não fica apenas no texto. A evidência prevista é: Procedimento: preparar dados sintéticos, executar a ação, verificar resposta HTTP ou evento, conferir banco/ledger quando aplicável e registrar versão do código. O critério de aceite deve ser escrito antes do piloto, com entrada, resultado esperado, responsável, data e ferramenta de verificação. Assim, a documentação separa conformidade técnica, observação de uso e hipótese de impacto.

A validação combina teste automatizado, inspeção manual e uma tarefa representativa. Vamos registrar tempo, erros, conclusão, compreensão e barreiras observadas por teclado, leitor de tela, zoom, contraste, texto, áudio ou conexão instável. O objetivo não é criar uma aparência de certeza, e sim manter um ciclo curto entre observar, corrigir e conferir novamente.

Também registramos uma limitação, porque ela faz parte do resultado: Classificação: aprovado, reprovado ou bloqueado. Não marcar como aprovado sem evidência reproduzível; bloquear quando depender de credencial ou serviço externo ainda não configurado. O próximo passo é transformar a pendência em issue, teste, decisão de reunião ou evidência de piloto. Este capítulo não promete superioridade nem receita: ele deixa um contrato verificável para que a equipe consiga descobrir se a hipótese se sustenta.

\section{Teste 08: WebSocket sender falso}
A discussão de teste 08: websocket sender falso começou por uma situação concreta: objetivo: validar websocket sender falso no ambiente de demonstração.

Nossa escolha para este ponto foi direta: Resultado esperado: Payload com senderId divergente encerra a sessão.

Na implementação, este requisito é tratado como parte do sistema e não como uma tela isolada. A interface orienta a próxima ação; a API confere identidade, papel, estado e limites; a persistência guarda somente o necessário; e as integrações externas ficam condicionadas às variáveis do ambiente. O escopo continua exclusivamente de software, sem depender de escopo físico.

O fluxo começa quando a pessoa encontra uma necessidade ou uma habilidade que deseja compartilhar. Em seguida, ela confirma a intenção, executa a ação principal e recebe uma resposta que pode ser consultada depois. Para nós, teste 08: websocket sender falso só está concluído quando o caminho é compreensível, existe recuperação para erro e o servidor registra uma evidência coerente.

O aceite não fica apenas no texto. A evidência prevista é: Procedimento: preparar dados sintéticos, executar a ação, verificar resposta HTTP ou evento, conferir banco/ledger quando aplicável e registrar versão do código. O critério de aceite deve ser escrito antes do piloto, com entrada, resultado esperado, responsável, data e ferramenta de verificação. Assim, a documentação separa conformidade técnica, observação de uso e hipótese de impacto.

A validação combina teste automatizado, inspeção manual e uma tarefa representativa. Vamos registrar tempo, erros, conclusão, compreensão e barreiras observadas por teclado, leitor de tela, zoom, contraste, texto, áudio ou conexão instável. O objetivo não é criar uma aparência de certeza, e sim manter um ciclo curto entre observar, corrigir e conferir novamente.

Também registramos uma limitação, porque ela faz parte do resultado: Classificação: aprovado, reprovado ou bloqueado. Não marcar como aprovado sem evidência reproduzível; bloquear quando depender de credencial ou serviço externo ainda não configurado. O próximo passo é transformar a pendência em issue, teste, decisão de reunião ou evidência de piloto. Este capítulo não promete superioridade nem receita: ele deixa um contrato verificável para que a equipe consiga descobrir se a hipótese se sustenta.

\section{Teste 09: Wallet URL limpa}
Na leitura que fizemos do projeto, objetivo: validar wallet url limpa no ambiente de demonstração.

Depois de comparar as alternativas, a equipe decidiu o seguinte: Resultado esperado: Nenhum token aparece em query string ou rota.

Na implementação, este requisito é tratado como parte do sistema e não como uma tela isolada. A interface orienta a próxima ação; a API confere identidade, papel, estado e limites; a persistência guarda somente o necessário; e as integrações externas ficam condicionadas às variáveis do ambiente. O escopo continua exclusivamente de software, sem depender de escopo físico.

O fluxo começa quando a pessoa encontra uma necessidade ou uma habilidade que deseja compartilhar. Em seguida, ela confirma a intenção, executa a ação principal e recebe uma resposta que pode ser consultada depois. Para nós, teste 09: wallet url limpa só está concluído quando o caminho é compreensível, existe recuperação para erro e o servidor registra uma evidência coerente.

Para conferir se a decisão funciona, vamos procurar esta evidência: Procedimento: preparar dados sintéticos, executar a ação, verificar resposta HTTP ou evento, conferir banco/ledger quando aplicável e registrar versão do código. O critério de aceite deve ser escrito antes do piloto, com entrada, resultado esperado, responsável, data e ferramenta de verificação. Assim, a documentação separa conformidade técnica, observação de uso e hipótese de impacto.

A validação combina teste automatizado, inspeção manual e uma tarefa representativa. Vamos registrar tempo, erros, conclusão, compreensão e barreiras observadas por teclado, leitor de tela, zoom, contraste, texto, áudio ou conexão instável. O objetivo não é criar uma aparência de certeza, e sim manter um ciclo curto entre observar, corrigir e conferir novamente.

O ponto que ainda merece cuidado é o seguinte: Classificação: aprovado, reprovado ou bloqueado. Não marcar como aprovado sem evidência reproduzível; bloquear quando depender de credencial ou serviço externo ainda não configurado. O próximo passo é transformar a pendência em issue, teste, decisão de reunião ou evidência de piloto. Este capítulo não promete superioridade nem receita: ele deixa um contrato verificável para que a equipe consiga descobrir se a hipótese se sustenta.

\section{Teste 10: FCM token vazio}
Na leitura que fizemos do projeto, objetivo: validar fcm token vazio no ambiente de demonstração.

Depois de comparar as alternativas, a equipe decidiu o seguinte: Resultado esperado: Validação retorna 400 e não altera banco.

Na implementação, este requisito é tratado como parte do sistema e não como uma tela isolada. A interface orienta a próxima ação; a API confere identidade, papel, estado e limites; a persistência guarda somente o necessário; e as integrações externas ficam condicionadas às variáveis do ambiente. O escopo continua exclusivamente de software, sem depender de escopo físico.

O fluxo começa quando a pessoa encontra uma necessidade ou uma habilidade que deseja compartilhar. Em seguida, ela confirma a intenção, executa a ação principal e recebe uma resposta que pode ser consultada depois. Para nós, teste 10: fcm token vazio só está concluído quando o caminho é compreensível, existe recuperação para erro e o servidor registra uma evidência coerente.

Para conferir se a decisão funciona, vamos procurar esta evidência: Procedimento: preparar dados sintéticos, executar a ação, verificar resposta HTTP ou evento, conferir banco/ledger quando aplicável e registrar versão do código. O critério de aceite deve ser escrito antes do piloto, com entrada, resultado esperado, responsável, data e ferramenta de verificação. Assim, a documentação separa conformidade técnica, observação de uso e hipótese de impacto.

A validação combina teste automatizado, inspeção manual e uma tarefa representativa. Vamos registrar tempo, erros, conclusão, compreensão e barreiras observadas por teclado, leitor de tela, zoom, contraste, texto, áudio ou conexão instável. O objetivo não é criar uma aparência de certeza, e sim manter um ciclo curto entre observar, corrigir e conferir novamente.

O ponto que ainda merece cuidado é o seguinte: Classificação: aprovado, reprovado ou bloqueado. Não marcar como aprovado sem evidência reproduzível; bloquear quando depender de credencial ou serviço externo ainda não configurado. O próximo passo é transformar a pendência em issue, teste, decisão de reunião ou evidência de piloto. Este capítulo não promete superioridade nem receita: ele deixa um contrato verificável para que a equipe consiga descobrir se a hipótese se sustenta.

\section{Teste 11: FCM revogado}
Durante a consolidação do MVP, percebemos que objetivo: validar fcm revogado no ambiente de demonstração.

No protótipo, esse entendimento aparece assim: Resultado esperado: DELETE remove o token do usuário autenticado.

Na implementação, este requisito é tratado como parte do sistema e não como uma tela isolada. A interface orienta a próxima ação; a API confere identidade, papel, estado e limites; a persistência guarda somente o necessário; e as integrações externas ficam condicionadas às variáveis do ambiente. O escopo continua exclusivamente de software, sem depender de escopo físico.

O fluxo começa quando a pessoa encontra uma necessidade ou uma habilidade que deseja compartilhar. Em seguida, ela confirma a intenção, executa a ação principal e recebe uma resposta que pode ser consultada depois. Para nós, teste 11: fcm revogado só está concluído quando o caminho é compreensível, existe recuperação para erro e o servidor registra uma evidência coerente.

A forma mais honesta de verificar este requisito é registrar: Procedimento: preparar dados sintéticos, executar a ação, verificar resposta HTTP ou evento, conferir banco/ledger quando aplicável e registrar versão do código. O critério de aceite deve ser escrito antes do piloto, com entrada, resultado esperado, responsável, data e ferramenta de verificação. Assim, a documentação separa conformidade técnica, observação de uso e hipótese de impacto.

A validação combina teste automatizado, inspeção manual e uma tarefa representativa. Vamos registrar tempo, erros, conclusão, compreensão e barreiras observadas por teclado, leitor de tela, zoom, contraste, texto, áudio ou conexão instável. O objetivo não é criar uma aparência de certeza, e sim manter um ciclo curto entre observar, corrigir e conferir novamente.

A equipe não trata este item como resolvido antes de enfrentar o risco: Classificação: aprovado, reprovado ou bloqueado. Não marcar como aprovado sem evidência reproduzível; bloquear quando depender de credencial ou serviço externo ainda não configurado. O próximo passo é transformar a pendência em issue, teste, decisão de reunião ou evidência de piloto. Este capítulo não promete superioridade nem receita: ele deixa um contrato verificável para que a equipe consiga descobrir se a hipótese se sustenta.

\section{Teste 12: Webhook sem assinatura}
A discussão de teste 12: webhook sem assinatura começou por uma situação concreta: objetivo: validar webhook sem assinatura no ambiente de demonstração.

Nossa escolha para este ponto foi direta: Resultado esperado: recursos comerciais permanecem bloqueados.

Na implementação, este requisito é tratado como parte do sistema e não como uma tela isolada. A interface orienta a próxima ação; a API confere identidade, papel, estado e limites; a persistência guarda somente o necessário; e as integrações externas ficam condicionadas às variáveis do ambiente. O escopo continua exclusivamente de software, sem depender de escopo físico.

O fluxo começa quando a pessoa encontra uma necessidade ou uma habilidade que deseja compartilhar. Em seguida, ela confirma a intenção, executa a ação principal e recebe uma resposta que pode ser consultada depois. Para nós, teste 12: webhook sem assinatura só está concluído quando o caminho é compreensível, existe recuperação para erro e o servidor registra uma evidência coerente.

O aceite não fica apenas no texto. A evidência prevista é: Procedimento: preparar dados sintéticos, executar a ação, verificar resposta HTTP ou evento, conferir banco/ledger quando aplicável e registrar versão do código. O critério de aceite deve ser escrito antes do piloto, com entrada, resultado esperado, responsável, data e ferramenta de verificação. Assim, a documentação separa conformidade técnica, observação de uso e hipótese de impacto.

A validação combina teste automatizado, inspeção manual e uma tarefa representativa. Vamos registrar tempo, erros, conclusão, compreensão e barreiras observadas por teclado, leitor de tela, zoom, contraste, texto, áudio ou conexão instável. O objetivo não é criar uma aparência de certeza, e sim manter um ciclo curto entre observar, corrigir e conferir novamente.

Também registramos uma limitação, porque ela faz parte do resultado: Classificação: aprovado, reprovado ou bloqueado. Não marcar como aprovado sem evidência reproduzível; bloquear quando depender de credencial ou serviço externo ainda não configurado. O próximo passo é transformar a pendência em issue, teste, decisão de reunião ou evidência de piloto. Este capítulo não promete superioridade nem receita: ele deixa um contrato verificável para que a equipe consiga descobrir se a hipótese se sustenta.

\section{Teste 13: Webhook repetido}
A discussão de teste 13: webhook repetido começou por uma situação concreta: objetivo: validar webhook repetido no ambiente de demonstração.

Nossa escolha para este ponto foi direta: Resultado esperado: Segunda entrega é no-op idempotente.

Na implementação, este requisito é tratado como parte do sistema e não como uma tela isolada. A interface orienta a próxima ação; a API confere identidade, papel, estado e limites; a persistência guarda somente o necessário; e as integrações externas ficam condicionadas às variáveis do ambiente. O escopo continua exclusivamente de software, sem depender de escopo físico.

O fluxo começa quando a pessoa encontra uma necessidade ou uma habilidade que deseja compartilhar. Em seguida, ela confirma a intenção, executa a ação principal e recebe uma resposta que pode ser consultada depois. Para nós, teste 13: webhook repetido só está concluído quando o caminho é compreensível, existe recuperação para erro e o servidor registra uma evidência coerente.

O aceite não fica apenas no texto. A evidência prevista é: Procedimento: preparar dados sintéticos, executar a ação, verificar resposta HTTP ou evento, conferir banco/ledger quando aplicável e registrar versão do código. O critério de aceite deve ser escrito antes do piloto, com entrada, resultado esperado, responsável, data e ferramenta de verificação. Assim, a documentação separa conformidade técnica, observação de uso e hipótese de impacto.

A validação combina teste automatizado, inspeção manual e uma tarefa representativa. Vamos registrar tempo, erros, conclusão, compreensão e barreiras observadas por teclado, leitor de tela, zoom, contraste, texto, áudio ou conexão instável. O objetivo não é criar uma aparência de certeza, e sim manter um ciclo curto entre observar, corrigir e conferir novamente.

Também registramos uma limitação, porque ela faz parte do resultado: Classificação: aprovado, reprovado ou bloqueado. Não marcar como aprovado sem evidência reproduzível; bloquear quando depender de credencial ou serviço externo ainda não configurado. O próximo passo é transformar a pendência em issue, teste, decisão de reunião ou evidência de piloto. Este capítulo não promete superioridade nem receita: ele deixa um contrato verificável para que a equipe consiga descobrir se a hipótese se sustenta.

\section{Teste 14: Webhook valor incorreto}
A discussão de teste 14: webhook valor incorreto começou por uma situação concreta: objetivo: validar webhook valor incorreto no ambiente de demonstração.

Nossa escolha para este ponto foi direta: Resultado esperado: Evento é rejeitado e auditado.

Na implementação, este requisito é tratado como parte do sistema e não como uma tela isolada. A interface orienta a próxima ação; a API confere identidade, papel, estado e limites; a persistência guarda somente o necessário; e as integrações externas ficam condicionadas às variáveis do ambiente. O escopo continua exclusivamente de software, sem depender de escopo físico.

O fluxo começa quando a pessoa encontra uma necessidade ou uma habilidade que deseja compartilhar. Em seguida, ela confirma a intenção, executa a ação principal e recebe uma resposta que pode ser consultada depois. Para nós, teste 14: webhook valor incorreto só está concluído quando o caminho é compreensível, existe recuperação para erro e o servidor registra uma evidência coerente.

O aceite não fica apenas no texto. A evidência prevista é: Procedimento: preparar dados sintéticos, executar a ação, verificar resposta HTTP ou evento, conferir banco/ledger quando aplicável e registrar versão do código. O critério de aceite deve ser escrito antes do piloto, com entrada, resultado esperado, responsável, data e ferramenta de verificação. Assim, a documentação separa conformidade técnica, observação de uso e hipótese de impacto.

A validação combina teste automatizado, inspeção manual e uma tarefa representativa. Vamos registrar tempo, erros, conclusão, compreensão e barreiras observadas por teclado, leitor de tela, zoom, contraste, texto, áudio ou conexão instável. O objetivo não é criar uma aparência de certeza, e sim manter um ciclo curto entre observar, corrigir e conferir novamente.

Também registramos uma limitação, porque ela faz parte do resultado: Classificação: aprovado, reprovado ou bloqueado. Não marcar como aprovado sem evidência reproduzível; bloquear quando depender de credencial ou serviço externo ainda não configurado. O próximo passo é transformar a pendência em issue, teste, decisão de reunião ou evidência de piloto. Este capítulo não promete superioridade nem receita: ele deixa um contrato verificável para que a equipe consiga descobrir se a hipótese se sustenta.

\section{Teste 15: Compra paga}
A discussão de teste 15: compra paga começou por uma situação concreta: objetivo: validar compra paga no ambiente de demonstração.

Nossa escolha para este ponto foi direta: Resultado esperado: Ledger recebe crédito apenas uma vez.

Na implementação, este requisito é tratado como parte do sistema e não como uma tela isolada. A interface orienta a próxima ação; a API confere identidade, papel, estado e limites; a persistência guarda somente o necessário; e as integrações externas ficam condicionadas às variáveis do ambiente. O escopo continua exclusivamente de software, sem depender de escopo físico.

O fluxo começa quando a pessoa encontra uma necessidade ou uma habilidade que deseja compartilhar. Em seguida, ela confirma a intenção, executa a ação principal e recebe uma resposta que pode ser consultada depois. Para nós, teste 15: compra paga só está concluído quando o caminho é compreensível, existe recuperação para erro e o servidor registra uma evidência coerente.

O aceite não fica apenas no texto. A evidência prevista é: Procedimento: preparar dados sintéticos, executar a ação, verificar resposta HTTP ou evento, conferir banco/ledger quando aplicável e registrar versão do código. O critério de aceite deve ser escrito antes do piloto, com entrada, resultado esperado, responsável, data e ferramenta de verificação. Assim, a documentação separa conformidade técnica, observação de uso e hipótese de impacto.

A validação combina teste automatizado, inspeção manual e uma tarefa representativa. Vamos registrar tempo, erros, conclusão, compreensão e barreiras observadas por teclado, leitor de tela, zoom, contraste, texto, áudio ou conexão instável. O objetivo não é criar uma aparência de certeza, e sim manter um ciclo curto entre observar, corrigir e conferir novamente.

Também registramos uma limitação, porque ela faz parte do resultado: Classificação: aprovado, reprovado ou bloqueado. Não marcar como aprovado sem evidência reproduzível; bloquear quando depender de credencial ou serviço externo ainda não configurado. O próximo passo é transformar a pendência em issue, teste, decisão de reunião ou evidência de piloto. Este capítulo não promete superioridade nem receita: ele deixa um contrato verificável para que a equipe consiga descobrir se a hipótese se sustenta.

\section{Teste 16: Saldo negativo}
A discussão de teste 16: saldo negativo começou por uma situação concreta: objetivo: validar saldo negativo no ambiente de demonstração.

Nossa escolha para este ponto foi direta: Resultado esperado: Constraint e serviço rejeitam mutação.

Na implementação, este requisito é tratado como parte do sistema e não como uma tela isolada. A interface orienta a próxima ação; a API confere identidade, papel, estado e limites; a persistência guarda somente o necessário; e as integrações externas ficam condicionadas às variáveis do ambiente. O escopo continua exclusivamente de software, sem depender de escopo físico.

O fluxo começa quando a pessoa encontra uma necessidade ou uma habilidade que deseja compartilhar. Em seguida, ela confirma a intenção, executa a ação principal e recebe uma resposta que pode ser consultada depois. Para nós, teste 16: saldo negativo só está concluído quando o caminho é compreensível, existe recuperação para erro e o servidor registra uma evidência coerente.

O aceite não fica apenas no texto. A evidência prevista é: Procedimento: preparar dados sintéticos, executar a ação, verificar resposta HTTP ou evento, conferir banco/ledger quando aplicável e registrar versão do código. O critério de aceite deve ser escrito antes do piloto, com entrada, resultado esperado, responsável, data e ferramenta de verificação. Assim, a documentação separa conformidade técnica, observação de uso e hipótese de impacto.

A validação combina teste automatizado, inspeção manual e uma tarefa representativa. Vamos registrar tempo, erros, conclusão, compreensão e barreiras observadas por teclado, leitor de tela, zoom, contraste, texto, áudio ou conexão instável. O objetivo não é criar uma aparência de certeza, e sim manter um ciclo curto entre observar, corrigir e conferir novamente.

Também registramos uma limitação, porque ela faz parte do resultado: Classificação: aprovado, reprovado ou bloqueado. Não marcar como aprovado sem evidência reproduzível; bloquear quando depender de credencial ou serviço externo ainda não configurado. O próximo passo é transformar a pendência em issue, teste, decisão de reunião ou evidência de piloto. Este capítulo não promete superioridade nem receita: ele deixa um contrato verificável para que a equipe consiga descobrir se a hipótese se sustenta.

\section{Teste 17: Transferência concorrente}
Durante a consolidação do MVP, percebemos que objetivo: validar transferência concorrente no ambiente de demonstração.

No protótipo, esse entendimento aparece assim: Resultado esperado: Transação preserva consistência.

Na implementação, este requisito é tratado como parte do sistema e não como uma tela isolada. A interface orienta a próxima ação; a API confere identidade, papel, estado e limites; a persistência guarda somente o necessário; e as integrações externas ficam condicionadas às variáveis do ambiente. O escopo continua exclusivamente de software, sem depender de escopo físico.

O fluxo começa quando a pessoa encontra uma necessidade ou uma habilidade que deseja compartilhar. Em seguida, ela confirma a intenção, executa a ação principal e recebe uma resposta que pode ser consultada depois. Para nós, teste 17: transferência concorrente só está concluído quando o caminho é compreensível, existe recuperação para erro e o servidor registra uma evidência coerente.

A forma mais honesta de verificar este requisito é registrar: Procedimento: preparar dados sintéticos, executar a ação, verificar resposta HTTP ou evento, conferir banco/ledger quando aplicável e registrar versão do código. O critério de aceite deve ser escrito antes do piloto, com entrada, resultado esperado, responsável, data e ferramenta de verificação. Assim, a documentação separa conformidade técnica, observação de uso e hipótese de impacto.

A validação combina teste automatizado, inspeção manual e uma tarefa representativa. Vamos registrar tempo, erros, conclusão, compreensão e barreiras observadas por teclado, leitor de tela, zoom, contraste, texto, áudio ou conexão instável. O objetivo não é criar uma aparência de certeza, e sim manter um ciclo curto entre observar, corrigir e conferir novamente.

A equipe não trata este item como resolvido antes de enfrentar o risco: Classificação: aprovado, reprovado ou bloqueado. Não marcar como aprovado sem evidência reproduzível; bloquear quando depender de credencial ou serviço externo ainda não configurado. O próximo passo é transformar a pendência em issue, teste, decisão de reunião ou evidência de piloto. Este capítulo não promete superioridade nem receita: ele deixa um contrato verificável para que a equipe consiga descobrir se a hipótese se sustenta.

\section{Teste 18: Recompensa repetida}
A discussão de teste 18: recompensa repetida começou por uma situação concreta: objetivo: validar recompensa repetida no ambiente de demonstração.

Nossa escolha para este ponto foi direta: Resultado esperado: Chave de referência impede duplicidade.

Na implementação, este requisito é tratado como parte do sistema e não como uma tela isolada. A interface orienta a próxima ação; a API confere identidade, papel, estado e limites; a persistência guarda somente o necessário; e as integrações externas ficam condicionadas às variáveis do ambiente. O escopo continua exclusivamente de software, sem depender de escopo físico.

O fluxo começa quando a pessoa encontra uma necessidade ou uma habilidade que deseja compartilhar. Em seguida, ela confirma a intenção, executa a ação principal e recebe uma resposta que pode ser consultada depois. Para nós, teste 18: recompensa repetida só está concluído quando o caminho é compreensível, existe recuperação para erro e o servidor registra uma evidência coerente.

O aceite não fica apenas no texto. A evidência prevista é: Procedimento: preparar dados sintéticos, executar a ação, verificar resposta HTTP ou evento, conferir banco/ledger quando aplicável e registrar versão do código. O critério de aceite deve ser escrito antes do piloto, com entrada, resultado esperado, responsável, data e ferramenta de verificação. Assim, a documentação separa conformidade técnica, observação de uso e hipótese de impacto.

A validação combina teste automatizado, inspeção manual e uma tarefa representativa. Vamos registrar tempo, erros, conclusão, compreensão e barreiras observadas por teclado, leitor de tela, zoom, contraste, texto, áudio ou conexão instável. O objetivo não é criar uma aparência de certeza, e sim manter um ciclo curto entre observar, corrigir e conferir novamente.

Também registramos uma limitação, porque ela faz parte do resultado: Classificação: aprovado, reprovado ou bloqueado. Não marcar como aprovado sem evidência reproduzível; bloquear quando depender de credencial ou serviço externo ainda não configurado. O próximo passo é transformar a pendência em issue, teste, decisão de reunião ou evidência de piloto. Este capítulo não promete superioridade nem receita: ele deixa um contrato verificável para que a equipe consiga descobrir se a hipótese se sustenta.

\section{Teste 19: Cota de IA}
Na leitura que fizemos do projeto, objetivo: validar cota de ia no ambiente de demonstração.

Depois de comparar as alternativas, a equipe decidiu o seguinte: Resultado esperado: Limite é aplicado no servidor e não no contador do cliente.

Na implementação, este requisito é tratado como parte do sistema e não como uma tela isolada. A interface orienta a próxima ação; a API confere identidade, papel, estado e limites; a persistência guarda somente o necessário; e as integrações externas ficam condicionadas às variáveis do ambiente. O escopo continua exclusivamente de software, sem depender de escopo físico.

O fluxo começa quando a pessoa encontra uma necessidade ou uma habilidade que deseja compartilhar. Em seguida, ela confirma a intenção, executa a ação principal e recebe uma resposta que pode ser consultada depois. Para nós, teste 19: cota de ia só está concluído quando o caminho é compreensível, existe recuperação para erro e o servidor registra uma evidência coerente.

Para conferir se a decisão funciona, vamos procurar esta evidência: Procedimento: preparar dados sintéticos, executar a ação, verificar resposta HTTP ou evento, conferir banco/ledger quando aplicável e registrar versão do código. O critério de aceite deve ser escrito antes do piloto, com entrada, resultado esperado, responsável, data e ferramenta de verificação. Assim, a documentação separa conformidade técnica, observação de uso e hipótese de impacto.

A validação combina teste automatizado, inspeção manual e uma tarefa representativa. Vamos registrar tempo, erros, conclusão, compreensão e barreiras observadas por teclado, leitor de tela, zoom, contraste, texto, áudio ou conexão instável. O objetivo não é criar uma aparência de certeza, e sim manter um ciclo curto entre observar, corrigir e conferir novamente.

O ponto que ainda merece cuidado é o seguinte: Classificação: aprovado, reprovado ou bloqueado. Não marcar como aprovado sem evidência reproduzível; bloquear quando depender de credencial ou serviço externo ainda não configurado. O próximo passo é transformar a pendência em issue, teste, decisão de reunião ou evidência de piloto. Este capítulo não promete superioridade nem receita: ele deixa um contrato verificável para que a equipe consiga descobrir se a hipótese se sustenta.

\section{Teste 20: Falha FCM}
A discussão de teste 20: falha fcm começou por uma situação concreta: objetivo: validar falha fcm no ambiente de demonstração.

Nossa escolha para este ponto foi direta: Resultado esperado: Notificação persistida não é perdida se push falhar.

Na implementação, este requisito é tratado como parte do sistema e não como uma tela isolada. A interface orienta a próxima ação; a API confere identidade, papel, estado e limites; a persistência guarda somente o necessário; e as integrações externas ficam condicionadas às variáveis do ambiente. O escopo continua exclusivamente de software, sem depender de escopo físico.

O fluxo começa quando a pessoa encontra uma necessidade ou uma habilidade que deseja compartilhar. Em seguida, ela confirma a intenção, executa a ação principal e recebe uma resposta que pode ser consultada depois. Para nós, teste 20: falha fcm só está concluído quando o caminho é compreensível, existe recuperação para erro e o servidor registra uma evidência coerente.

O aceite não fica apenas no texto. A evidência prevista é: Procedimento: preparar dados sintéticos, executar a ação, verificar resposta HTTP ou evento, conferir banco/ledger quando aplicável e registrar versão do código. O critério de aceite deve ser escrito antes do piloto, com entrada, resultado esperado, responsável, data e ferramenta de verificação. Assim, a documentação separa conformidade técnica, observação de uso e hipótese de impacto.

A validação combina teste automatizado, inspeção manual e uma tarefa representativa. Vamos registrar tempo, erros, conclusão, compreensão e barreiras observadas por teclado, leitor de tela, zoom, contraste, texto, áudio ou conexão instável. O objetivo não é criar uma aparência de certeza, e sim manter um ciclo curto entre observar, corrigir e conferir novamente.

Também registramos uma limitação, porque ela faz parte do resultado: Classificação: aprovado, reprovado ou bloqueado. Não marcar como aprovado sem evidência reproduzível; bloquear quando depender de credencial ou serviço externo ainda não configurado. O próximo passo é transformar a pendência em issue, teste, decisão de reunião ou evidência de piloto. Este capítulo não promete superioridade nem receita: ele deixa um contrato verificável para que a equipe consiga descobrir se a hipótese se sustenta.

\section{Teste 21: Banco vazio}
Na leitura que fizemos do projeto, objetivo: validar banco vazio no ambiente de demonstração.

Depois de comparar as alternativas, a equipe decidiu o seguinte: Resultado esperado: Flyway cria schema em ordem e Hibernate valida.

Na implementação, este requisito é tratado como parte do sistema e não como uma tela isolada. A interface orienta a próxima ação; a API confere identidade, papel, estado e limites; a persistência guarda somente o necessário; e as integrações externas ficam condicionadas às variáveis do ambiente. O escopo continua exclusivamente de software, sem depender de escopo físico.

O fluxo começa quando a pessoa encontra uma necessidade ou uma habilidade que deseja compartilhar. Em seguida, ela confirma a intenção, executa a ação principal e recebe uma resposta que pode ser consultada depois. Para nós, teste 21: banco vazio só está concluído quando o caminho é compreensível, existe recuperação para erro e o servidor registra uma evidência coerente.

Para conferir se a decisão funciona, vamos procurar esta evidência: Procedimento: preparar dados sintéticos, executar a ação, verificar resposta HTTP ou evento, conferir banco/ledger quando aplicável e registrar versão do código. O critério de aceite deve ser escrito antes do piloto, com entrada, resultado esperado, responsável, data e ferramenta de verificação. Assim, a documentação separa conformidade técnica, observação de uso e hipótese de impacto.

A validação combina teste automatizado, inspeção manual e uma tarefa representativa. Vamos registrar tempo, erros, conclusão, compreensão e barreiras observadas por teclado, leitor de tela, zoom, contraste, texto, áudio ou conexão instável. O objetivo não é criar uma aparência de certeza, e sim manter um ciclo curto entre observar, corrigir e conferir novamente.

O ponto que ainda merece cuidado é o seguinte: Classificação: aprovado, reprovado ou bloqueado. Não marcar como aprovado sem evidência reproduzível; bloquear quando depender de credencial ou serviço externo ainda não configurado. O próximo passo é transformar a pendência em issue, teste, decisão de reunião ou evidência de piloto. Este capítulo não promete superioridade nem receita: ele deixa um contrato verificável para que a equipe consiga descobrir se a hipótese se sustenta.

\section{Teste 22: Migration repetida}
Durante a consolidação do MVP, percebemos que objetivo: validar migration repetida no ambiente de demonstração.

No protótipo, esse entendimento aparece assim: Resultado esperado: Flyway informa estado consistente sem recriar dados.

Na implementação, este requisito é tratado como parte do sistema e não como uma tela isolada. A interface orienta a próxima ação; a API confere identidade, papel, estado e limites; a persistência guarda somente o necessário; e as integrações externas ficam condicionadas às variáveis do ambiente. O escopo continua exclusivamente de software, sem depender de escopo físico.

O fluxo começa quando a pessoa encontra uma necessidade ou uma habilidade que deseja compartilhar. Em seguida, ela confirma a intenção, executa a ação principal e recebe uma resposta que pode ser consultada depois. Para nós, teste 22: migration repetida só está concluído quando o caminho é compreensível, existe recuperação para erro e o servidor registra uma evidência coerente.

A forma mais honesta de verificar este requisito é registrar: Procedimento: preparar dados sintéticos, executar a ação, verificar resposta HTTP ou evento, conferir banco/ledger quando aplicável e registrar versão do código. O critério de aceite deve ser escrito antes do piloto, com entrada, resultado esperado, responsável, data e ferramenta de verificação. Assim, a documentação separa conformidade técnica, observação de uso e hipótese de impacto.

A validação combina teste automatizado, inspeção manual e uma tarefa representativa. Vamos registrar tempo, erros, conclusão, compreensão e barreiras observadas por teclado, leitor de tela, zoom, contraste, texto, áudio ou conexão instável. O objetivo não é criar uma aparência de certeza, e sim manter um ciclo curto entre observar, corrigir e conferir novamente.

A equipe não trata este item como resolvido antes de enfrentar o risco: Classificação: aprovado, reprovado ou bloqueado. Não marcar como aprovado sem evidência reproduzível; bloquear quando depender de credencial ou serviço externo ainda não configurado. O próximo passo é transformar a pendência em issue, teste, decisão de reunião ou evidência de piloto. Este capítulo não promete superioridade nem receita: ele deixa um contrato verificável para que a equipe consiga descobrir se a hipótese se sustenta.

\section{Teste 23: Backup}
Na leitura que fizemos do projeto, objetivo: validar backup no ambiente de demonstração.

Depois de comparar as alternativas, a equipe decidiu o seguinte: Resultado esperado: Restore produz banco operacional em ambiente isolado.

Na implementação, este requisito é tratado como parte do sistema e não como uma tela isolada. A interface orienta a próxima ação; a API confere identidade, papel, estado e limites; a persistência guarda somente o necessário; e as integrações externas ficam condicionadas às variáveis do ambiente. O escopo continua exclusivamente de software, sem depender de escopo físico.

O fluxo começa quando a pessoa encontra uma necessidade ou uma habilidade que deseja compartilhar. Em seguida, ela confirma a intenção, executa a ação principal e recebe uma resposta que pode ser consultada depois. Para nós, teste 23: backup só está concluído quando o caminho é compreensível, existe recuperação para erro e o servidor registra uma evidência coerente.

Para conferir se a decisão funciona, vamos procurar esta evidência: Procedimento: preparar dados sintéticos, executar a ação, verificar resposta HTTP ou evento, conferir banco/ledger quando aplicável e registrar versão do código. O critério de aceite deve ser escrito antes do piloto, com entrada, resultado esperado, responsável, data e ferramenta de verificação. Assim, a documentação separa conformidade técnica, observação de uso e hipótese de impacto.

A validação combina teste automatizado, inspeção manual e uma tarefa representativa. Vamos registrar tempo, erros, conclusão, compreensão e barreiras observadas por teclado, leitor de tela, zoom, contraste, texto, áudio ou conexão instável. O objetivo não é criar uma aparência de certeza, e sim manter um ciclo curto entre observar, corrigir e conferir novamente.

O ponto que ainda merece cuidado é o seguinte: Classificação: aprovado, reprovado ou bloqueado. Não marcar como aprovado sem evidência reproduzível; bloquear quando depender de credencial ou serviço externo ainda não configurado. O próximo passo é transformar a pendência em issue, teste, decisão de reunião ou evidência de piloto. Este capítulo não promete superioridade nem receita: ele deixa um contrato verificável para que a equipe consiga descobrir se a hipótese se sustenta.

\section{Teste 24: Rollback}
Na leitura que fizemos do projeto, objetivo: validar rollback no ambiente de demonstração.

Depois de comparar as alternativas, a equipe decidiu o seguinte: Resultado esperado: Imagem anterior sobe com contrato compatível.

Na implementação, este requisito é tratado como parte do sistema e não como uma tela isolada. A interface orienta a próxima ação; a API confere identidade, papel, estado e limites; a persistência guarda somente o necessário; e as integrações externas ficam condicionadas às variáveis do ambiente. O escopo continua exclusivamente de software, sem depender de escopo físico.

O fluxo começa quando a pessoa encontra uma necessidade ou uma habilidade que deseja compartilhar. Em seguida, ela confirma a intenção, executa a ação principal e recebe uma resposta que pode ser consultada depois. Para nós, teste 24: rollback só está concluído quando o caminho é compreensível, existe recuperação para erro e o servidor registra uma evidência coerente.

Para conferir se a decisão funciona, vamos procurar esta evidência: Procedimento: preparar dados sintéticos, executar a ação, verificar resposta HTTP ou evento, conferir banco/ledger quando aplicável e registrar versão do código. O critério de aceite deve ser escrito antes do piloto, com entrada, resultado esperado, responsável, data e ferramenta de verificação. Assim, a documentação separa conformidade técnica, observação de uso e hipótese de impacto.

A validação combina teste automatizado, inspeção manual e uma tarefa representativa. Vamos registrar tempo, erros, conclusão, compreensão e barreiras observadas por teclado, leitor de tela, zoom, contraste, texto, áudio ou conexão instável. O objetivo não é criar uma aparência de certeza, e sim manter um ciclo curto entre observar, corrigir e conferir novamente.

O ponto que ainda merece cuidado é o seguinte: Classificação: aprovado, reprovado ou bloqueado. Não marcar como aprovado sem evidência reproduzível; bloquear quando depender de credencial ou serviço externo ainda não configurado. O próximo passo é transformar a pendência em issue, teste, decisão de reunião ou evidência de piloto. Este capítulo não promete superioridade nem receita: ele deixa um contrato verificável para que a equipe consiga descobrir se a hipótese se sustenta.

\section{Teste 25: CORS inválido}
Na leitura que fizemos do projeto, objetivo: validar cors inválido no ambiente de demonstração.

Depois de comparar as alternativas, a equipe decidiu o seguinte: Resultado esperado: Origem não autorizada não recebe resposta autenticada.

Na implementação, este requisito é tratado como parte do sistema e não como uma tela isolada. A interface orienta a próxima ação; a API confere identidade, papel, estado e limites; a persistência guarda somente o necessário; e as integrações externas ficam condicionadas às variáveis do ambiente. O escopo continua exclusivamente de software, sem depender de escopo físico.

O fluxo começa quando a pessoa encontra uma necessidade ou uma habilidade que deseja compartilhar. Em seguida, ela confirma a intenção, executa a ação principal e recebe uma resposta que pode ser consultada depois. Para nós, teste 25: cors inválido só está concluído quando o caminho é compreensível, existe recuperação para erro e o servidor registra uma evidência coerente.

Para conferir se a decisão funciona, vamos procurar esta evidência: Procedimento: preparar dados sintéticos, executar a ação, verificar resposta HTTP ou evento, conferir banco/ledger quando aplicável e registrar versão do código. O critério de aceite deve ser escrito antes do piloto, com entrada, resultado esperado, responsável, data e ferramenta de verificação. Assim, a documentação separa conformidade técnica, observação de uso e hipótese de impacto.

A validação combina teste automatizado, inspeção manual e uma tarefa representativa. Vamos registrar tempo, erros, conclusão, compreensão e barreiras observadas por teclado, leitor de tela, zoom, contraste, texto, áudio ou conexão instável. O objetivo não é criar uma aparência de certeza, e sim manter um ciclo curto entre observar, corrigir e conferir novamente.

O ponto que ainda merece cuidado é o seguinte: Classificação: aprovado, reprovado ou bloqueado. Não marcar como aprovado sem evidência reproduzível; bloquear quando depender de credencial ou serviço externo ainda não configurado. O próximo passo é transformar a pendência em issue, teste, decisão de reunião ou evidência de piloto. Este capítulo não promete superioridade nem receita: ele deixa um contrato verificável para que a equipe consiga descobrir se a hipótese se sustenta.

\section{Teste 26: Upload grande}
Na leitura que fizemos do projeto, objetivo: validar upload grande no ambiente de demonstração.

Depois de comparar as alternativas, a equipe decidiu o seguinte: Resultado esperado: Limite de multipart evita abuso e retorna erro legível.

Na implementação, este requisito é tratado como parte do sistema e não como uma tela isolada. A interface orienta a próxima ação; a API confere identidade, papel, estado e limites; a persistência guarda somente o necessário; e as integrações externas ficam condicionadas às variáveis do ambiente. O escopo continua exclusivamente de software, sem depender de escopo físico.

O fluxo começa quando a pessoa encontra uma necessidade ou uma habilidade que deseja compartilhar. Em seguida, ela confirma a intenção, executa a ação principal e recebe uma resposta que pode ser consultada depois. Para nós, teste 26: upload grande só está concluído quando o caminho é compreensível, existe recuperação para erro e o servidor registra uma evidência coerente.

Para conferir se a decisão funciona, vamos procurar esta evidência: Procedimento: preparar dados sintéticos, executar a ação, verificar resposta HTTP ou evento, conferir banco/ledger quando aplicável e registrar versão do código. O critério de aceite deve ser escrito antes do piloto, com entrada, resultado esperado, responsável, data e ferramenta de verificação. Assim, a documentação separa conformidade técnica, observação de uso e hipótese de impacto.

A validação combina teste automatizado, inspeção manual e uma tarefa representativa. Vamos registrar tempo, erros, conclusão, compreensão e barreiras observadas por teclado, leitor de tela, zoom, contraste, texto, áudio ou conexão instável. O objetivo não é criar uma aparência de certeza, e sim manter um ciclo curto entre observar, corrigir e conferir novamente.

O ponto que ainda merece cuidado é o seguinte: Classificação: aprovado, reprovado ou bloqueado. Não marcar como aprovado sem evidência reproduzível; bloquear quando depender de credencial ou serviço externo ainda não configurado. O próximo passo é transformar a pendência em issue, teste, decisão de reunião ou evidência de piloto. Este capítulo não promete superioridade nem receita: ele deixa um contrato verificável para que a equipe consiga descobrir se a hipótese se sustenta.

\section{Teste 27: XSS em post}
A discussão de teste 27: xss em post começou por uma situação concreta: objetivo: validar xss em post no ambiente de demonstração.

Nossa escolha para este ponto foi direta: Resultado esperado: Conteúdo é renderizado como texto e não executa script.

Na implementação, este requisito é tratado como parte do sistema e não como uma tela isolada. A interface orienta a próxima ação; a API confere identidade, papel, estado e limites; a persistência guarda somente o necessário; e as integrações externas ficam condicionadas às variáveis do ambiente. O escopo continua exclusivamente de software, sem depender de escopo físico.

O fluxo começa quando a pessoa encontra uma necessidade ou uma habilidade que deseja compartilhar. Em seguida, ela confirma a intenção, executa a ação principal e recebe uma resposta que pode ser consultada depois. Para nós, teste 27: xss em post só está concluído quando o caminho é compreensível, existe recuperação para erro e o servidor registra uma evidência coerente.

O aceite não fica apenas no texto. A evidência prevista é: Procedimento: preparar dados sintéticos, executar a ação, verificar resposta HTTP ou evento, conferir banco/ledger quando aplicável e registrar versão do código. O critério de aceite deve ser escrito antes do piloto, com entrada, resultado esperado, responsável, data e ferramenta de verificação. Assim, a documentação separa conformidade técnica, observação de uso e hipótese de impacto.

A validação combina teste automatizado, inspeção manual e uma tarefa representativa. Vamos registrar tempo, erros, conclusão, compreensão e barreiras observadas por teclado, leitor de tela, zoom, contraste, texto, áudio ou conexão instável. O objetivo não é criar uma aparência de certeza, e sim manter um ciclo curto entre observar, corrigir e conferir novamente.

Também registramos uma limitação, porque ela faz parte do resultado: Classificação: aprovado, reprovado ou bloqueado. Não marcar como aprovado sem evidência reproduzível; bloquear quando depender de credencial ou serviço externo ainda não configurado. O próximo passo é transformar a pendência em issue, teste, decisão de reunião ou evidência de piloto. Este capítulo não promete superioridade nem receita: ele deixa um contrato verificável para que a equipe consiga descobrir se a hipótese se sustenta.

\section{Teste 28: Teclado}
Durante a consolidação do MVP, percebemos que objetivo: validar teclado no ambiente de demonstração.

No protótipo, esse entendimento aparece assim: Resultado esperado: Fluxo de proposta termina sem mouse.

Na implementação, este requisito é tratado como parte do sistema e não como uma tela isolada. A interface orienta a próxima ação; a API confere identidade, papel, estado e limites; a persistência guarda somente o necessário; e as integrações externas ficam condicionadas às variáveis do ambiente. O escopo continua exclusivamente de software, sem depender de escopo físico.

O fluxo começa quando a pessoa encontra uma necessidade ou uma habilidade que deseja compartilhar. Em seguida, ela confirma a intenção, executa a ação principal e recebe uma resposta que pode ser consultada depois. Para nós, teste 28: teclado só está concluído quando o caminho é compreensível, existe recuperação para erro e o servidor registra uma evidência coerente.

A forma mais honesta de verificar este requisito é registrar: Procedimento: preparar dados sintéticos, executar a ação, verificar resposta HTTP ou evento, conferir banco/ledger quando aplicável e registrar versão do código. O critério de aceite deve ser escrito antes do piloto, com entrada, resultado esperado, responsável, data e ferramenta de verificação. Assim, a documentação separa conformidade técnica, observação de uso e hipótese de impacto.

A validação combina teste automatizado, inspeção manual e uma tarefa representativa. Vamos registrar tempo, erros, conclusão, compreensão e barreiras observadas por teclado, leitor de tela, zoom, contraste, texto, áudio ou conexão instável. O objetivo não é criar uma aparência de certeza, e sim manter um ciclo curto entre observar, corrigir e conferir novamente.

A equipe não trata este item como resolvido antes de enfrentar o risco: Classificação: aprovado, reprovado ou bloqueado. Não marcar como aprovado sem evidência reproduzível; bloquear quando depender de credencial ou serviço externo ainda não configurado. O próximo passo é transformar a pendência em issue, teste, decisão de reunião ou evidência de piloto. Este capítulo não promete superioridade nem receita: ele deixa um contrato verificável para que a equipe consiga descobrir se a hipótese se sustenta.

\section{Teste 29: Leitor de tela}
A discussão de teste 29: leitor de tela começou por uma situação concreta: objetivo: validar leitor de tela no ambiente de demonstração.

Nossa escolha para este ponto foi direta: Resultado esperado: Headings, labels, foco e feedback são anunciados.

Na implementação, este requisito é tratado como parte do sistema e não como uma tela isolada. A interface orienta a próxima ação; a API confere identidade, papel, estado e limites; a persistência guarda somente o necessário; e as integrações externas ficam condicionadas às variáveis do ambiente. O escopo continua exclusivamente de software, sem depender de escopo físico.

O fluxo começa quando a pessoa encontra uma necessidade ou uma habilidade que deseja compartilhar. Em seguida, ela confirma a intenção, executa a ação principal e recebe uma resposta que pode ser consultada depois. Para nós, teste 29: leitor de tela só está concluído quando o caminho é compreensível, existe recuperação para erro e o servidor registra uma evidência coerente.

O aceite não fica apenas no texto. A evidência prevista é: Procedimento: preparar dados sintéticos, executar a ação, verificar resposta HTTP ou evento, conferir banco/ledger quando aplicável e registrar versão do código. O critério de aceite deve ser escrito antes do piloto, com entrada, resultado esperado, responsável, data e ferramenta de verificação. Assim, a documentação separa conformidade técnica, observação de uso e hipótese de impacto.

A validação combina teste automatizado, inspeção manual e uma tarefa representativa. Vamos registrar tempo, erros, conclusão, compreensão e barreiras observadas por teclado, leitor de tela, zoom, contraste, texto, áudio ou conexão instável. O objetivo não é criar uma aparência de certeza, e sim manter um ciclo curto entre observar, corrigir e conferir novamente.

Também registramos uma limitação, porque ela faz parte do resultado: Classificação: aprovado, reprovado ou bloqueado. Não marcar como aprovado sem evidência reproduzível; bloquear quando depender de credencial ou serviço externo ainda não configurado. O próximo passo é transformar a pendência em issue, teste, decisão de reunião ou evidência de piloto. Este capítulo não promete superioridade nem receita: ele deixa um contrato verificável para que a equipe consiga descobrir se a hipótese se sustenta.

\section{Teste 30: Voz indisponível}
A discussão de teste 30: voz indisponível começou por uma situação concreta: objetivo: validar voz indisponível no ambiente de demonstração.

Nossa escolha para este ponto foi direta: Resultado esperado: Teclado e texto continuam disponíveis.

Na implementação, este requisito é tratado como parte do sistema e não como uma tela isolada. A interface orienta a próxima ação; a API confere identidade, papel, estado e limites; a persistência guarda somente o necessário; e as integrações externas ficam condicionadas às variáveis do ambiente. O escopo continua exclusivamente de software, sem depender de escopo físico.

O fluxo começa quando a pessoa encontra uma necessidade ou uma habilidade que deseja compartilhar. Em seguida, ela confirma a intenção, executa a ação principal e recebe uma resposta que pode ser consultada depois. Para nós, teste 30: voz indisponível só está concluído quando o caminho é compreensível, existe recuperação para erro e o servidor registra uma evidência coerente.

O aceite não fica apenas no texto. A evidência prevista é: Procedimento: preparar dados sintéticos, executar a ação, verificar resposta HTTP ou evento, conferir banco/ledger quando aplicável e registrar versão do código. O critério de aceite deve ser escrito antes do piloto, com entrada, resultado esperado, responsável, data e ferramenta de verificação. Assim, a documentação separa conformidade técnica, observação de uso e hipótese de impacto.

A validação combina teste automatizado, inspeção manual e uma tarefa representativa. Vamos registrar tempo, erros, conclusão, compreensão e barreiras observadas por teclado, leitor de tela, zoom, contraste, texto, áudio ou conexão instável. O objetivo não é criar uma aparência de certeza, e sim manter um ciclo curto entre observar, corrigir e conferir novamente.

Também registramos uma limitação, porque ela faz parte do resultado: Classificação: aprovado, reprovado ou bloqueado. Não marcar como aprovado sem evidência reproduzível; bloquear quando depender de credencial ou serviço externo ainda não configurado. O próximo passo é transformar a pendência em issue, teste, decisão de reunião ou evidência de piloto. Este capítulo não promete superioridade nem receita: ele deixa um contrato verificável para que a equipe consiga descobrir se a hipótese se sustenta.

\section{Teste 31: Libras indisponível}
Na leitura que fizemos do projeto, objetivo: validar libras indisponível no ambiente de demonstração.

Depois de comparar as alternativas, a equipe decidiu o seguinte: Resultado esperado: Conteúdo alternativo continua completo.

Na implementação, este requisito é tratado como parte do sistema e não como uma tela isolada. A interface orienta a próxima ação; a API confere identidade, papel, estado e limites; a persistência guarda somente o necessário; e as integrações externas ficam condicionadas às variáveis do ambiente. O escopo continua exclusivamente de software, sem depender de escopo físico.

O fluxo começa quando a pessoa encontra uma necessidade ou uma habilidade que deseja compartilhar. Em seguida, ela confirma a intenção, executa a ação principal e recebe uma resposta que pode ser consultada depois. Para nós, teste 31: libras indisponível só está concluído quando o caminho é compreensível, existe recuperação para erro e o servidor registra uma evidência coerente.

Para conferir se a decisão funciona, vamos procurar esta evidência: Procedimento: preparar dados sintéticos, executar a ação, verificar resposta HTTP ou evento, conferir banco/ledger quando aplicável e registrar versão do código. O critério de aceite deve ser escrito antes do piloto, com entrada, resultado esperado, responsável, data e ferramenta de verificação. Assim, a documentação separa conformidade técnica, observação de uso e hipótese de impacto.

A validação combina teste automatizado, inspeção manual e uma tarefa representativa. Vamos registrar tempo, erros, conclusão, compreensão e barreiras observadas por teclado, leitor de tela, zoom, contraste, texto, áudio ou conexão instável. O objetivo não é criar uma aparência de certeza, e sim manter um ciclo curto entre observar, corrigir e conferir novamente.

O ponto que ainda merece cuidado é o seguinte: Classificação: aprovado, reprovado ou bloqueado. Não marcar como aprovado sem evidência reproduzível; bloquear quando depender de credencial ou serviço externo ainda não configurado. O próximo passo é transformar a pendência em issue, teste, decisão de reunião ou evidência de piloto. Este capítulo não promete superioridade nem receita: ele deixa um contrato verificável para que a equipe consiga descobrir se a hipótese se sustenta.

\section{Teste 32: Railway health}
A discussão de teste 32: railway health começou por uma situação concreta: objetivo: validar railway health no ambiente de demonstração.

Nossa escolha para este ponto foi direta: Resultado esperado: Deploy é considerado saudável apenas com banco e API respondendo.

Na implementação, este requisito é tratado como parte do sistema e não como uma tela isolada. A interface orienta a próxima ação; a API confere identidade, papel, estado e limites; a persistência guarda somente o necessário; e as integrações externas ficam condicionadas às variáveis do ambiente. O escopo continua exclusivamente de software, sem depender de escopo físico.

O fluxo começa quando a pessoa encontra uma necessidade ou uma habilidade que deseja compartilhar. Em seguida, ela confirma a intenção, executa a ação principal e recebe uma resposta que pode ser consultada depois. Para nós, teste 32: railway health só está concluído quando o caminho é compreensível, existe recuperação para erro e o servidor registra uma evidência coerente.

O aceite não fica apenas no texto. A evidência prevista é: Procedimento: preparar dados sintéticos, executar a ação, verificar resposta HTTP ou evento, conferir banco/ledger quando aplicável e registrar versão do código. O critério de aceite deve ser escrito antes do piloto, com entrada, resultado esperado, responsável, data e ferramenta de verificação. Assim, a documentação separa conformidade técnica, observação de uso e hipótese de impacto.

A validação combina teste automatizado, inspeção manual e uma tarefa representativa. Vamos registrar tempo, erros, conclusão, compreensão e barreiras observadas por teclado, leitor de tela, zoom, contraste, texto, áudio ou conexão instável. O objetivo não é criar uma aparência de certeza, e sim manter um ciclo curto entre observar, corrigir e conferir novamente.

Também registramos uma limitação, porque ela faz parte do resultado: Classificação: aprovado, reprovado ou bloqueado. Não marcar como aprovado sem evidência reproduzível; bloquear quando depender de credencial ou serviço externo ainda não configurado. O próximo passo é transformar a pendência em issue, teste, decisão de reunião ou evidência de piloto. Este capítulo não promete superioridade nem receita: ele deixa um contrato verificável para que a equipe consiga descobrir se a hipótese se sustenta.

\section{Teste 33: Secret scan}
Durante a consolidação do MVP, percebemos que objetivo: validar secret scan no ambiente de demonstração.

No protótipo, esse entendimento aparece assim: Resultado esperado: Pipeline falha ao detectar chave real em arquivo versionado.

Na implementação, este requisito é tratado como parte do sistema e não como uma tela isolada. A interface orienta a próxima ação; a API confere identidade, papel, estado e limites; a persistência guarda somente o necessário; e as integrações externas ficam condicionadas às variáveis do ambiente. O escopo continua exclusivamente de software, sem depender de escopo físico.

O fluxo começa quando a pessoa encontra uma necessidade ou uma habilidade que deseja compartilhar. Em seguida, ela confirma a intenção, executa a ação principal e recebe uma resposta que pode ser consultada depois. Para nós, teste 33: secret scan só está concluído quando o caminho é compreensível, existe recuperação para erro e o servidor registra uma evidência coerente.

A forma mais honesta de verificar este requisito é registrar: Procedimento: preparar dados sintéticos, executar a ação, verificar resposta HTTP ou evento, conferir banco/ledger quando aplicável e registrar versão do código. O critério de aceite deve ser escrito antes do piloto, com entrada, resultado esperado, responsável, data e ferramenta de verificação. Assim, a documentação separa conformidade técnica, observação de uso e hipótese de impacto.

A validação combina teste automatizado, inspeção manual e uma tarefa representativa. Vamos registrar tempo, erros, conclusão, compreensão e barreiras observadas por teclado, leitor de tela, zoom, contraste, texto, áudio ou conexão instável. O objetivo não é criar uma aparência de certeza, e sim manter um ciclo curto entre observar, corrigir e conferir novamente.

A equipe não trata este item como resolvido antes de enfrentar o risco: Classificação: aprovado, reprovado ou bloqueado. Não marcar como aprovado sem evidência reproduzível; bloquear quando depender de credencial ou serviço externo ainda não configurado. O próximo passo é transformar a pendência em issue, teste, decisão de reunião ou evidência de piloto. Este capítulo não promete superioridade nem receita: ele deixa um contrato verificável para que a equipe consiga descobrir se a hipótese se sustenta.

\section{Apêndice 01: Contrato de ambiente}
Durante a consolidação do MVP, percebemos que toda variável crítica tem nome, finalidade, ambiente, dono e estratégia de rotação.

No protótipo, esse entendimento aparece assim: A equipe deve transformar esta orientação em artefato operacional: arquivo, issue, checklist, teste ou registro de reunião.

Na implementação, este requisito é tratado como parte do sistema e não como uma tela isolada. A interface orienta a próxima ação; a API confere identidade, papel, estado e limites; a persistência guarda somente o necessário; e as integrações externas ficam condicionadas às variáveis do ambiente. O escopo continua exclusivamente de software, sem depender de escopo físico.

O fluxo começa quando a pessoa encontra uma necessidade ou uma habilidade que deseja compartilhar. Em seguida, ela confirma a intenção, executa a ação principal e recebe uma resposta que pode ser consultada depois. Para nós, apêndice 01: contrato de ambiente só está concluído quando o caminho é compreensível, existe recuperação para erro e o servidor registra uma evidência coerente.

A forma mais honesta de verificar este requisito é registrar: A documentação não substitui a execução. Ela fornece o contrato para que a equipe consiga executar, revisar e melhorar o produto. O critério de aceite deve ser escrito antes do piloto, com entrada, resultado esperado, responsável, data e ferramenta de verificação. Assim, a documentação separa conformidade técnica, observação de uso e hipótese de impacto.

A validação combina teste automatizado, inspeção manual e uma tarefa representativa. Vamos registrar tempo, erros, conclusão, compreensão e barreiras observadas por teclado, leitor de tela, zoom, contraste, texto, áudio ou conexão instável. O objetivo não é criar uma aparência de certeza, e sim manter um ciclo curto entre observar, corrigir e conferir novamente.

A equipe não trata este item como resolvido antes de enfrentar o risco: Responsáveis, datas e resultados devem ser preenchidos pela equipe antes da banca e novamente antes de qualquer piloto comercial. O próximo passo é transformar a pendência em issue, teste, decisão de reunião ou evidência de piloto. Este capítulo não promete superioridade nem receita: ele deixa um contrato verificável para que a equipe consiga descobrir se a hipótese se sustenta.

\section{Apêndice 02: Contrato de dados}
A discussão de apêndice 02: contrato de dados começou por uma situação concreta: dTOs de entrada têm limites, validação e mensagens estáveis para o frontend.

Nossa escolha para este ponto foi direta: A equipe deve transformar esta orientação em artefato operacional: arquivo, issue, checklist, teste ou registro de reunião.

Na implementação, este requisito é tratado como parte do sistema e não como uma tela isolada. A interface orienta a próxima ação; a API confere identidade, papel, estado e limites; a persistência guarda somente o necessário; e as integrações externas ficam condicionadas às variáveis do ambiente. O escopo continua exclusivamente de software, sem depender de escopo físico.

O fluxo começa quando a pessoa encontra uma necessidade ou uma habilidade que deseja compartilhar. Em seguida, ela confirma a intenção, executa a ação principal e recebe uma resposta que pode ser consultada depois. Para nós, apêndice 02: contrato de dados só está concluído quando o caminho é compreensível, existe recuperação para erro e o servidor registra uma evidência coerente.

O aceite não fica apenas no texto. A evidência prevista é: A documentação não substitui a execução. Ela fornece o contrato para que a equipe consiga executar, revisar e melhorar o produto. O critério de aceite deve ser escrito antes do piloto, com entrada, resultado esperado, responsável, data e ferramenta de verificação. Assim, a documentação separa conformidade técnica, observação de uso e hipótese de impacto.

A validação combina teste automatizado, inspeção manual e uma tarefa representativa. Vamos registrar tempo, erros, conclusão, compreensão e barreiras observadas por teclado, leitor de tela, zoom, contraste, texto, áudio ou conexão instável. O objetivo não é criar uma aparência de certeza, e sim manter um ciclo curto entre observar, corrigir e conferir novamente.

Também registramos uma limitação, porque ela faz parte do resultado: Responsáveis, datas e resultados devem ser preenchidos pela equipe antes da banca e novamente antes de qualquer piloto comercial. O próximo passo é transformar a pendência em issue, teste, decisão de reunião ou evidência de piloto. Este capítulo não promete superioridade nem receita: ele deixa um contrato verificável para que a equipe consiga descobrir se a hipótese se sustenta.

\section{Apêndice 03: Contrato de erro}
Na leitura que fizemos do projeto, erros de autorização, validação, conflito e integração são distinguíveis pelo cliente.

Depois de comparar as alternativas, a equipe decidiu o seguinte: A equipe deve transformar esta orientação em artefato operacional: arquivo, issue, checklist, teste ou registro de reunião.

Na implementação, este requisito é tratado como parte do sistema e não como uma tela isolada. A interface orienta a próxima ação; a API confere identidade, papel, estado e limites; a persistência guarda somente o necessário; e as integrações externas ficam condicionadas às variáveis do ambiente. O escopo continua exclusivamente de software, sem depender de escopo físico.

O fluxo começa quando a pessoa encontra uma necessidade ou uma habilidade que deseja compartilhar. Em seguida, ela confirma a intenção, executa a ação principal e recebe uma resposta que pode ser consultada depois. Para nós, apêndice 03: contrato de erro só está concluído quando o caminho é compreensível, existe recuperação para erro e o servidor registra uma evidência coerente.

Para conferir se a decisão funciona, vamos procurar esta evidência: A documentação não substitui a execução. Ela fornece o contrato para que a equipe consiga executar, revisar e melhorar o produto. O critério de aceite deve ser escrito antes do piloto, com entrada, resultado esperado, responsável, data e ferramenta de verificação. Assim, a documentação separa conformidade técnica, observação de uso e hipótese de impacto.

A validação combina teste automatizado, inspeção manual e uma tarefa representativa. Vamos registrar tempo, erros, conclusão, compreensão e barreiras observadas por teclado, leitor de tela, zoom, contraste, texto, áudio ou conexão instável. O objetivo não é criar uma aparência de certeza, e sim manter um ciclo curto entre observar, corrigir e conferir novamente.

O ponto que ainda merece cuidado é o seguinte: Responsáveis, datas e resultados devem ser preenchidos pela equipe antes da banca e novamente antes de qualquer piloto comercial. O próximo passo é transformar a pendência em issue, teste, decisão de reunião ou evidência de piloto. Este capítulo não promete superioridade nem receita: ele deixa um contrato verificável para que a equipe consiga descobrir se a hipótese se sustenta.

\section{Apêndice 04: Contrato de auditoria}
A discussão de apêndice 04: contrato de auditoria começou por uma situação concreta: ações administrativas informam ator, alvo, motivo, timestamp e correlação.

Nossa escolha para este ponto foi direta: A equipe deve transformar esta orientação em artefato operacional: arquivo, issue, checklist, teste ou registro de reunião.

Na implementação, este requisito é tratado como parte do sistema e não como uma tela isolada. A interface orienta a próxima ação; a API confere identidade, papel, estado e limites; a persistência guarda somente o necessário; e as integrações externas ficam condicionadas às variáveis do ambiente. O escopo continua exclusivamente de software, sem depender de escopo físico.

O fluxo começa quando a pessoa encontra uma necessidade ou uma habilidade que deseja compartilhar. Em seguida, ela confirma a intenção, executa a ação principal e recebe uma resposta que pode ser consultada depois. Para nós, apêndice 04: contrato de auditoria só está concluído quando o caminho é compreensível, existe recuperação para erro e o servidor registra uma evidência coerente.

O aceite não fica apenas no texto. A evidência prevista é: A documentação não substitui a execução. Ela fornece o contrato para que a equipe consiga executar, revisar e melhorar o produto. O critério de aceite deve ser escrito antes do piloto, com entrada, resultado esperado, responsável, data e ferramenta de verificação. Assim, a documentação separa conformidade técnica, observação de uso e hipótese de impacto.

A validação combina teste automatizado, inspeção manual e uma tarefa representativa. Vamos registrar tempo, erros, conclusão, compreensão e barreiras observadas por teclado, leitor de tela, zoom, contraste, texto, áudio ou conexão instável. O objetivo não é criar uma aparência de certeza, e sim manter um ciclo curto entre observar, corrigir e conferir novamente.

Também registramos uma limitação, porque ela faz parte do resultado: Responsáveis, datas e resultados devem ser preenchidos pela equipe antes da banca e novamente antes de qualquer piloto comercial. O próximo passo é transformar a pendência em issue, teste, decisão de reunião ou evidência de piloto. Este capítulo não promete superioridade nem receita: ele deixa um contrato verificável para que a equipe consiga descobrir se a hipótese se sustenta.

\section{Apêndice 05: Contrato de consentimento}
A discussão de apêndice 05: contrato de consentimento começou por uma situação concreta: o produto só coleta preferência ou dado adicional depois de explicar finalidade.

Nossa escolha para este ponto foi direta: A equipe deve transformar esta orientação em artefato operacional: arquivo, issue, checklist, teste ou registro de reunião.

Na implementação, este requisito é tratado como parte do sistema e não como uma tela isolada. A interface orienta a próxima ação; a API confere identidade, papel, estado e limites; a persistência guarda somente o necessário; e as integrações externas ficam condicionadas às variáveis do ambiente. O escopo continua exclusivamente de software, sem depender de escopo físico.

O fluxo começa quando a pessoa encontra uma necessidade ou uma habilidade que deseja compartilhar. Em seguida, ela confirma a intenção, executa a ação principal e recebe uma resposta que pode ser consultada depois. Para nós, apêndice 05: contrato de consentimento só está concluído quando o caminho é compreensível, existe recuperação para erro e o servidor registra uma evidência coerente.

O aceite não fica apenas no texto. A evidência prevista é: A documentação não substitui a execução. Ela fornece o contrato para que a equipe consiga executar, revisar e melhorar o produto. O critério de aceite deve ser escrito antes do piloto, com entrada, resultado esperado, responsável, data e ferramenta de verificação. Assim, a documentação separa conformidade técnica, observação de uso e hipótese de impacto.

A validação combina teste automatizado, inspeção manual e uma tarefa representativa. Vamos registrar tempo, erros, conclusão, compreensão e barreiras observadas por teclado, leitor de tela, zoom, contraste, texto, áudio ou conexão instável. O objetivo não é criar uma aparência de certeza, e sim manter um ciclo curto entre observar, corrigir e conferir novamente.

Também registramos uma limitação, porque ela faz parte do resultado: Responsáveis, datas e resultados devem ser preenchidos pela equipe antes da banca e novamente antes de qualquer piloto comercial. O próximo passo é transformar a pendência em issue, teste, decisão de reunião ou evidência de piloto. Este capítulo não promete superioridade nem receita: ele deixa um contrato verificável para que a equipe consiga descobrir se a hipótese se sustenta.

\section{Apêndice 06: Contrato de recomendação}
Na leitura que fizemos do projeto, toda sugestão devolve item, pontuação, fatores e alternativa de feedback.

Depois de comparar as alternativas, a equipe decidiu o seguinte: A equipe deve transformar esta orientação em artefato operacional: arquivo, issue, checklist, teste ou registro de reunião.

Na implementação, este requisito é tratado como parte do sistema e não como uma tela isolada. A interface orienta a próxima ação; a API confere identidade, papel, estado e limites; a persistência guarda somente o necessário; e as integrações externas ficam condicionadas às variáveis do ambiente. O escopo continua exclusivamente de software, sem depender de escopo físico.

O fluxo começa quando a pessoa encontra uma necessidade ou uma habilidade que deseja compartilhar. Em seguida, ela confirma a intenção, executa a ação principal e recebe uma resposta que pode ser consultada depois. Para nós, apêndice 06: contrato de recomendação só está concluído quando o caminho é compreensível, existe recuperação para erro e o servidor registra uma evidência coerente.

Para conferir se a decisão funciona, vamos procurar esta evidência: A documentação não substitui a execução. Ela fornece o contrato para que a equipe consiga executar, revisar e melhorar o produto. O critério de aceite deve ser escrito antes do piloto, com entrada, resultado esperado, responsável, data e ferramenta de verificação. Assim, a documentação separa conformidade técnica, observação de uso e hipótese de impacto.

A validação combina teste automatizado, inspeção manual e uma tarefa representativa. Vamos registrar tempo, erros, conclusão, compreensão e barreiras observadas por teclado, leitor de tela, zoom, contraste, texto, áudio ou conexão instável. O objetivo não é criar uma aparência de certeza, e sim manter um ciclo curto entre observar, corrigir e conferir novamente.

O ponto que ainda merece cuidado é o seguinte: Responsáveis, datas e resultados devem ser preenchidos pela equipe antes da banca e novamente antes de qualquer piloto comercial. O próximo passo é transformar a pendência em issue, teste, decisão de reunião ou evidência de piloto. Este capítulo não promete superioridade nem receita: ele deixa um contrato verificável para que a equipe consiga descobrir se a hipótese se sustenta.

\section{Apêndice 07: Contrato de evolução}
A discussão de apêndice 07: contrato de evolução começou por uma situação concreta: feature flag ou migration deve permitir liberar por ambiente e voltar com segurança.

Nossa escolha para este ponto foi direta: A equipe deve transformar esta orientação em artefato operacional: arquivo, issue, checklist, teste ou registro de reunião.

Na implementação, este requisito é tratado como parte do sistema e não como uma tela isolada. A interface orienta a próxima ação; a API confere identidade, papel, estado e limites; a persistência guarda somente o necessário; e as integrações externas ficam condicionadas às variáveis do ambiente. O escopo continua exclusivamente de software, sem depender de escopo físico.

O fluxo começa quando a pessoa encontra uma necessidade ou uma habilidade que deseja compartilhar. Em seguida, ela confirma a intenção, executa a ação principal e recebe uma resposta que pode ser consultada depois. Para nós, apêndice 07: contrato de evolução só está concluído quando o caminho é compreensível, existe recuperação para erro e o servidor registra uma evidência coerente.

O aceite não fica apenas no texto. A evidência prevista é: A documentação não substitui a execução. Ela fornece o contrato para que a equipe consiga executar, revisar e melhorar o produto. O critério de aceite deve ser escrito antes do piloto, com entrada, resultado esperado, responsável, data e ferramenta de verificação. Assim, a documentação separa conformidade técnica, observação de uso e hipótese de impacto.

A validação combina teste automatizado, inspeção manual e uma tarefa representativa. Vamos registrar tempo, erros, conclusão, compreensão e barreiras observadas por teclado, leitor de tela, zoom, contraste, texto, áudio ou conexão instável. O objetivo não é criar uma aparência de certeza, e sim manter um ciclo curto entre observar, corrigir e conferir novamente.

Também registramos uma limitação, porque ela faz parte do resultado: Responsáveis, datas e resultados devem ser preenchidos pela equipe antes da banca e novamente antes de qualquer piloto comercial. O próximo passo é transformar a pendência em issue, teste, decisão de reunião ou evidência de piloto. Este capítulo não promete superioridade nem receita: ele deixa um contrato verificável para que a equipe consiga descobrir se a hipótese se sustenta.

\section{Apêndice 08: Contrato de acessibilidade}
Na leitura que fizemos do projeto, cada componente novo traz nome acessível, foco, estado e fallback.

Depois de comparar as alternativas, a equipe decidiu o seguinte: A equipe deve transformar esta orientação em artefato operacional: arquivo, issue, checklist, teste ou registro de reunião.

Na implementação, este requisito é tratado como parte do sistema e não como uma tela isolada. A interface orienta a próxima ação; a API confere identidade, papel, estado e limites; a persistência guarda somente o necessário; e as integrações externas ficam condicionadas às variáveis do ambiente. O escopo continua exclusivamente de software, sem depender de escopo físico.

O fluxo começa quando a pessoa encontra uma necessidade ou uma habilidade que deseja compartilhar. Em seguida, ela confirma a intenção, executa a ação principal e recebe uma resposta que pode ser consultada depois. Para nós, apêndice 08: contrato de acessibilidade só está concluído quando o caminho é compreensível, existe recuperação para erro e o servidor registra uma evidência coerente.

Para conferir se a decisão funciona, vamos procurar esta evidência: A documentação não substitui a execução. Ela fornece o contrato para que a equipe consiga executar, revisar e melhorar o produto. O critério de aceite deve ser escrito antes do piloto, com entrada, resultado esperado, responsável, data e ferramenta de verificação. Assim, a documentação separa conformidade técnica, observação de uso e hipótese de impacto.

A validação combina teste automatizado, inspeção manual e uma tarefa representativa. Vamos registrar tempo, erros, conclusão, compreensão e barreiras observadas por teclado, leitor de tela, zoom, contraste, texto, áudio ou conexão instável. O objetivo não é criar uma aparência de certeza, e sim manter um ciclo curto entre observar, corrigir e conferir novamente.

O ponto que ainda merece cuidado é o seguinte: Responsáveis, datas e resultados devem ser preenchidos pela equipe antes da banca e novamente antes de qualquer piloto comercial. O próximo passo é transformar a pendência em issue, teste, decisão de reunião ou evidência de piloto. Este capítulo não promete superioridade nem receita: ele deixa um contrato verificável para que a equipe consiga descobrir se a hipótese se sustenta.

\section{Apêndice 09: Roteiro de demo}
Durante a consolidação do MVP, percebemos que uma apresentação de cinco minutos mostra cadastro, troca, saldo, auditoria e acessibilidade.

No protótipo, esse entendimento aparece assim: A equipe deve transformar esta orientação em artefato operacional: arquivo, issue, checklist, teste ou registro de reunião.

Na implementação, este requisito é tratado como parte do sistema e não como uma tela isolada. A interface orienta a próxima ação; a API confere identidade, papel, estado e limites; a persistência guarda somente o necessário; e as integrações externas ficam condicionadas às variáveis do ambiente. O escopo continua exclusivamente de software, sem depender de escopo físico.

O fluxo começa quando a pessoa encontra uma necessidade ou uma habilidade que deseja compartilhar. Em seguida, ela confirma a intenção, executa a ação principal e recebe uma resposta que pode ser consultada depois. Para nós, apêndice 09: roteiro de demo só está concluído quando o caminho é compreensível, existe recuperação para erro e o servidor registra uma evidência coerente.

A forma mais honesta de verificar este requisito é registrar: A documentação não substitui a execução. Ela fornece o contrato para que a equipe consiga executar, revisar e melhorar o produto. O critério de aceite deve ser escrito antes do piloto, com entrada, resultado esperado, responsável, data e ferramenta de verificação. Assim, a documentação separa conformidade técnica, observação de uso e hipótese de impacto.

A validação combina teste automatizado, inspeção manual e uma tarefa representativa. Vamos registrar tempo, erros, conclusão, compreensão e barreiras observadas por teclado, leitor de tela, zoom, contraste, texto, áudio ou conexão instável. O objetivo não é criar uma aparência de certeza, e sim manter um ciclo curto entre observar, corrigir e conferir novamente.

A equipe não trata este item como resolvido antes de enfrentar o risco: Responsáveis, datas e resultados devem ser preenchidos pela equipe antes da banca e novamente antes de qualquer piloto comercial. O próximo passo é transformar a pendência em issue, teste, decisão de reunião ou evidência de piloto. Este capítulo não promete superioridade nem receita: ele deixa um contrato verificável para que a equipe consiga descobrir se a hipótese se sustenta.

\section{Apêndice 10: Roteiro de piloto}
Durante a consolidação do MVP, percebemos que instituição escolhe turma, mede baseline, testa a plataforma e revisa indicadores.

No protótipo, esse entendimento aparece assim: A equipe deve transformar esta orientação em artefato operacional: arquivo, issue, checklist, teste ou registro de reunião.

Na implementação, este requisito é tratado como parte do sistema e não como uma tela isolada. A interface orienta a próxima ação; a API confere identidade, papel, estado e limites; a persistência guarda somente o necessário; e as integrações externas ficam condicionadas às variáveis do ambiente. O escopo continua exclusivamente de software, sem depender de escopo físico.

O fluxo começa quando a pessoa encontra uma necessidade ou uma habilidade que deseja compartilhar. Em seguida, ela confirma a intenção, executa a ação principal e recebe uma resposta que pode ser consultada depois. Para nós, apêndice 10: roteiro de piloto só está concluído quando o caminho é compreensível, existe recuperação para erro e o servidor registra uma evidência coerente.

A forma mais honesta de verificar este requisito é registrar: A documentação não substitui a execução. Ela fornece o contrato para que a equipe consiga executar, revisar e melhorar o produto. O critério de aceite deve ser escrito antes do piloto, com entrada, resultado esperado, responsável, data e ferramenta de verificação. Assim, a documentação separa conformidade técnica, observação de uso e hipótese de impacto.

A validação combina teste automatizado, inspeção manual e uma tarefa representativa. Vamos registrar tempo, erros, conclusão, compreensão e barreiras observadas por teclado, leitor de tela, zoom, contraste, texto, áudio ou conexão instável. O objetivo não é criar uma aparência de certeza, e sim manter um ciclo curto entre observar, corrigir e conferir novamente.

A equipe não trata este item como resolvido antes de enfrentar o risco: Responsáveis, datas e resultados devem ser preenchidos pela equipe antes da banca e novamente antes de qualquer piloto comercial. O próximo passo é transformar a pendência em issue, teste, decisão de reunião ou evidência de piloto. Este capítulo não promete superioridade nem receita: ele deixa um contrato verificável para que a equipe consiga descobrir se a hipótese se sustenta.

\section{Apêndice 11: Riscos abertos}
A discussão de apêndice 11: riscos abertos começou por uma situação concreta: rotação externa, revisão jurídica, testes com usuários e custo de produção ficam registrados.

Nossa escolha para este ponto foi direta: A equipe deve transformar esta orientação em artefato operacional: arquivo, issue, checklist, teste ou registro de reunião.

Na implementação, este requisito é tratado como parte do sistema e não como uma tela isolada. A interface orienta a próxima ação; a API confere identidade, papel, estado e limites; a persistência guarda somente o necessário; e as integrações externas ficam condicionadas às variáveis do ambiente. O escopo continua exclusivamente de software, sem depender de escopo físico.

O fluxo começa quando a pessoa encontra uma necessidade ou uma habilidade que deseja compartilhar. Em seguida, ela confirma a intenção, executa a ação principal e recebe uma resposta que pode ser consultada depois. Para nós, apêndice 11: riscos abertos só está concluído quando o caminho é compreensível, existe recuperação para erro e o servidor registra uma evidência coerente.

O aceite não fica apenas no texto. A evidência prevista é: A documentação não substitui a execução. Ela fornece o contrato para que a equipe consiga executar, revisar e melhorar o produto. O critério de aceite deve ser escrito antes do piloto, com entrada, resultado esperado, responsável, data e ferramenta de verificação. Assim, a documentação separa conformidade técnica, observação de uso e hipótese de impacto.

A validação combina teste automatizado, inspeção manual e uma tarefa representativa. Vamos registrar tempo, erros, conclusão, compreensão e barreiras observadas por teclado, leitor de tela, zoom, contraste, texto, áudio ou conexão instável. O objetivo não é criar uma aparência de certeza, e sim manter um ciclo curto entre observar, corrigir e conferir novamente.

Também registramos uma limitação, porque ela faz parte do resultado: Responsáveis, datas e resultados devem ser preenchidos pela equipe antes da banca e novamente antes de qualquer piloto comercial. O próximo passo é transformar a pendência em issue, teste, decisão de reunião ou evidência de piloto. Este capítulo não promete superioridade nem receita: ele deixa um contrato verificável para que a equipe consiga descobrir se a hipótese se sustenta.

\section{Apêndice 12: Conclusão}
A discussão de apêndice 12: conclusão começou por uma situação concreta: a SkillSwap é um produto de software promissor, mas a validação de impacto e receita continua trabalho da equipe.

Nossa escolha para este ponto foi direta: A equipe deve transformar esta orientação em artefato operacional: arquivo, issue, checklist, teste ou registro de reunião.

Na implementação, este requisito é tratado como parte do sistema e não como uma tela isolada. A interface orienta a próxima ação; a API confere identidade, papel, estado e limites; a persistência guarda somente o necessário; e as integrações externas ficam condicionadas às variáveis do ambiente. O escopo continua exclusivamente de software, sem depender de escopo físico.

O fluxo começa quando a pessoa encontra uma necessidade ou uma habilidade que deseja compartilhar. Em seguida, ela confirma a intenção, executa a ação principal e recebe uma resposta que pode ser consultada depois. Para nós, apêndice 12: conclusão só está concluído quando o caminho é compreensível, existe recuperação para erro e o servidor registra uma evidência coerente.

O aceite não fica apenas no texto. A evidência prevista é: A documentação não substitui a execução. Ela fornece o contrato para que a equipe consiga executar, revisar e melhorar o produto. O critério de aceite deve ser escrito antes do piloto, com entrada, resultado esperado, responsável, data e ferramenta de verificação. Assim, a documentação separa conformidade técnica, observação de uso e hipótese de impacto.

A validação combina teste automatizado, inspeção manual e uma tarefa representativa. Vamos registrar tempo, erros, conclusão, compreensão e barreiras observadas por teclado, leitor de tela, zoom, contraste, texto, áudio ou conexão instável. O objetivo não é criar uma aparência de certeza, e sim manter um ciclo curto entre observar, corrigir e conferir novamente.

Também registramos uma limitação, porque ela faz parte do resultado: Responsáveis, datas e resultados devem ser preenchidos pela equipe antes da banca e novamente antes de qualquer piloto comercial. O próximo passo é transformar a pendência em issue, teste, decisão de reunião ou evidência de piloto. Este capítulo não promete superioridade nem receita: ele deixa um contrato verificável para que a equipe consiga descobrir se a hipótese se sustenta.

% Adendo autoral da equipe. Escopo final exclusivamente de software.
\section{Reposicionamento para o SkillSwap Access}

O SkillSwap Access amplia a plataforma de troca de habilidades para uma jornada digital de aprendizagem, prática, evidência, recomendação e oportunidade. A mesma conta serve a usuários comuns e a pessoas com deficiência; a prioridade de acessibilidade é uma decisão de produto, não uma separação de público.

Este reposicionamento responde a uma fragmentação concreta. Cursos, comunidades, portfólios, vagas, mentorias, comunicação e recursos de acessibilidade normalmente exigem cadastros e contextos distintos. A plataforma propõe uma camada de continuidade: a pessoa informa interesses e habilidades, realiza uma atividade, registra evidência, recebe uma sugestão explicada e pode iniciar uma conexão com um projeto ou oportunidade.

\section{Declaração de autoria e limites de interpretação}

A equipe utilizou inteligência artificial somente para correção ortográfica, pontuação e revisão formal da redação. A autoria, as escolhas de requisitos, o código, a seleção das referências, os dados, os testes e a validação pertencem aos cinco autores. Nenhum resultado de usuário, métrica de impacto ou conclusão de superioridade é apresentado sem evidência produzida pelo protocolo.

\section{Problema, objetivo e hipótese}

O problema de produto é a perda de contexto entre aprender, demonstrar uma habilidade e acessar uma oportunidade. O problema de inclusão é que barreiras de navegação, contraste, leitura, comunicação e negociação de adaptações podem aparecer antes mesmo de a pessoa conseguir participar.

O objetivo é especificar e implementar uma plataforma acessível que una aprendizagem entre pares, evidência prática, recomendação explicável e oportunidades, com controles de privacidade e uma visão institucional agregada.

A hipótese de engenharia é que uma jornada contínua e explicável pode reduzir fricção de descoberta e aumentar a compreensão da próxima ação. A hipótese será considerada somente após piloto, métricas e análise da equipe; o desenho técnico não é prova de impacto social.

\section{Público universal e valor institucional}

Qualquer usuário pode ensinar, aprender, participar de um projeto, procurar mentoria ou montar portfólio. Pessoas com deficiência orientam requisitos de acessibilidade desde o início, enquanto usuários sem deficiência também se beneficiam de texto simples, áudio, teclado, contraste, legendas e navegação previsível.

Universidades podem acompanhar trilhas e evidências. Empresas podem divulgar desafios e oportunidades com critérios de acessibilidade. Instituições de inclusão podem conectar formação e apoio. Em todos os casos, indicadores devem ser agregados, autorizados e úteis para uma decisão de melhoria, nunca para reduzir pessoas a uma taxa.

\section{Funcionalidades incrementadas}

\begin{itemize}
  \item formação em texto, áudio, vídeo opcional, legendas e linguagem simples;
  \item desafios e projetos com evidências práticas, comentários e portfólio;
  \item recomendação explicável de cursos, mentores, projetos e oportunidades;
  \item catálogo de bolsas, projetos, mentorias e vagas com requisitos acessíveis;
  \item painel institucional com capacitação, participação e barreiras agregadas;
  \item chat e arquivos como caminho principal, com fallback textual quando vídeo não estiver disponível;
  \item Firebase Cloud Messaging opcional para notificações, sem token de sessão no navegador;
  \item mapas baseados em Leaflet/OpenStreetMap ou MapLibre quando úteis, sem bloquear o fluxo remoto;
  \item módulos comerciais fora do escopo da versão gratuita;
  \item observabilidade por Actuator e Prometheus, com Grafana opcional.
\end{itemize}

\section{Arquitetura de software}

\begin{center}
\fbox{\parbox{0.84\textwidth}{\centering Usuários e tecnologias assistivas\\[2pt] $\downarrow$\\[2pt] React + TypeScript + Vite\\[2pt] $\downarrow$\\[2pt] Firebase Auth/FCM \quad Spring Boot — monólito modular\\[2pt] $\downarrow$\\[2pt] PostgreSQL + Flyway \quad Actuator + Prometheus}}
\end{center}

A decisão pelo monólito modular é deliberada. Os módulos de identidade, perfis, aprendizagem, comunidades, chat, notificações e administração possuem responsabilidades e contratos claros, mas compartilham o ciclo de build e o deploy inicial. Isso reduz custo operacional e facilita a banca reproduzir o sistema. Microsserviços são uma opção de evolução depois de existir carga real, fronteiras estáveis, observabilidade e equipe para operar falhas distribuídas.

\section{Contrato do módulo Access}

O frontend autenticado usa a rota \texttt{/acesso}. O backend expõe \texttt{GET /api/access/overview}, \texttt{GET /api/access/opportunities}, \texttt{GET /api/access/recommendations} e \texttt{POST /api/access/opportunities}. A criação é protegida por \texttt{ADMIN}. A migração \texttt{V5\_\_skillswap\_access.sql} cria \texttt{tb\_access\_opportunities} e insere dados demonstrativos com identificadores idempotentes.

A pontuação inicial normaliza habilidades e interesses, separa requisitos por vírgula, busca correspondência parcial ou exata e acrescenta pesos pequenos para apoio de acessibilidade e modalidade remota. A resposta contém pontuação e explicação. A regra é intencionalmente simples para que a pessoa possa entender, contestar e corrigir seus dados.

\section{Acessibilidade como aceite}

Os testes precisam cobrir teclado sem mouse, ordem de foco, rótulos, mensagens de erro, contraste, zoom de 200\%, leitor de tela, leitura em voz alta, legendas, texto alternativo, linguagem simples e conexão limitada. Voz e VLibras são complementares. Nenhuma tela pode depender exclusivamente de cor, áudio, hover ou gesto.

\section{Segurança e privacidade}

O Firebase ID token é obtido pelo SDK e enviado no cabeçalho HTTP; não é salvo em \texttt{localStorage} e não aparece em query string. WebSockets são autenticados pela primeira mensagem, sem credencial na URL. O servidor valida o token e desconsidera um identificador de remetente divergente. O webhook de pagamento valida assinatura, status e idempotência.

Os segredos de banco, Firebase Admin, vídeo, e-mail e monitoramento devem ficar no gerenciador do Railway. Os arquivos reais fornecidos durante a preparação não entram na entrega; devem ser revogados e rotacionados. Logs não devem conter tokens, senhas, áudio bruto, conteúdo de chat ou dados pessoais desnecessários.

\section{Observabilidade}

Prometheus coleta séries temporais técnicas: latência, erros, disponibilidade, pool de banco, conexões WebSocket, tempo de recomendação, memória e CPU. Grafana visualiza esses dados, mas é opcional. Labels não devem conter e-mail, ID pessoal ou texto livre. O endpoint de métricas deve estar protegido e possuir retenção compatível com o custo do MVP.

\section{IA do produto}

O chatbot, a tradução, a simplificação de texto e a moderação usam o modelo estável \texttt{gemini-2.5-flash}. A variável \texttt{GEMINI\_MODEL} permite configuração explícita e a chave \texttt{GEMINI\_API\_KEY} fica somente no backend. O código envia a chave no cabeçalho da requisição, não na URL. A geração de imagens de conquistas é uma integração separada, pois este modelo Flash não é apresentado como gerador de imagens.

\section{Deploy econômico}

O primeiro deploy pode utilizar Railway para API e PostgreSQL, Firebase para identidade e FCM, e alternativas gratuitas para mapas e recursos de acessibilidade. Vídeo, SMS e pagamentos são condicionais; texto e arquivos garantem degradação segura. O backend deve iniciar com migrations, health check, CORS restrito, HTTPS, backups e evidência de restauração.

\section{Estado da arte e diferencial proposto}

ONTODAPS sustenta a necessidade de considerar deficiência no desenho de e-learning [10]. O estudo de interfaces de voz justifica manter canais complementares e não tratar voz como solução universal [11]. A recomendação educacional interpretável e sua revisão sistemática apoiam o uso de explicações e regras auditáveis [3, 9]. A aprendizagem entre pares orienta a inclusão de atividades e colaboração [12].

Nossa hipótese de diferencial é a integração, em uma única jornada, de acessibilidade, aprendizagem entre pares, evidência prática, recomendação explicável e oportunidades para usuários comuns e pessoas com deficiência, com visão institucional. Isso não significa que cada componente seja inédito. A superioridade somente poderá ser discutida com estudo comparativo, tarefas, amostra, métricas e análise de vieses.

\section{Métricas de validação}

O protocolo deve registrar tempo de tarefa, erros, abandono, conclusão, compreensão da explicação, satisfação, acessibilidade percebida, latência, disponibilidade e incidentes de segurança. A equipe deve publicar amostra, instrumento, data, limitações e dados anonimizados ou agregados. Resultados inexistentes permanecem como trabalho futuro, não como números estimados.

\section{Riscos e mitigação}

Os principais riscos são baixa massa crítica, recomendação inadequada, exposição de dados, dependência de provedores, barreira de acessibilidade, fraude em créditos e custo de operação. Mitigações incluem regras explicáveis, revisão humana, minimização de dados, fallbacks, limites, ledger idempotente, backups, métricas e testes com tecnologias assistivas.

\section{Roadmap}

Na primeira fase, a equipe estabiliza autenticação, perfil, troca, comunidade, chat, Access e acessibilidade. Na segunda, valida evidências práticas e painel institucional com dados agregados. Na terceira, mede retenção, qualidade de recomendação e conexão com oportunidades. Escala, modelos avançados e extração de microsserviços só devem ocorrer depois de evidência de uso e capacidade operacional.

\section{Conclusão do adendo}

O SkillSwap Access é uma proposta de produto e de pesquisa aplicada. Seu valor está na integração verificável e na acessibilidade como requisito do caminho principal. A documentação diferencia implementação, integração condicional, hipótese e resultado medido. O escopo físico não faz parte do TCC; o pacote entregue é software, documentação e protocolo de validação.

\section{Referências}
\begin{thebibliography}{99}
\bibitem{chang2019} CHANG, J.; LI, H.; SUN, B. Matching knowledge suppliers and demanders on a digital platform: a novel method. \textit{IEEE Access}, v. 7, p. 21331--21342, 2019. Disponível em: \href{https://ieeexplore.ieee.org/document/8628950}{IEEE Xplore}.

\bibitem{amin2023} AMIN, S. et al. Smart e-learning framework for personalized adaptive learning and sequential path recommendations using reinforcement learning. \textit{IEEE Access}, v. 11, p. 89769--89790, 2023. Disponível em: \href{https://ieeexplore.ieee.org/document/10220065}{IEEE Xplore}.

\bibitem{yun2022} YUN, Y. et al. Interpretable educational recommendation: an open framework based on Bayesian principal component analysis. In: \textit{2022 IEEE International Conference on Systems, Man, and Cybernetics}, 2022. p. 3409--3414. Disponível em: \href{https://ieeexplore.ieee.org/document/9945498}{IEEE Xplore}.

\bibitem{lorgat2022} LORGAT, M. G.; PAREDES, H.; ROCHA, T. A systematic review of the promotion of accessible software development. In: \textit{2022 11th International Conference on Computer Technologies and Development}, 2022. p. 17--23. Disponível em: \href{https://ieeexplore.ieee.org/document/10153876}{IEEE Xplore}.

\bibitem{trielsa2023} TRIELSA, C. J.; ANGELINE, M. The effect of inclusive design on easy accessibility for disabled e-commerce users in Indonesia. In: \textit{2023 17th International Conference on Ubiquitous Information Management and Communication}, 2023. p. 1--4. Disponível em: \href{https://ieeexplore.ieee.org/document/10035640}{IEEE Xplore}.

\bibitem{kondoro2023} KONDORO, A. et al. Towards improving the accessibility and usability of a mobile-based learning management system for blind primary school teachers. In: \textit{2023 IST-Africa Conference}, 2023. p. 1--10. Disponível em: \href{https://ieeexplore.ieee.org/document/10187797}{IEEE Xplore}.

\bibitem{srivastava2022} SRIVASTAVA, S.; SRIVASTAVA, R.; TIWARI, N. Developing an e-learning and job portal for IT aspirants. In: \textit{2022 International Conference on Applied Artificial Intelligence and Computing}, 2022. p. 160--165. Disponível em: \href{https://ieeexplore.ieee.org/document/9792727}{IEEE Xplore}.

\bibitem{jhalani2024} JHALANI, Y. et al. Platform for helping senior citizens and enhancing employment for specially-abled. In: \textit{2024 International Conference on Data Science and Network Security}, 2024. p. 220--225. Disponível em: \href{https://ieeexplore.ieee.org/document/10691250}{IEEE Xplore}.

\bibitem{pesovski2022} PESOVSKI, I.; BOGDANOVA, A. M.; TRAJKOVIK, V. Systematic review of the published explainable educational recommendation systems. In: \textit{2022 20th International Conference on Information Technology Based Higher Education and Training}, 2022. p. 1--8. Disponível em: \href{https://ieeexplore.ieee.org/document/10032029}{IEEE Xplore}.

\bibitem{nganji2014} NGANJI, J. T.; BRAYSHAW, M. Designing and reflecting on disability-aware e-learning systems: the case of ONTODAPS. In: \textit{2014 IEEE 14th International Conference on Advanced Learning Technologies}, 2014. p. 571--575. Disponível em: \href{https://ieeexplore.ieee.org/document/6901543}{IEEE Xplore}.

\bibitem{jin2025} JIN, D. H.; KIM, Y. M.; LEE, Y. Investigating the acceptance of voice user interfaces for users with disabilities. \textit{IEEE Access}, v. 13, p. 1055--1069, 2025. Disponível em: \href{https://ieeexplore.ieee.org/document/10807228}{IEEE Xplore}.

\bibitem{gopal2021} GOPAL, B.; COOPER, S. Peer instruction in online synchronous software engineering: findings from fine-grained clicker data. In: \textit{2021 IEEE Frontiers in Education Conference}, 2021. Disponível em: \href{https://ieeexplore.ieee.org/document/9637353}{IEEE Xplore}.
\end{thebibliography}
\chapter{Atualização do escopo gratuito e da acessibilidade}

A versão final da SkillSwap não cobra dos usuários. Cadastro, troca de habilidades, conversa, agendamento, aula e geração do comprovante de conclusão permanecem disponíveis sem pagamento obrigatório. As rotas relacionadas a carteira, SwapCoins, premium e Stripe foram retiradas da navegação e bloqueadas na API. O código legado dessas integrações é mantido somente para rastreabilidade da evolução e não representa uma funcionalidade disponível.

\section{Preferências escolhidas pela pessoa}

Após o cadastro com Firebase, a plataforma apresenta uma etapa opcional de configuração. A pessoa pode escolher contraste intenso, tamanho do texto, canal preferido de comunicação, legendas quando disponíveis, descrição de imagens, VLibras, linguagem simplificada e operação por teclado. Essas respostas são preferências de uso e não servem para diagnosticar, classificar ou expor uma deficiência.

As cores, formas e marcadores visuais são apenas temas de interface. Não existe uma regra que associe automaticamente uma cor ou uma forma a autismo, daltonismo, cegueira, deficiência motora ou ausência de deficiência. Essa decisão preserva a autonomia da pessoa e evita que uma característica sensível seja inferida de maneira incorreta.

\section{Contraste e comunicação equivalente}

O contraste intenso utiliza fundo escuro, texto claro, foco visível, bordas destacadas e informações textuais junto aos elementos visuais. O vermelho não é usado como único indicador de erro. O sistema mantém texto, teclado e mensagens de estado como caminhos alternativos quando áudio, vídeo, câmera, microfone ou conexão externa não estiverem disponíveis.

VLibras e legendas são recursos complementares. A presença de um widget ou de uma opção de legenda não é apresentada como certificação WCAG. A avaliação completa depende de testes com teclado, leitores de tela, zoom, tecnologias assistivas e participantes voluntários.

\section{Conclusão da aula e comprovante}

A conclusão da aula depende da confirmação do aluno e do professor. Depois das duas confirmações, o aluno pode gerar uma página acessível de certificado contendo seu nome, a aula, o professor, a duração, a data e um identificador. O arquivo pode ser impresso ou salvo como PDF pelo navegador.

O comprovante registra participação na atividade da SkillSwap. Ele não equivale a diploma, certificado profissional ou validação externa de competência.


\end{document}
